% Qu’est-ce que la Philosophie ? — Philosophie 1re Tech — Cours signé OnBuch+
% Programme officiel MINESEC. Compiler avec : tectonic N20-qu-est-ce-que-philosophie.tex
\newcommand{\DOCMATIERE}{Philosophie}
\newcommand{\DOCNIVEAU}{1re Tech}
\newcommand{\DOCMODULE}{Philosophie – classe de Première technique}
\newcommand{\DOCLECON}{N20}
\newcommand{\DOCTITRE}{Qu’est-ce que la Philosophie ?}
% OnBuch+ — préambule « cours signés » ENRICHI (compilé avec tectonic / XeLaTeX)
% Système visuel « Pop » d'OnBuch+ : fond crème (jamais blanc), bordures encre
% épaisses, ombres « dures » décalées, titres Archivo Black en capitales.
% Version MATHÉMATIQUES : graphes (pgfplots/tikz), boîtes pédagogiques
% (définition, propriété, méthode, attention, le savais-tu, exercice résolu,
% exercice, TP, bilan), page de garde et signature OnBuch+ ancrée en bas.
%
% Macros attendues dans main.tex AVANT \input{preamble} :
%   \DOCMATIERE  \DOCNIVEAU  \DOCMODULE  \DOCLECON (numéro)  \DOCTITRE
\documentclass[11pt]{article}

\usepackage{fontspec}
\usepackage[french]{babel}
\usepackage[a4paper,margin=2cm,top=3.4cm,bottom=2.8cm,headheight=1.4cm,footskip=1.3cm]{geometry}
\usepackage{amsmath,amssymb,amsfonts}
\usepackage{mathtools} % \xmapsto, \coloneqq... souvent utilisés par les modèles
\usepackage{cancel} % simplifications barrées (bilans de forces)
% \abs{z} (module, valeur absolue) et \norm{u} (norme), utilisés par les modèles
\providecommand{\abs}{}\let\abs\relax\DeclarePairedDelimiter{\abs}{\lvert}{\rvert}
\providecommand{\norm}{}\let\norm\relax\DeclarePairedDelimiter{\norm}{\lVert}{\rVert}
\usepackage[table]{xcolor} % [table] : fournit aussi \rowcolors (alternance de lignes)
\usepackage{pagecolor}
\usepackage{tikz}
\usetikzlibrary{arrows.meta,positioning,calc,shapes.geometric,shapes.misc,shapes.arrows,decorations.pathmorphing,decorations.pathreplacing,decorations.markings,patterns,shadows,mindmap,trees,fit,backgrounds,matrix,intersections,angles,quotes}
\usepackage{pgfplots}
\usepackage[version=4]{mhchem} % équations nucléaires \ce{^{235}_{92}U}
\usepackage[european,siunitx]{circuitikz} % schémas électriques
\pgfplotsset{compat=1.18}
\usepgfplotslibrary{fillbetween,groupplots} % aires entre courbes (calcul intégral)
\usepackage{enumitem}
\usepackage{fancyhdr}
% [most] charge toutes les bibliothèques tcolorbox (listings, minted, raster,
% documentation...) et gonfle inutilement la mémoire TeX sur un document long
% (~300 boîtes) : on ne charge que ce qui est réellement utilisé.
\usepackage[skins,breakable]{tcolorbox}
\usepackage{graphicx}
\usepackage{colortbl}
\usepackage{booktabs}
\usepackage{tabularx}
\usepackage{multirow}
\usepackage{diagbox} % case d'en-tête diagonale des tableaux à double entrée
% Colonnes extensibles centrée / gauche / droite (C, L, R), souvent utilisées par les modèles
\newcolumntype{C}{>{\centering\arraybackslash}X}
\newcolumntype{L}{>{\raggedright\arraybackslash}X}
\newcolumntype{R}{>{\raggedleft\arraybackslash}X}
\usepackage{array}
\usepackage{multicol}
\usepackage{siunitx}
% parse-numbers=false : accepte aussi \SI{2,5 \times 10^{-3}}{...}
\sisetup{locale=FR,output-decimal-marker={,},group-minimum-digits=4,parse-numbers=false}
\DeclareSIUnit\ohm{\ensuremath{\symup{\Omega}}} % U+2126 absent de la police
\DeclareSIUnit\division{div} % réglages d'oscilloscope (ms/div, V/div)
\DeclareSIUnit\tr{tr} % tours (vitesse de rotation, ex. \SI{1500}{\tr\per\minute})
\DeclareSIUnit\molar{M} % concentration molaire (ex. \SI{3}{\molar}, \SI{1}{\micro\molar})
\DeclareSIUnit\an{an} % année (le modèle confond parfois avec \year, un compteur TeX) — ex. \SI{5}{\kilo\meter\per\an}
\DeclareSIUnit\FCFA{FCFA} % franc CFA (ex. \SI{500}{\FCFA\per\kilo\gram})
\DeclareSIUnit\degreeBrix{\textdegree Brix} % teneur en sucre d'un sirop/jus (réfractométrie)
\DeclareSIUnit\month{mois} % le modèle utilise parfois \month, un compteur TeX, comme unité de durée
\DeclareSIUnit\microliter{\micro\liter} % le modèle écrit parfois \microliter en un seul token
\DeclareSIUnit\million{millions} % dénombrement (ex. \SI{15}{\million\per\mL} spermatozoïdes)
\usepackage{qrcode}
% \degree : commande fréquemment utilisée par les modèles pour un angle ou une
% température en dehors de \SI{}{} (non fournie par les paquets chargés).
\providecommand{\degree}{^{\circ}}
% \qed (fin de démonstration), utilisé par les modèles sans amsthm
\providecommand{\qed}{\hfill$\square$}
% \barre : double barre des valeurs interdites dans un tableau de variations (tkz-tab)
\providecommand{\barre}{\ensuremath{\|}}
% environnement proof (amsthm non chargé) utilisé par les modèles
\ifdefined\proof\else\newenvironment{proof}[1][Démonstration]{\par\noindent\textit{#1.}\ }{\qed\par}\fi
\usepackage{lastpage}
\usepackage{hyperref}
\hypersetup{hidelinks}

% ============================================================
% Palette « Pop » OnBuch+ — mêmes hex que la classe P (plus_ui.dart)
% ============================================================
\definecolor{poporange}{HTML}{F59321}
\definecolor{popdark}{HTML}{E07A0C}
\definecolor{popcream}{HTML}{FFF6E8}
\definecolor{popink}{HTML}{1C1714}
\definecolor{poppurple}{HTML}{7A5AE0}
\definecolor{popgreen}{HTML}{2E9E6A}
\definecolor{poppink}{HTML}{E0457B}
\definecolor{popblue}{HTML}{2F7DE1}
\definecolor{popgold}{HTML}{C98A12}
\definecolor{popmuted}{HTML}{8A7F76}
\definecolor{popline}{HTML}{EADFD2}
\definecolor{popgray}{HTML}{8A7F76}
\colorlet{popgrey}{popgray}
% Alias tolérants (le modèle nomme parfois les couleurs différemment)
\colorlet{popwhite}{white}
\colorlet{popblack}{popink}
\colorlet{popred}{poppink}
\colorlet{popyellow}{popgold}
% Teintes claires (fonds de tableaux/encadrés) — le modèle nomme parfois
% les couleurs avec un suffixe L (ex. popblueL) sans les avoir définies.
\colorlet{poporangeL}{poporange!20}
\colorlet{popdarkL}{popdark!20}
\colorlet{poppurpleL}{poppurple!20}
\colorlet{popgreenL}{popgreen!20}
\colorlet{poppinkL}{poppink!20}
\colorlet{popblueL}{popblue!20}
\colorlet{popgoldL}{popgold!20}
\colorlet{popredL}{popred!20}
\colorlet{popgrayL}{popgray!20}
% Teintes claires pour fonds de tableaux / zones
\colorlet{popblueL}{popblue!10!white}
\colorlet{poporangeL}{poporange!14!white}
\colorlet{popgreenL}{popgreen!12!white}
\colorlet{poppurpleL}{poppurple!12!white}
\colorlet{poppinkL}{poppink!12!white}
\colorlet{popgoldL}{popgold!14!white}

\pagecolor{popcream}

% ============================================================
% Typographie
% ============================================================
\defaultfontfeatures{Ligatures=TeX}
\setmainfont{PlusJakartaSans-Medium.ttf}[Path=../../../fonts/,BoldFont=PlusJakartaSans-Bold.ttf]
% Maths en sans-serif (Fira Math) avec chiffres et lettres droites de Plus
% Jakarta, pour que les formules se fondent dans le texte.
\usepackage[math-style=ISO,bold-style=ISO]{unicode-math}
\setmathfont{FiraMath-Regular.otf}[Path=../../../fonts/]
\setmathfont{PlusJakartaSans-Medium.ttf}[Path=../../../fonts/,range={up/num,up/{Latin,latin},\mathup}]
\usepackage{icomma} % virgule décimale sans espace en mode math
% Caractères mathématiques Unicode absents des polices de texte (Plus Jakarta,
% Archivo Black) : rendus par leur équivalent mathématique, en texte comme en
% formule — sinon ils s'affichent en carré vide (ex. « Arithmétique dans ℤ »).
\usepackage{newunicodechar}
\newunicodechar{ℝ}{\ensuremath{\mathbb{R}}}
\newunicodechar{ℤ}{\ensuremath{\mathbb{Z}}}
\newunicodechar{ℂ}{\ensuremath{\mathbb{C}}}
\newunicodechar{ℕ}{\ensuremath{\mathbb{N}}}
\newunicodechar{ℚ}{\ensuremath{\mathbb{Q}}}
\newunicodechar{𝔻}{\ensuremath{\mathbb{D}}}
\newunicodechar{α}{\ensuremath{\alpha}}
\newunicodechar{β}{\ensuremath{\beta}}
\newunicodechar{γ}{\ensuremath{\gamma}}
\newunicodechar{δ}{\ensuremath{\delta}}
\newunicodechar{ε}{\ensuremath{\varepsilon}}
\newunicodechar{θ}{\ensuremath{\theta}}
\newunicodechar{λ}{\ensuremath{\lambda}}
\newunicodechar{μ}{\ensuremath{\mu}}
\newunicodechar{σ}{\ensuremath{\sigma}}
\newunicodechar{τ}{\ensuremath{\tau}}
\newunicodechar{φ}{\ensuremath{\varphi}}
\newunicodechar{ψ}{\ensuremath{\psi}}
\newunicodechar{ω}{\ensuremath{\omega}}
\newunicodechar{Σ}{\ensuremath{\Sigma}}
% majuscules produites par \MakeUppercase dans les titres (α-aminés -> Α)
\newunicodechar{Α}{\ensuremath{\alpha}}
\newunicodechar{Β}{\ensuremath{\beta}}
\newunicodechar{Γ}{\ensuremath{\Gamma}}
\newunicodechar{Δ}{\ensuremath{\Delta}}
\newunicodechar{Ε}{\ensuremath{\varepsilon}}
\newunicodechar{Θ}{\ensuremath{\Theta}}
\newunicodechar{Λ}{\ensuremath{\Lambda}}
\newunicodechar{Μ}{\ensuremath{\mu}}
\newunicodechar{Τ}{\ensuremath{\tau}}
\newunicodechar{Φ}{\ensuremath{\Phi}}
\newunicodechar{Ψ}{\ensuremath{\Psi}}
\newunicodechar{Ω}{\ensuremath{\Omega}}
\newunicodechar{↦}{\ensuremath{\mapsto}}
\newunicodechar{∈}{\ensuremath{\in}}
\newunicodechar{∉}{\ensuremath{\notin}}
\newunicodechar{∧}{\ensuremath{\wedge}}
\newunicodechar{⇔}{\ensuremath{\Leftrightarrow}}
\newunicodechar{⟺}{\ensuremath{\Longleftrightarrow}}
\newunicodechar{⇒}{\ensuremath{\Rightarrow}}
\newunicodechar{⟹}{\ensuremath{\Longrightarrow}}
\newunicodechar{∘}{\ensuremath{\circ}}
\newunicodechar{∪}{\ensuremath{\cup}}
\newunicodechar{∩}{\ensuremath{\cap}}
\newunicodechar{≡}{\ensuremath{\equiv}}
\newunicodechar{⊂}{\ensuremath{\subset}}
\newunicodechar{⊥}{\ensuremath{\perp}}
\newunicodechar{∃}{\ensuremath{\exists}}
\newunicodechar{∀}{\ensuremath{\forall}}
\newunicodechar{∛}{\ensuremath{\sqrt[3]{\,}}}
\newunicodechar{∜}{\ensuremath{\sqrt[4]{\,}}}
\newunicodechar{⃗}{\ensuremath{{}^{\to}}}
\newunicodechar{ⁿ}{\ifmmode^{n}\else\textsuperscript{\ensuremath{n}}\fi}
\newunicodechar{ˣ}{\ifmmode^{x}\else\textsuperscript{\ensuremath{x}}\fi}
\newunicodechar{ᵇ}{\ifmmode^{b}\else\textsuperscript{\ensuremath{b}}\fi}
\newunicodechar{ʳ}{\ifmmode^{r}\else\textsuperscript{\ensuremath{r}}\fi}
\newunicodechar{ᵘ}{\ifmmode^{u}\else\textsuperscript{\ensuremath{u}}\fi}
\newunicodechar{ʸ}{\ifmmode^{y}\else\textsuperscript{\ensuremath{y}}\fi}
\newunicodechar{ᵃ}{\ifmmode^{a}\else\textsuperscript{\ensuremath{a}}\fi}
\newunicodechar{ᵉ}{\ifmmode^{e}\else\textsuperscript{\ensuremath{e}}\fi}
\newunicodechar{ᵏ}{\ifmmode^{k}\else\textsuperscript{\ensuremath{k}}\fi}
\newunicodechar{ᵖ}{\ifmmode^{p}\else\textsuperscript{\ensuremath{p}}\fi}
\newunicodechar{ᴺ}{\ifmmode^{N}\else\textsuperscript{\ensuremath{N}}\fi}
\newunicodechar{ᶜ}{\ifmmode^{c}\else\textsuperscript{\ensuremath{c}}\fi}
\newunicodechar{ʰ}{\ifmmode^{h}\else\textsuperscript{\ensuremath{h}}\fi}
\newunicodechar{ᵠ}{\ifmmode^{q}\else\textsuperscript{\ensuremath{q}}\fi}
\newunicodechar{ₙ}{\ifmmode_{n}\else\textsubscript{\ensuremath{n}}\fi}
\newunicodechar{ₐ}{\ifmmode_{a}\else\textsubscript{\ensuremath{a}}\fi}
\newunicodechar{ᵢ}{\ifmmode_{i}\else\textsubscript{\ensuremath{i}}\fi}
\newunicodechar{ᵣ}{\ifmmode_{r}\else\textsubscript{\ensuremath{r}}\fi}
\newunicodechar{ₖ}{\ifmmode_{k}\else\textsubscript{\ensuremath{k}}\fi}
\newunicodechar{ᵦ}{\ifmmode_{\beta}\else\textsubscript{\ensuremath{\beta}}\fi}
\newunicodechar{ₓ}{\ifmmode_{x}\else\textsubscript{\ensuremath{x}}\fi}
\newfontfamily\popdisplay{ArchivoBlack-Regular.ttf}[Path=../../../fonts/]
\newfontfamily\poplogo{SpaceGrotesk-Bold.ttf}[Path=../../../fonts/]
\newfontfamily\popsemi{PlusJakartaSans-SemiBold.ttf}[Path=../../../fonts/]
\newfontfamily\popxbold{PlusJakartaSans-ExtraBold.ttf}[Path=../../../fonts/]
\newfontfamily\popxboldls{PlusJakartaSans-ExtraBold.ttf}[Path=../../../fonts/,LetterSpace=3.0]

\setlist[enumerate]{leftmargin=1.7em,itemsep=5pt,topsep=4pt}
\setlist[itemize]{leftmargin=1.7em,itemsep=4pt,topsep=4pt,label={\color{poporange}\textbullet}}
\setlength{\parindent}{0pt}
\setlength{\parskip}{5pt}
\renewcommand{\arraystretch}{1.25}

% pgfplots : style Pop par défaut
\pgfplotsset{
  every axis/.append style={axis line style={popink,line width=1pt},tick style={popink},
    grid style={popline,line width=0.5pt},label style={font=\small},tick label style={font=\footnotesize}},
  popcurve/.style={poporange,line width=1.8pt,no markers,smooth},
}
% Même style disponible dans un \draw TikZ simple (hors pgfplots)
\tikzset{popcurve/.style={poporange,line width=1.8pt}}
\tikzset{popgrid/.style={popline,very thin}}
\colorlet{popcurve}{poporange} % le nom du style est parfois utilisé comme couleur

% ============================================================
% Pill / squiggle
% ============================================================
\newcommand{\poppill}[3]{%
  \tikz[baseline=-0.55ex]\node[draw=popink,line width=1.2pt,rounded corners=2.6pt,%
    fill=#2,inner xsep=6.5pt,inner ysep=3pt]{%
    \popxboldls\fontsize{7}{7}\selectfont\color{#3}\MakeUppercase{#1}};%
}
\newcommand{\popsquiggle}[1]{%
  \tikz[baseline=-0.4ex]{
    \draw[#1,line width=2.6pt,line cap=round]
      plot [smooth] coordinates {(0,0.09) (0.34,0.01) (0.68,0.21) (1.02,0.01) (1.36,0.10)};
  }%
}
% Mot-clé surligné (marqueur orange sous le texte)
\newcommand{\cle}[1]{{\popxbold\color{popdark}#1}}

% ============================================================
% En-tête / pied
% ============================================================
\pagestyle{fancy}
\fancyhf{}
\renewcommand{\headrulewidth}{0pt}
\renewcommand{\footrulewidth}{0pt}
\lhead{\raisebox{-0.15ex}{\poplogo\fontsize{13}{13}\selectfont\color{popink}OnBuch\color{poporange}+}}
\rhead{\poppill{\DOCMATIERE\ \textbullet\ \DOCNIVEAU\ \textbullet\ Leçon \DOCLECON}{white}{popink}}
\lfoot{\popxbold\fontsize{7.6}{7.6}\selectfont\color{popmuted}\MakeUppercase{Cours signé OnBuch+}\ \textbullet\ {\popsemi\DOCTITRE}}
\rfoot{\tikz[baseline=-0.6ex]\node[rounded corners=8.5pt,fill=popink,minimum height=17pt,minimum width=17pt,inner xsep=5pt,inner sep=0]{\popxbold\fontsize{8}{8}\selectfont\color{popcream}\thepage\,/\,\pageref*{LastPage}};}

% ============================================================
% Titres
% ============================================================
\newcommand{\chaptitre}[2]{%
  \begin{tcolorbox}[enhanced,colback=poporange,colframe=popink,boxrule=2.6pt,arc=11pt,
      drop shadow=popink,left=15pt,right=15pt,top=12pt,bottom=11pt,before skip=3pt,after skip=18pt]
    \poppill{#2}{popink}{popcream}\par\vspace{9pt}
    {\raggedright\hyphenpenalty=10000\exhyphenpenalty=10000\popdisplay\fontsize{22}{24}\selectfont\color{popcream}\MakeUppercase{#1}\par}
  \end{tcolorbox}
}
% Section numérotée : pastille orange avec le numéro + titre Archivo
\newcounter{popsec}
\newcommand{\coursec}[1]{%
  \stepcounter{popsec}\setcounter{popsub}{0}%
  \par\vspace{16pt}\noindent
  \begin{minipage}{\linewidth}
  \tikz[baseline=-0.7ex]\node[draw=popink,line width=1.6pt,fill=poporange,rounded corners=4pt,minimum size=22pt,inner sep=0]{\popdisplay\fontsize{12}{12}\selectfont\color{popcream}\thepopsec};%
  \hspace{8pt}{\popdisplay\fontsize{15}{17}\selectfont\color{popink}\MakeUppercase{#1}}\par
  \vspace{4pt}\noindent{\color{poporange}\rule{36pt}{3.6pt}}
  \end{minipage}\par\nopagebreak\vspace{8pt}
}
\newcounter{popsub}
\newcommand{\courssub}[1]{%
  \stepcounter{popsub}%
  \par\vspace{9pt}\noindent{\popxbold\fontsize{11.5}{13}\selectfont\color{popink}{\color{poporange}\thepopsec.\thepopsub}\hspace{6pt}#1}\par\nopagebreak\vspace{4pt}
}

% ============================================================
% Boîtes pédagogiques — même recette : fond blanc, bordure épaisse colorée,
% ombre dure encre, badge plein. Titre optionnel [#1].
% ============================================================
\tcbset{popbase/.style={breakable,enhanced,colback=white,boxrule=2.3pt,arc=9pt,
  shadow={3pt}{-3pt}{0pt}{popink},left=14pt,right=14pt,top=11pt,bottom=11pt,before skip=11pt,after skip=15pt,
  fontupper=\popsemi\fontsize{10.6}{14.5}\selectfont\color{popink},
  fontlower=\popsemi\fontsize{10.6}{14.5}\selectfont\color{popink}}}
% Coche dessinée (le glyphe ✓ n'existe pas dans les polices du document)
\newcommand{\popcheck}[1]{\tikz[baseline=-0.5ex]\draw[#1,line width=1.6pt,line cap=round,line join=round](-0.13,0.0)--(-0.04,-0.09)--(0.13,0.1);}
\providecommand{\coche}{\popcheck{popgreen}}
\providecommand{\resultat}[1]{\par\smallskip\noindent\fcolorbox{popgreen}{popgreenL}{\parbox{\dimexpr\linewidth-2\fboxsep-2\fboxrule\relax}{\textbf{Résultat.} #1}}\par\smallskip}
\newcommand{\popboxhead}[3]{\poppill{#1}{#2}{white}\ifx\relax#3\relax\else\hspace{6pt}{\popxbold\fontsize{10.6}{12}\selectfont\color{popink}#3}\fi\par\vspace{7pt}}

\newtcolorbox{aretenir}[1][]{popbase,colframe=popgreen,before upper={\popboxhead{À retenir}{popgreen}{#1}}}
\newtcolorbox{exemplebox}[1][]{popbase,colframe=poporange,before upper={\popboxhead{Exemple}{poporange}{#1}}}
\newtcolorbox{definition}[1][]{popbase,colframe=popblue,before upper={\popboxhead{Définition}{popblue}{#1}}}
\newtcolorbox{propriete}[1][]{popbase,colframe=poppurple,before upper={\popboxhead{Propriété}{poppurple}{#1}}}
\newtcolorbox{methode}[1][]{popbase,colframe=popgold,before upper={\popboxhead{Méthode}{popgold}{#1}}}
\newtcolorbox{attention}[1][]{popbase,colframe=poppink,before upper={\popboxhead{Attention}{poppink}{#1}}}
\newtcolorbox{experience}[1][]{popbase,colframe=popblue,colback=popblueL,before upper={\popboxhead{Expérience}{popblue}{#1}}}
% « Le savais-tu ? » : carte violette pleine (contexte, culture, Cameroun)
\newtcolorbox{savaistu}[1][]{popbase,colback=poppurple,colframe=popink,
  fontupper=\popsemi\fontsize{10.4}{14.2}\selectfont\color{white},
  before upper={\poppill{Le savais-tu\,?}{white}{poppurple}\ifx\relax#1\relax\else\hspace{6pt}{\popxbold\color{white}#1}\fi\par\vspace{7pt}}}
% Exercice résolu : énoncé en haut, \tcblower puis la solution sur fond crème
\newcounter{popexo}\newcounter{popexr}
\newtcolorbox{exoresolu}[1][]{popbase,colframe=poporange,colbacklower=popcream,
  segmentation style={popink,line width=1.2pt,dashed},
  before upper={\stepcounter{popexr}\popboxhead{Exercice résolu \thepopexr}{poporange}{#1}},
  before lower={\poppill{Solution}{popink}{popcream}\par\vspace{6pt}}}
% Exercice (non résolu en place) — la correction est dans la partie Corrigés
\newtcolorbox{exercice}[1][]{popbase,colframe=popink,
  before upper={\stepcounter{popexo}\popboxhead{Exercice \thepopexo}{popink}{#1}}}
% Corrigé
\newtcolorbox{corrige}[1][]{popbase,colframe=popgreen,colback=popgreenL,
  before upper={\popboxhead{Corrigé}{popgreen}{#1}}}
% Objectifs / prérequis / bilan
\newtcolorbox{objectifs}{popbase,colframe=popink,colback=poporangeL,
  before upper={\popboxhead{Objectifs de la leçon}{popink}{}}}
\newtcolorbox{prerequis}{popbase,colframe=popink,colback=popblueL,
  before upper={\popboxhead{Prérequis}{popblue}{}}}
\newtcolorbox{bilan}{popbase,colframe=popink,colback=white,boxrule=2.8pt,
  before upper={\popboxhead{Fiche bilan}{poporange}{L'essentiel à connaître pour l'examen}}}
% Figure encadrée avec légende
\newtcolorbox{popfigure}[1][]{enhanced,breakable=false,colback=white,colframe=popline,boxrule=1.4pt,arc=8pt,
  left=10pt,right=10pt,top=10pt,bottom=8pt,before skip=10pt,after skip=12pt,halign=center,
  lower separated=false,halign lower=center,
  fontlower=\popsemi\fontsize{9}{11.5}\selectfont\color{popmuted},
  #1}
\newcommand{\legende}[1]{\par\vspace{6pt}{\popsemi\fontsize{9}{11.5}\selectfont\color{popmuted}#1}}

% ============================================================
% Signature OnBuch+
% ============================================================
\newcommand{\poplogoinline}[1]{{\poplogo\fontsize{#1}{#1}\selectfont\color{popink}OnBuch\color{poporange}+}}
\newcommand{\signatureonbuch}{%
  \begin{tcolorbox}[enhanced,colback=popink,colframe=popink,boxrule=2.2pt,arc=10pt,
      drop shadow=poporange,left=14pt,right=14pt,top=10pt,bottom=10pt,before skip=0pt,after skip=0pt]
    \tikz[baseline=-0.7ex]\node[circle,fill=popgreen,draw=popcream,line width=1.3pt,minimum size=22pt,inner sep=0]{\popcheck{white}};%
    \hspace{9pt}\begin{minipage}[c]{0.62\linewidth}
      {\popxbold\fontsize{12}{13}\selectfont\color{popcream}Signé\ \textbullet\ {\poplogo OnBuch\color{poporange}+}}\par\vspace{2pt}
      {\raggedright\popsemi\fontsize{8}{10.5}\selectfont\color{popline}\MakeUppercase{Équipe pédagogique OnBuch+}\\\MakeUppercase{Conforme au programme officiel MINESEC}\par}
    \end{minipage}\hfill
    {\poplogo\fontsize{22}{22}\selectfont\color{popcream}OnBuch\color{poporange}+}
  \end{tcolorbox}%
}

% ============================================================
% Page de garde
% #1 titre · #2 matière · #3 niveau · #4 module · #5 n° leçon · #6 durée ·
% #7 description / objectifs (texte ou itemize)
% ============================================================
\newcommand{\pagegarde}[7]{%
  \thispagestyle{empty}
  \newgeometry{a4paper,margin=1.7cm,top=1.6cm,bottom=1.4cm}%
  \begin{tikzpicture}[remember picture,overlay]
    \fill[poporange!18] ([xshift=-3.2cm,yshift=-3.6cm]current page.north east) circle (4.6cm);
    \fill[poppurple!14] ([xshift=2.4cm,yshift=7.6cm]current page.south west) circle (3.8cm);
  \end{tikzpicture}%
  {\poplogo\fontsize{30}{30}\selectfont\color{popink}OnBuch\color{poporange}+}\hfill
  \raisebox{0.6ex}{\poppill{Cours signé \textbullet\ Premium}{popink}{popcream}}\par
  \vspace{1pt}\popsquiggle{poppurple}\par
  \vspace{8pt}
  \begin{tcolorbox}[enhanced,colback=poporange,colframe=popink,boxrule=2.8pt,arc=13pt,
      drop shadow=popink,left=18pt,right=18pt,top=12pt,bottom=13pt,after skip=10pt]
    \poppill{#2 \textbullet\ #3}{popink}{popcream}\hspace{5pt}\poppill{#4}{white}{popink}\par\vspace{8pt}
    {\popdisplay\fontsize{14}{14}\selectfont\color{popink}LEÇON #5}\par\vspace{4pt}
    {\raggedright\hyphenpenalty=10000\exhyphenpenalty=10000\popdisplay\fontsize{25}{27}\selectfont\color{popcream}\MakeUppercase{#1}\par}
    \vspace{8pt}
    {\popxbold\fontsize{9.5}{9.5}\selectfont\color{popink}\MakeUppercase{Durée conseillée : #6}}
  \end{tcolorbox}
  \begin{tcolorbox}[enhanced,colback=white,colframe=popink,boxrule=2.2pt,arc=10pt,
      drop shadow=popink,left=16pt,right=16pt,top=10pt,bottom=10pt,after skip=8pt]
    \poppill{Dans cette leçon}{poppurple}{white}\par\vspace{5pt}
    \popsemi\fontsize{9.3}{12.6}\selectfont\color{popink}#7
  \end{tcolorbox}
  \noindent\begin{minipage}[t]{0.487\linewidth}
    \begin{tcolorbox}[enhanced,colback=popgreen,colframe=popink,boxrule=2.2pt,arc=10pt,
        drop shadow=popink,left=11pt,right=11pt,top=10pt,bottom=10pt,height=3.1cm,valign=center]
      \tcbox[colback=white,colframe=popink,boxrule=1.3pt,arc=5pt,boxsep=0pt,nobeforeafter,
        left=3pt,right=3pt,top=3pt,bottom=3pt,valign=center]{\qrcode[height=1.9cm]{https://play.google.com/store/apps/details?id=com.luvvix.onbuch&hl=fr}}%
      \hspace{8pt}\begin{minipage}[c]{0.5\linewidth}
        {\popdisplay\fontsize{11}{12.5}\selectfont\color{white}\MakeUppercase{Télécharge OnBuch}}\par\vspace{4pt}
        {\raggedright\popsemi\fontsize{8.4}{11}\selectfont\color{white}Gratuit \textbullet\ Google Play\\Tuteur IA, annales, concours, résultats\par}
      \end{minipage}
    \end{tcolorbox}
  \end{minipage}\hfill
  \begin{minipage}[t]{0.487\linewidth}
    \begin{tcolorbox}[enhanced,colback=poppurple,colframe=popink,boxrule=2.2pt,arc=10pt,
        drop shadow=popink,left=11pt,right=11pt,top=10pt,bottom=10pt,height=3.1cm,valign=center]
      \tikz[baseline=-0.7ex]\node[circle,fill=white,draw=popink,line width=1.3pt,minimum size=1.6cm,inner sep=0]{\popdisplay\fontsize{24}{24}\selectfont\color{poppurple}+};%
      \hspace{8pt}\begin{minipage}[c]{0.6\linewidth}
        {\popdisplay\fontsize{11}{12.5}\selectfont\color{white}\MakeUppercase{Passe à OnBuch+}}\par\vspace{4pt}
        {\raggedright\popsemi\fontsize{8.4}{11}\selectfont\color{white}Tous les cours signés, Tuteur IA illimité, fiches PDF — onglet OnBuch+ de l'app\par}
      \end{minipage}
    \end{tcolorbox}
  \end{minipage}
  \vspace{8pt}
  \popcontenu
  \vfill
  \signatureonbuch
  \restoregeometry
  \newpage
}

% Rangée « Ce cours contient » (page de garde)
\newcommand{\poptile}[3]{% #1 couleur #2 gros texte #3 légende
  \begin{tcolorbox}[enhanced,colback=#1,colframe=popink,boxrule=2pt,arc=9pt,shadow={2.5pt}{-2.5pt}{0pt}{popink},
      left=6pt,right=6pt,top=9pt,bottom=9pt,halign=center,valign=center,height=1.75cm,width=\linewidth,nobeforeafter]
    {\popdisplay\fontsize{12.5}{13}\selectfont\color{white}\MakeUppercase{#2}}\par\vspace{3pt}
    {\centering\popsemi\fontsize{7.8}{9.5}\selectfont\color{white}#3\par}
  \end{tcolorbox}}
\newcommand{\popcontenu}{%
  \par\noindent{\popxboldls\fontsize{7.5}{7.5}\selectfont\color{popmuted}CE COURS SIGNÉ CONTIENT}\par\vspace{7pt}
  \noindent\begin{minipage}[t]{0.235\linewidth}\poptile{popblue}{Cours}{complet, illustré, pas à pas}\end{minipage}\hfill
  \begin{minipage}[t]{0.235\linewidth}\poptile{poporange}{Méthodes}{et exercices résolus}\end{minipage}\hfill
  \begin{minipage}[t]{0.235\linewidth}\poptile{poppink}{Exercices}{type examen + corrigés}\end{minipage}\hfill
  \begin{minipage}[t]{0.235\linewidth}\poptile{popgreen}{Fiche bilan}{l'essentiel à retenir}\end{minipage}\par}

% Sommaire
\newtcolorbox{sommaire}{enhanced,colback=white,colframe=popink,boxrule=2.2pt,arc=10pt,
  drop shadow=popink,left=18pt,right=18pt,top=13pt,bottom=13pt,before skip=6pt,after skip=20pt,
  before upper={\poppill{Sommaire}{poppurple}{white}\par\vspace{9pt}\popsemi\fontsize{10.6}{14.5}\selectfont\color{popink}}}

% Signature de fin de cours — poussée tout en bas de la dernière page
\newcommand{\findecours}{%
  \par\vfill\nopagebreak
  \begin{center}\popsquiggle{poporange}\end{center}\vspace{4pt}
  \signatureonbuch
}
\AtBeginDocument{\def\llbracket{\mathopen{[\mkern-3.5mu[}}\def\rrbracket{\mathclose{]\mkern-3.5mu]}}} % intervalles d'entiers (glyphe ⟦ absent de la police)
% Glyphes absents de Fira Math : redéfinis à partir de glyphes disponibles
\AtBeginDocument{%
  \def\setminus{\mathbin{\backslash}}\let\smallsetminus\setminus
  \def\vdots{\mathord{\vbox{\baselineskip3.2pt\lineskiplimit0pt\kern4pt\hbox{.}\hbox{.}\hbox{.}}}}%
  \def\ddots{\mathinner{\mkern1mu\raise6.5pt\hbox{.}\mkern2mu\raise3.6pt\hbox{.}\mkern2mu\raise0.7pt\hbox{.}\mkern1mu}}%
  \def\triangle{\mathord{\scalebox{0.9}{$\symup{\Delta}$}}}%
  \def\nabla{\mathord{\raisebox{\depth}{\scalebox{1}[-1]{$\symup{\Delta}$}}}}%
  \def\checkmark{\popcheck{popdark}}%
  \def\star{\mathbin{\ast}}\def\bigstar{\mathord{\ast}}%
  \def\diamond{\mathbin{\scalebox{0.6}{$\Diamond$}}}%
  \def\bigcup{\mathop{\mathchoice{\raisebox{-0.3em}{\scalebox{1.7}{$\cup$}}}{\scalebox{1.25}{$\cup$}}{\cup}{\cup}}\displaylimits}%
  \def\bigcap{\mathop{\mathchoice{\raisebox{-0.3em}{\scalebox{1.7}{$\cap$}}}{\scalebox{1.25}{$\cap$}}{\cap}{\cap}}\displaylimits}%
  \def\lesseqqgtr{\mathrel{\lessgtr}}\def\bigoplus{\mathop{\oplus}\displaylimits}%
}
\providecommand{\ssi}{si et seulement si} % abréviation fréquente
\AtBeginDocument{\providecommand{\wideparen}{\overparen}} % arc de cercle

%% FIGFIX-BEGIN v3 — correctifs globaux des figures (docs/onbuch-generation/tools/figfix)
\usepackage{adjustbox}\usepackage{etoolbox}
% toute figure TikZ/pgfplots plus large que la ligne est réduite à la largeur disponible
\BeforeBeginEnvironment{tikzpicture}{\begin{adjustbox}{max width=\linewidth}}
\AfterEndEnvironment{tikzpicture}{\end{adjustbox}}
% légendes pgfplots lisibles par-dessus les courbes
\pgfplotsset{every axis/.append style={legend style={fill=white,fill opacity=0.92,text opacity=1,draw=black!15,inner sep=3pt}}}
% étiquettes d'arêtes (syntaxe edge["..."]) avec fond blanc
\tikzset{every edge quotes/.append style={fill=white,inner sep=1.2pt}}
% bulles de cartes mentales : texte à la ligne dans la bulle, bulle un peu plus grande
\tikzset{every concept/.append style={text width=2.0cm,align=center,minimum size=2.3cm,inner sep=2pt}}
%% FIGFIX-END
\begin{document}

\pagegarde{Qu’est-ce que la Philosophie ?}{Philosophie}{1re Tech}{Philosophie – classe de Première technique}{N20}{5 séances (fiche de progression harmonisée)}{Cette leçon introductive définit la philosophie, retrace son origine grecque et précise son objet, sa méthode et son importance. Elle établit enfin le rapport entre philosophie et science, en montrant au technicien camerounais que philosopher est une démarche rationnelle utile à sa formation et à sa vie quotidienne.
\vspace{4pt}
\begin{multicols}{2}\small
\begin{itemize}
\item À la fin de cette leçon, je sais définir la philosophie à partir de son étymologie et de plusieurs définitions d'auteurs
\item À la fin de cette leçon, je sais situer l'origine historique de la philosophie (Grèce, VIe siècle av. J.-C.) et identifier ses conditions d'émergence
\item À la fin de cette leçon, je sais dégager l'objet propre de la philosophie (l'interrogation radicale sur le tout, l'homme, la connaissance, les valeurs)
\item À la fin de cette leçon, je sais décrire la méthode philosophique (doute, analyse, argumentation, critique) et la distinguer d'une simple opinion
\item À la fin de cette leçon, je sais justifier l'importance de la philosophie pour l'individu, la société et le technicien
\item À la fin de cette leçon, je sais comparer philosophie et science (points communs, différences, complémentarité)
\end{itemize}\end{multicols}}

\begin{sommaire}
\begin{multicols}{2}
\begin{enumerate}
\item Définitions de la philosophie
\item Origine et objet de la philosophie
\item La méthode de la philosophie
\item L'importance de la philosophie
\item Philosophie et science
\item TP guidé : philosopher sur des situations camerounaises
\item Activité d'intégration
\item Méthodes et exercices résolus
\item Exercices
\item Corrigés des exercices
\item Fiche bilan
\end{enumerate}
\end{multicols}
\end{sommaire}

\begin{objectifs}
\begin{itemize}
\item À la fin de cette leçon, je sais définir la philosophie à partir de son étymologie et de plusieurs définitions d'auteurs
\item À la fin de cette leçon, je sais situer l'origine historique de la philosophie (Grèce, VIe siècle av. J.-C.) et identifier ses conditions d'émergence
\item À la fin de cette leçon, je sais dégager l'objet propre de la philosophie (l'interrogation radicale sur le tout, l'homme, la connaissance, les valeurs)
\item À la fin de cette leçon, je sais décrire la méthode philosophique (doute, analyse, argumentation, critique) et la distinguer d'une simple opinion
\item À la fin de cette leçon, je sais justifier l'importance de la philosophie pour l'individu, la société et le technicien
\item À la fin de cette leçon, je sais comparer philosophie et science (points communs, différences, complémentarité)
\item À la fin de cette leçon, je sais appliquer la démarche philosophique à des situations concrètes de la vie courante au Cameroun
\item À la fin de cette leçon, je sais rédiger et justifier une réponse argumentée en suivant une démarche logique (thèse, arguments, exemples, conclusion)
\end{itemize}
\end{objectifs}

\begin{prerequis}
\begin{itemize}
\item Notions de culture générale sur la Grèce antique (histoire du collège)
\item Vocabulaire de base : opinion, vérité, raisonnement, argument
\item Notion de démarche scientifique vue en SVT/Physique (observation, hypothèse, vérification)
\item Capacité à lire et à commenter un texte court (français, Seconde)
\item Connaissance des grands repères chronologiques (avant/après J.-C.)
\end{itemize}
\end{prerequis}

\newpage

\coursec{Définitions de la philosophie}

\textbf{Situation d'accroche.} Lors d'un débat au foyer culturel de Yaoundé, un élève de Première affirme : « La philosophie, c'est seulement pour les intellectuels ; le technicien n'en a pas besoin. » Un autre réplique : « Pourtant, quand tu répares un moteur, tu cherches la cause de la panne, tu raisonnes, tu vérifies : n'est-ce pas déjà philosopher ? » Pendant ce temps, à l'hôpital de Bafoussam, un médecin hésite à annoncer une mauvaise nouvelle à un patient : faut-il toujours dire la vérité ? Ces situations quotidiennes montrent que la philosophie n'est pas un luxe réservé aux lettrés. Qu'est-ce donc que la philosophie, d'où vient-elle, et à quoi sert-elle réellement pour le technicien camerounais ?

Cette leçon répondra à ces questions. Mais avant de savoir d'où vient la philosophie et à quoi elle sert, il faut d'abord la \emph{définir}. Comment définir un concept philosophique ? C'est la première compétence que nous allons acquérir.

\begin{methode}[Comment définir un concept philosophique ?]
Pour définir rigoureusement un concept en philosophie, on procède en trois étapes :
\begin{enumerate}
\item \textbf{L'étymologie} : on remonte au sens premier du mot (souvent grec ou latin) pour dégager son sens originaire.
\item \textbf{Les définitions d'auteurs} : on compare les définitions proposées par les grands philosophes, qui éclairent chacun un aspect du concept.
\item \textbf{La synthèse personnelle} : on rassemble ces éléments en une définition complète, claire et justifiée, que l'on peut défendre dans une dissertation.
\end{enumerate}
Appliquons cette méthode au concept de \cle{philosophie}.
\end{methode}

\courssub{L'étymologie : l'amour de la sagesse}

\textbf{Intuition.} Quand on ne comprend pas un mot, le premier réflexe du bon élève est de le découper. « Philosophie » n'est pas un mot français d'origine : il vient du grec ancien. Le découper, c'est déjà commencer à philosopher, car c'est refuser d'accepter le mot comme une évidence.

Le mot \emph{philosophie} vient du grec \emph{philosophia}, formé de deux racines :
\begin{itemize}
\item \cle{philein}, qui signifie \emph{aimer}, \emph{désirer}, \emph{rechercher avec affection} ;
\item \cle{sophia}, qui signifie \emph{sagesse} ou \emph{savoir}.
\end{itemize}

\begin{definition}[La philosophie]
La \cle{philosophie} (du grec \emph{philosophia} : \emph{philein}, aimer, et \emph{sophia}, sagesse) est littéralement l'\textbf{amour de la sagesse}. Par extension, c'est une \textbf{démarche rationnelle et critique de recherche de la vérité} sur le monde, sur l'homme et sur les valeurs (le bien, le juste, le beau).
\end{definition}

Ce sens étymologique est précieux : il dit que le philosophe n'est pas celui qui \emph{possède} la sagesse, mais celui qui \emph{l'aime} et la \emph{cherche}. La philosophie est donc un mouvement, une quête, jamais une possession définitive.

\begin{popfigure}
\begin{center}
\begin{tikzpicture}[
  grow=down,
  level 1/.style={sibling distance=5.2cm, level distance=1.4cm},
  level 2/.style={sibling distance=2.9cm, level distance=1.5cm},
  every node/.style={draw, rounded corners=2pt, align=center, font=\small, inner sep=4pt},
  edge from parent/.style={draw=popmuted, -{Stealth[length=2mm]}}
]
\node[fill=popblue, text=white, font=\small\bfseries, align=center]{PHILOSOPHIA\\ (philosophie)}
  child { node[fill=popblueL, align=center]{\emph{philein}\\ aimer, rechercher}
    child { node[fill=popgreenL, align=center]{Étonnement\\ devant le monde} }
    child { node[fill=popgreenL, align=center]{Critique\\ des opinions} }
  }
  child { node[fill=poppurpleL, align=center]{\emph{sophia}\\ sagesse, savoir}
    child { node[fill=popgoldL, align=center]{Sagesse :\\ vérité et bonheur} }
  };
\end{tikzpicture}
\end{center}
\legende{Arbre étymologique du mot \emph{philosophia} : de la racine \emph{philein} (aimer) naissent l'étonnement et la critique ; de la racine \emph{sophia} (sagesse) naît la recherche de la vérité et du bonheur.}
\end{popfigure}

\courssub{Pythagore : le philosophe n'est pas le sage}

\textbf{Intuition.} Imagine un élève qui prétendrait tout savoir : on le trouverait prétentieux, et il cesserait d'apprendre. Celui qui sait qu'il ne sait pas tout reste, lui, ouvert à la recherche. C'est exactement l'attitude que le mot « philosophe » exprime.

Selon la tradition rapportée par les Anciens, c'est \cle{Pythagore} (VI\textsuperscript{e} siècle av. J.-C.) qui aurait inventé le mot \emph{philosophos}. Les Grecs appelaient \emph{sophos} (sage) celui qui possède le savoir. Pythagore, par humilité, refusait ce titre : seuls les dieux, disait-il, sont véritablement sages ; l'homme, lui, ne peut qu'\emph{aimer} et \emph{rechercher} la sagesse. Il se déclarait donc \emph{philosophe}, c'est-à-dire « ami de la sagesse ».

\begin{savaistu}[Pythagore, premier « philosophe » ?]
Pythagore de Samos (vers 580 -- vers 495 av. J.-C.), célèbre pour son théorème de géométrie, aurait été le premier à se dire \emph{philosophe} plutôt que \emph{sage}. Cette humilité devant le savoir est fondatrice : elle fait de la philosophie une recherche inachevée, un chemin plutôt qu'un point d'arrivée. Le même Pythagore illustre d'ailleurs le lien ancien entre mathématiques et philosophie : pour lui, les nombres étaient la clé de l'ordre du monde.
\end{savaistu}

Cette leçon d'humilité a une conséquence directe pour l'élève : on n'a jamais « fini » de philosopher. Même le bachelier, même le professeur, reste un \emph{chercheur} de sagesse.

\courssub{La définition par l'étonnement : Platon et Aristote}

\textbf{Intuition.} Pourquoi un enfant pose-t-il sans cesse des questions : « Pourquoi le ciel est-il bleu ? Pourquoi la nuit tombe-t-elle ? » Parce qu'il s'\emph{étonne}. L'adulte, habitué au monde, ne s'étonne plus. Or, disent Platon et Aristote, c'est précisément cet étonnement qui fait naître la philosophie.

Dans le \emph{Théétète}, \cle{Platon} (427--347 av. J.-C.) affirme que l'étonnement (\emph{thaumazein}) est le sentiment propre du philosophe : « la philosophie n'a pas d'autre origine ». De même, \cle{Aristote} (384--322 av. J.-C.), dans la \emph{Métaphysique}, explique que « c'est l'étonnement qui poussa les premiers penseurs à philosopher » : devant les phénomènes du monde (le jour et la nuit, les astres, la vie et la mort), l'homme s'est d'abord émerveillé, puis il a cherché à \emph{comprendre}.

\begin{aretenir}[L'étonnement, commencement de la philosophie]
Pour Platon et Aristote, philosopher commence par \textbf{s'étonner} devant ce qui semble aller de soi : le monde, la vie, la justice, le langage. L'étonnement n'est pas une ignorance qui subit : c'est une \emph{question} qui cherche. Il transforme l'évidence en problème.
\end{aretenir}

\begin{exemplebox}[L'étonnement du technicien]
Un élève en génie électrique à Douala remarque qu'un moteur chauffe anormalement. Il pourrait se contenter de dire « c'est comme ça ». Mais il s'étonne : \emph{pourquoi} chauffe-t-il ? Il mesure, compare, teste les hypothèses (mauvaise ventilation ? surcharge ? roulement usé ?). Cet étonnement qui devient recherche de la cause est le même mouvement que celui du philosophe devant le monde : refuser l'évidence pour chercher l'explication.
\end{exemplebox}

\courssub{La définition par la critique : Socrate et l'examen de soi}

\textbf{Intuition.} Nous répétons tous des opinions reçues : « on dit que... », « tout le monde sait que... ». La philosophie commence quand on cesse de répéter et qu'on commence à \emph{examiner} : cette opinion est-elle vraie ? Est-elle justifiée ?

\cle{Socrate} (470--399 av. J.-C.) est le maître de cette définition critique. Condamné à mort par le tribunal d'Athènes, il refuse de renoncer à philosopher et déclare devant ses juges : « Une vie sans examen ne vaut pas la peine d'être vécue. » Pour Socrate, philosopher, c'est \textbf{soumettre nos opinions, nos certitudes et nos valeurs à l'examen rationnel}. Sa méthode, le dialogue question-réponse, consiste à interroger ses concitoyens (qu'est-ce que la justice ? le courage ? la piété ?) pour révéler les contradictions de leurs réponses et les conduire à reconnaître leur ignorance : « Je sais que je ne sais rien. »

\begin{exemplebox}[Le marché Mokolo]
Au marché Mokolo à Yaoundé, le client avisé ne se contente jamais du premier prix annoncé : il questionne (« pourquoi si cher ? »), compare les étals, vérifie la qualité, négocie. De même, le philosophe ne se contente pas des opinions reçues : il les interroge, les compare, en exige les preuves. Accepter une idée sans l'examiner, c'est payer le premier prix sans marchander.
\end{exemplebox}

\courssub{La définition par la recherche de la vérité et du bonheur}

\textbf{Intuition.} À quoi sert de chercher la vérité ? Les philosophes de l'Antiquité répondent : à vivre bien. La philosophie n'est pas un jeu de l'esprit ; elle est un \emph{art de vivre} qui conduit au bonheur.

Pour \cle{Épicure} (341--270 av. J.-C.), la philosophie est une \emph{médecine de l'âme} : elle guérit les peurs inutiles (peur des dieux, peur de la mort) et les désirs désordonnés qui rendent malheureux. « Que nul, quand il est jeune, ne remette à plus tard de philosopher », écrit-il, car philosopher, c'est apprendre à être heureux par la raison. Pour les \cle{stoïciens} (comme Épictète et Sénèque), la philosophie enseigne à distinguer ce qui dépend de nous (nos jugements, nos actions) de ce qui n'en dépend pas (la maladie, la mort, l'opinion d'autrui), afin d'atteindre la tranquillité de l'âme. Dans les deux cas, la philosophie est définie comme \textbf{recherche de la vérité en vue du bonheur et de la sagesse pratique}.

\courssub{Les définitions des grands auteurs}

Chaque époque a reformulé la définition de la philosophie. Retenons cinq définitions essentielles :

\begin{itemize}
\item \textbf{Platon} : la philosophie est l'\emph{ascension} de l'âme, qui s'arrache aux apparences sensibles (le monde des ombres de la caverne) pour contempler les vérités éternelles, en particulier l'Idée du Bien.
\item \textbf{Aristote} : la philosophie est la \emph{science des premières causes et des premiers principes} ; elle naît de l'étonnement et recherche le « pourquoi » de toutes choses.
\item \textbf{Descartes} (1596--1650) : la philosophie est la recherche méthodique de la vérité par la raison ; dans le \emph{Discours de la méthode}, il définit la bonne méthode : « bien conduire sa raison et chercher la vérité dans les sciences », en commençant par le doute qui écarte les préjugés.
\item \textbf{Kant} (1724--1804) : la philosophie est l'\emph{examen critique de la raison par elle-même} ; elle répond à quatre questions : « Que puis-je savoir ? Que dois-je faire ? Que m'est-il permis d'espérer ? Qu'est-ce que l'homme ? » Philosopher, pour Kant, c'est aussi oser penser par soi-même : « Aie le courage de te servir de ton propre entendement ! »
\item \textbf{Hegel} (1770--1831) : la philosophie est « son temps saisi dans la pensée » ; elle est la réflexion la plus haute qu'une époque produit sur elle-même, « l'oiseau de Minerve [qui] prend son envol au crépuscule » : elle comprend le monde après coup, en le pensant.
\end{itemize}

\begin{popfigure}
\begin{center}
\begin{tikzpicture}[scale=1, transform shape]
\draw[popline, line width=1.2pt, -{Stealth[length=3mm]}] (0,0) -- (15.2,0);
\foreach \x/\date/\nom/\couleur in {
  1/{VI\textsuperscript{e} s. av. J.-C.}/{Pythagore}/popgreen,
  3.6/{470--399 av. J.-C.}/{Socrate}/popblue,
  5.6/{427--347 av. J.-C.}/{Platon}/poppurple,
  7.6/{384--322 av. J.-C.}/{Aristote}/poporange,
  10.6/{1596--1650 ap. J.-C.}/{Descartes}/popgold,
  13.2/{1724--1804 ap. J.-C.}/{Kant}/popink}
{
  \fill[\couleur] (\x,0) circle (0.09);
  \draw[\couleur] (\x,0) -- (\x,0.45);
  \node[above, align=center, font=\scriptsize, text=\couleur] at (\x,0.45) {\textbf{\nom}\\ \date};
}
\node[below, font=\scriptsize\itshape, text=popmuted] at (7.6,-0.25) {Antiquité grecque \hspace{3.5cm} Temps modernes};
\end{tikzpicture}
\end{center}
\legende{Frise chronologique des auteurs cités. L'échelle n'est pas proportionnelle.}
\end{popfigure}

\courssub{Synthèse : une définition complète de la philosophie}

Rassemblons les éléments dégagés. L'étymologie donne l'\emph{amour de la sagesse} ; Pythagore ajoute l'\emph{humilité} du chercheur ; Platon et Aristote, l'\emph{étonnement} ; Socrate, l'\emph{examen critique} ; Épicure et les stoïciens, la \emph{recherche du bonheur} ; Descartes et Kant, la \emph{rigueur rationnelle} ; Hegel, la \emph{réflexion sur son époque}.

\begin{aretenir}[Définition de synthèse]
La \cle{philosophie} est une \textbf{démarche rationnelle, critique et universelle}, née de l'étonnement, qui interroge \textbf{le monde} (qu'est-ce que la réalité ?), \textbf{l'homme} (qui suis-je ? que puis-je savoir ?) et \textbf{les valeurs} (qu'est-ce que le bien, le juste, le vrai ?), en vue de la vérité et d'une vie éclairée. Elle ne possède pas la sagesse : elle la recherche.
\end{aretenir}

\begin{attention}[Erreur fréquente : la « philosophie de vie »]
Au sens courant, on dit « c'est ma philosophie » pour exprimer une simple attitude de résignation ou de fatalisme (« il faut prendre la vie comme elle vient »). Ce n'est \textbf{pas} la philosophie au sens du cours : la philosophie est une \textbf{démarche rigoureuse} de questionnement et d'argumentation rationnelle, non une attitude passive. À l'examen, ne définis jamais la philosophie comme une « manière de voir la vie » sans préciser la rationalité critique qui la constitue.
\end{attention}

\textbf{Complément Bac : sagesse, savoir et opinion (\emph{doxa}).} Pour affiner la définition, distinguons trois termes que l'étymologie rapproche. Le \textbf{savoir} est un ensemble de connaissances vraies et démontrées (comme en mathématiques ou en mécanique). La \textbf{sagesse} est plus large : c'est un savoir \emph{mis au service de la vie}, une manière éclairée d'exister et d'agir. L'\textbf{opinion} (en grec \emph{doxa}) est une croyance non examinée, variable d'une personne à l'autre : « je pense que... », « on dit que... ». La philosophie part de l'opinion (qu'elle critique), passe par le savoir (qu'elle interroge dans ses fondements) et vise la sagesse (qu'elle ne possède jamais totalement). Ainsi, le philosophe n'est ni l'homme d'opinion, ni le simple savant : il est le chercheur de sagesse.

\begin{exoresolu}[Définir la philosophie à partir de son étymologie]
\textbf{Énoncé.} Un sujet de dissertation demande : « La philosophie est-elle un savoir comme les autres ? » En introduction, on te demande de définir la philosophie. Rédige une définition argumentée en quatre à six lignes, en mobilisant l'étymologie et au moins deux auteurs.
\tcblower
\textbf{Solution détaillée.} On applique la méthode en trois étapes.
\begin{enumerate}
\item \emph{Étymologie} : le mot « philosophie » vient du grec \emph{philosophia} (\emph{philein}, aimer ; \emph{sophia}, sagesse) : la philosophie est littéralement l'amour de la sagesse.
\item \emph{Auteurs} : selon la tradition, Pythagore aurait forgé ce mot pour marquer que l'homme ne possède pas la sagesse mais la recherche. Pour Platon et Aristote, cette recherche naît de l'étonnement devant le monde ; pour Socrate, elle consiste dans l'examen critique de nos opinions (« une vie sans examen ne vaut pas la peine d'être vécue »).
\item \emph{Synthèse} : la philosophie est donc une démarche rationnelle et critique de recherche de la vérité sur le monde, l'homme et les valeurs. Contrairement aux sciences, qui possèdent des savoirs démontrés, elle reste une \emph{quête} : ce qui ouvre le problème du sujet (est-elle alors un savoir ?).
\end{enumerate}
\textbf{Rédaction possible} : « Du grec \emph{philosophia} (\emph{philein}, aimer, et \emph{sophia}, sagesse), la philosophie signifie l'amour de la sagesse. Pythagore, à qui l'on attribue ce mot, entendait marquer l'humilité du philosophe, qui recherche la sagesse sans la posséder. Née de l'étonnement selon Platon et Aristote, elle est pour Socrate un examen critique de nos opinions. On peut ainsi la définir comme une démarche rationnelle et critique de recherche de la vérité sur le monde, l'homme et les valeurs. Mais une quête qui ne possède jamais son objet peut-elle encore être appelée un savoir ? »
\end{exoresolu}

\begin{exercice}[À toi de philosopher]
\begin{enumerate}
\item Donne l'étymologie du mot « philosophie » et explique en quoi elle distingue le philosophe du sage.
\item Pourquoi Socrate affirme-t-il qu'« une vie sans examen ne vaut pas la peine d'être vécue » ? Illustre ta réponse par une situation de la vie courante au Cameroun.
\item Un camarade déclare : « Ma philosophie, c'est de ne jamais me fatiguer pour rien. » S'agit-il de philosophie au sens du cours ? Justifie ta réponse.
\end{enumerate}
\end{exercice}

\begin{corrige}[Exercice : À toi de philosopher]
\begin{enumerate}
\item \emph{Philosophia} vient de \emph{philein} (aimer) et \emph{sophia} (sagesse). Le sage (\emph{sophos}) possède le savoir ; le philosophe, lui, l'aime et le recherche sans le posséder pleinement : c'est l'humilité fondatrice attribuée à Pythagore.
\item Pour Socrate, vivre sans examiner ses opinions, c'est vivre comme un dormeur, guidé par les préjugés et les « on-dit ». L'examen rationnel rend la vie humainement digne. Exemple : l'électeur camerounais qui examine les promesses des candidats au lieu de voter par habitude « examine » sa vie civique ; l'élève qui vérifie une rumeur avant de la partager sur WhatsApp pratique déjà l'examen socratique.
\item Non : il s'agit d'une « philosophie de vie » au sens courant, c'est-à-dire d'une attitude personnelle (ici, paresse ou résignation). La philosophie au sens propre est une démarche rationnelle et critique de recherche de la vérité ; elle exige questionnement, argumentation et rigueur, ce que cette simple attitude ne comporte pas.
\end{enumerate}
\end{corrige}

Définie comme amour de la sagesse, étonnement, examen critique et recherche du bonheur, la philosophie apparaît déjà comme une activité exigeante et universelle. Mais d'où vient-elle historiquement, et quel est précisément son objet ? C'est ce que nous verrons dans la section suivante.

\coursec{Origine et objet de la philosophie}

Dans la section précédente, nous avons défini la philosophie comme l'amour de la sagesse et comme un discours rationnel et critique. Mais une question demeure : \emph{d'où vient-elle ?} Pourquoi la philosophie est-elle née à un moment précis, dans un lieu précis ? Et de quoi parle-t-elle exactement ? C'est à ces deux questions --- celle de l'\cle{origine} et celle de l'\cle{objet} de la philosophie --- que cette section est consacrée.

\courssub{La naissance de la philosophie en Grèce}

\textbf{Un événement daté et localisé}\par

Les historiens de la philosophie s'accordent à dater la naissance de la philosophie au \textbf{VI\textsuperscript{e} siècle avant Jésus-Christ}, dans les cités grecques d'Ionie, sur la côte occidentale de l'Asie Mineure (l'actuelle Turquie), et plus précisément à \cle{Milet}, port commercial prospère. Le premier penseur que la tradition retient est \textbf{Thalès de Milet} (vers 625 -- vers 547 av. J.-C.). Celui-ci chercha à expliquer les phénomènes naturels --- les tremblements de terre, les éclipses, la crue des fleuves --- non par l'intervention des dieux, mais par des causes naturelles : il affirmait par exemple que l'eau est le principe (\emph{archè}) de toutes choses.

\begin{propriete}[Naissance de la philosophie (admise)]
La philosophie naît du \cle{passage du mythe au logos} en Grèce, au VI\textsuperscript{e} siècle avant J.-C., dans les cités ioniennes (Milet), avant de s'épanouir à Athènes avec Socrate, Platon et Aristote. Cette propriété est admise : elle s'appuie sur les témoignages historiques et ne fait pas l'objet d'une démonstration exigible.
\end{propriete}

\begin{savaistu}[Thalès et l'éclipse de 585 av. J.-C.]
Thalès de Milet, considéré comme le premier philosophe, aurait prédit une éclipse solaire en 585 av. J.-C., qui aurait même interrompu une bataille entre Mèdes et Lydiens. Qu'il ait réellement fait ce calcul ou non, la légende signifie quelque chose d'essentiel : la \cle{raison} humaine est capable de prévoir ce que la superstition attribuait à la colère des dieux. Le ciel cesse d'être un théâtre de caprices divins pour devenir un objet de savoir.
\end{savaistu}

\begin{popfigure}
\begin{tikzpicture}[scale=1.0]
% Mer Méditerranée
\fill[popblueL] (0,0) rectangle (13,6);
% Europe (nord)
\fill[popcream] (0,3.6) -- (1.5,4.2) -- (2.5,4.0) -- (3.2,4.6) -- (4.0,4.2) -- (4.6,5.0) -- (5.5,4.6) -- (6.5,5.2) -- (7.5,4.8) -- (8.5,5.4) -- (10,5.0) -- (11.5,5.5) -- (13,5.2) -- (13,6) -- (0,6) -- cycle;
% Grèce (péninsule)
\fill[popcream] (5.6,4.6) -- (5.9,3.4) -- (6.3,3.0) -- (6.7,3.5) -- (6.9,4.4) -- cycle;
% Asie Mineure (Anatolie)
\fill[popcream] (7.6,4.8) -- (8.2,3.6) -- (9.5,3.4) -- (11,3.8) -- (13,3.6) -- (13,5.2) -- (11.5,5.5) -- (10,5.0) -- (8.5,5.4) -- cycle;
% Afrique (sud)
\fill[popcream] (0,0) -- (13,0) -- (13,1.6) -- (11,1.8) -- (9,1.4) -- (7,1.8) -- (5,1.3) -- (3,1.7) -- (1,1.2) -- (0,1.5) -- cycle;
% Villes
\fill[popink] (6.3,3.15) circle (0.09);
\node[above, font=\small\bfseries, popink] at (6.3,3.25) {Athènes};
\fill[poporange] (8.15,3.75) circle (0.11);
\node[below, font=\small\bfseries, poporange] at (8.15,3.65) {Milet};
% Mer Égée label
\node[font=\small\itshape, popblue] at (7.3,3.0) {mer Égée};
\node[font=\small\itshape, popblue] at (3.5,2.6) {mer Méditerranée};
% Routes commerciales
\draw[->, popgreen, thick, dashed] (8.15,3.75) .. controls (7.4,3.4) .. (6.35,3.2);
\draw[->, popgreen, thick, dashed] (8.15,3.75) .. controls (9.5,2.8) .. (11,1.95);
\draw[->, popgreen, thick, dashed] (6.35,3.2) .. controls (5,2.5) .. (3.2,1.85);
% Flèche du nord
\draw[->, popdark, thick] (12.3,5.6) -- (12.3,6.0);
\node[right, font=\small, popdark] at (12.35,5.8) {N};
\end{tikzpicture}
\legende{Carte schématique de la Méditerranée antique. Milet (point orange), cité d'Ionie sur la côte d'Asie Mineure, est le berceau de la philosophie au VI\textsuperscript{e} siècle av. J.-C. ; Athènes (point noir) en devient le centre au V\textsuperscript{e} siècle. Les pointillés verts figurent les routes commerciales maritimes, vecteurs d'échanges d'idées. Carte très simplifiée, sans précision géographique.}
\end{popfigure}

\courssub{Les conditions d'émergence de la philosophie}

La philosophie n'est pas née par hasard en Grèce. Plusieurs conditions historiques se sont rencontrées au même moment. Retenons-en quatre.

\begin{enumerate}
\item \textbf{La cité démocratique et la liberté de parole.} Dans les cités grecques, et notamment à Athènes, les affaires publiques se discutent sur la place publique (\emph{agora}). Les citoyens débattent, argumentent, votent les lois. Cette liberté de tout dire, la \cle{parrhèsia}, habitue les esprits à soumettre toute opinion à la discussion : on ne croit plus, on \emph{argumente}.
\item \textbf{Les échanges commerciaux en Méditerranée.} Milet est un port qui communique avec l'Égypte, la Phénicie, Babylone. Les marchands grecs découvrent d'autres coutumes, d'autres dieux, d'autres récits. Ce choc des cultures relativise les traditions : si chaque peuple a ses propres mythes, aucun ne peut prétendre détenir seul la vérité. Il faut alors chercher une vérité \emph{universelle}, valable pour tous.
\item \textbf{L'écriture alphabétique.} Simple et accessible (une vingtaine de signes), elle permet de fixer, de transmettre et surtout de \emph{critiquer} les discours. Un texte écrit peut être relu, comparé, réfuté : l'écriture est l'outil de l'argumentation.
\item \textbf{Le loisir (\emph{skholè}).} Aristote le souligne : pour philosopher, il faut du temps libre, libéré du travail de survie. Le mot grec \emph{skholè}, qui signifie « loisir studieux », a d'ailleurs donné notre mot « école » : l'école est le lieu du temps libéré pour penser.
\end{enumerate}

\begin{attention}[Ce qu'il ne faut pas conclure]
Dire que la philosophie est née en Grèce ne signifie pas que les Grecs ont tout inventé de rien. Ils ont emprunté aux Égyptiens (géométrie) et aux Babyloniens (astronomie). Leur originalité est ailleurs : ils ont transformé ces savoirs pratiques en \emph{recherche désintéressée de la vérité} et en discours argumenté.
\end{attention}

\courssub{Le passage du mythe au logos}

\textbf{L'intuition d'abord.} Un enfant demande : « Pourquoi il tonne ? » On peut lui répondre : « C'est le dieu Sango qui frappe sa pierre à tonner », ou bien : « C'est la décharge électrique entre nuages chargés ». La première réponse est un \cle{mythe} ; la seconde relève du \cle{logos}. La naissance de la philosophie, c'est le moment où les hommes décident de préférer la seconde réponse.

\begin{definition}[Mythe et logos]
Le \cle{logos} est le discours rationnel et argumenté, qui explique les phénomènes par des causes naturelles et des raisons que chacun peut examiner et discuter. Il s'oppose au \cle{mythe}, récit imaginaire mettant en scène des dieux ou des héros pour expliquer le monde, et qui demande à être \emph{cru} plutôt que démontré.
\end{definition}

Le passage du mythe au logos comporte trois traits essentiels :

\begin{itemize}
\item \textbf{La naturalisation des causes} : le tonnerre, les séismes, les éclipses s'expliquent par la nature (\emph{phusis}) et non par la volonté des dieux ;
\item \textbf{L'exigence de preuve} : toute affirmation doit pouvoir être justifiée devant d'autres esprits rationnels ;
\item \textbf{L'universalité} : la vérité cherchée vaut pour tous les hommes, pas seulement pour une tribu ou une cité.
\end{itemize}

\begin{attention}[Erreur fréquente]
Ne pas écrire que le mythe est « stupide » ou « inutile ». Le mythe a une fonction sociale (il fonde l'unité d'un groupe) et il contient souvent une sagesse symbolique. Le passage du mythe au logos n'est pas une suppression, mais un \emph{changement de régime de vérité} : on passe de la croyance à la démonstration. D'ailleurs, Platon lui-même utilise encore des mythes dans ses dialogues.
\end{attention}

\courssub{L'étonnement : le point de départ de la philosophie}

Pourquoi les hommes se mettent-ils à philosopher ? Platon et Aristote répondent d'une même voix : par l'\cle{étonnement} (en grec, \emph{thaumazein}). Dans le \emph{Théétète}, Platon écrit que l'étonnement est « le sentiment du philosophe » ; Aristote, au début de la \emph{Métaphysique}, affirme que « c'est l'étonnement qui poussa les premiers penseurs à philosopher ».

\begin{exemplebox}[Retrouver l'étonnement]
Un enfant s'étonne de tout : « Pourquoi le ciel est-il bleu ? Où vais-je quand je dors ? Pourquoi faut-il obéir ? » L'adulte, lui, s'habitue : il ne voit plus rien d'étrange dans le fait que le soleil se lève, que les promesses engagent ou que les hommes meurent. Le philosophe est celui qui \emph{refuse l'habitude} et retrouve le regard de l'enfant : ce qui semblait aller de soi redevient un problème. L'étonnement n'est pas l'ignorance qui subit : c'est l'ignorance qui \emph{questionne} et qui veut savoir.
\end{exemplebox}

L'étonnement est donc à la fois un \emph{sentiment} (l'admiration mêlée d'inquiétude devant le monde) et un \emph{moteur} (il déclenche la recherche). Sans étonnement, pas de question ; sans question, pas de philosophie.

\courssub{L'objet de la philosophie : l'interrogation sur le tout}

De quoi la philosophie s'occupe-t-elle ? Sa réponse est vertigineuse : \emph{de tout}. Alors que chaque science découpe un domaine particulier (la biologie étudie le vivant, la physique la matière), la philosophie interroge \cle{le tout} : le monde dans son ensemble, l'être en général, et la place de l'homme dans cet ensemble. Kant a résumé cet objet en quatre questions célèbres :

\begin{enumerate}
\item \textbf{Que puis-je savoir ?} --- question de la \cle{connaissance} : quelles sont les sources, les limites et la valeur de notre savoir ? Puis-je me fier à mes sens ? À ma raison ?
\item \textbf{Que dois-je faire ?} --- question de l'\cle{action} : qu'est-ce qu'agir bien ? Le bonheur est-il le but de la vie ? Faut-il toujours dire la vérité ?
\item \textbf{Que m'est-il permis d'espérer ?} --- question du sens de l'existence et du destin humain ;
\item \textbf{Qu'est-ce que l'homme ?} --- question qui rassemble les trois précédentes : \emph{qui suis-je ?} Suis-je un corps, une âme, les deux ? Suis-je libre ou déterminé ?
\end{enumerate}

On peut regrouper ces interrogations autour de trois grands pôles : le \textbf{monde} (d'où vient l'univers ? qu'est-ce que l'être ?), l'\textbf{homme} (qui suis-je ? suis-je libre ?) et les \textbf{valeurs} (qu'est-ce que le \cle{bien}, le \cle{juste}, le \cle{beau} ?). La philosophie ne se contente pas de décrire ce qui est : elle demande aussi ce qui \emph{vaut} et ce qui \emph{doit être}.

\courssub{Les grandes branches de la philosophie}

Cet immense objet se divise en plusieurs disciplines, chacune attachée à une question centrale :

\begin{itemize}
\item La \cle{métaphysique} (ou ontologie) : qu'est-ce que l'être ? Le monde a-t-il une cause ? Dieu existe-t-il ?
\item L'\cle{épistémologie} (ou théorie de la connaissance) : qu'est-ce que savoir ? Qu'est-ce qu'une vérité scientifique ?
\item La \cle{morale} (ou éthique) : qu'est-ce que le bien ? Comment dois-je agir ?
\item La philosophie \cle{politique} : qu'est-ce qu'une société juste ? D'où vient l'autorité de l'État ?
\item L'\cle{esthétique} : qu'est-ce que le beau ? Qu'est-ce que l'art ?
\item La \cle{logique} : qu'est-ce qu'un raisonnement valide ? Comment distinguer un argument solide d'un sophisme ?
\end{itemize}

\begin{popfigure}
\begin{tikzpicture}[scale=1.0]
% Camembert : 6 secteurs de 60°
\fill[popblueL]   (0,0) -- (90:2.4) arc (90:150:2.4) -- cycle;
\fill[poporangeL] (0,0) -- (150:2.4) arc (150:210:2.4) -- cycle;
\fill[popgreenL]  (0,0) -- (210:2.4) arc (210:270:2.4) -- cycle;
\fill[poppurpleL] (0,0) -- (270:2.4) arc (270:330:2.4) -- cycle;
\fill[poppinkL]   (0,0) -- (330:2.4) arc (330:30:2.4) -- cycle;
\fill[popgoldL]   (0,0) -- (30:2.4) arc (30:90:2.4) -- cycle;
\draw[popdark, thick] (0,0) circle (2.4);
\foreach \a in {30,90,150,210,270,330} \draw[popdark] (0,0) -- (\a:2.4);
% Étiquettes
\node[font=\small\bfseries, popdark] at (120:1.55) {Métaphysique};
\node[font=\scriptsize, popdark, text width=2.6cm, align=center] at (120:3.35) {Qu'est-ce que l'être ?};
\node[font=\small\bfseries, popdark] at (180:1.5) {Épistémologie};
\node[font=\scriptsize, popdark, text width=2.8cm, align=center] at (180:3.45) {Que puis-je savoir ?};
\node[font=\small\bfseries, popdark] at (240:1.55) {Morale};
\node[font=\scriptsize, popdark, text width=2.6cm, align=center] at (240:3.35) {Que dois-je faire ?};
\node[font=\small\bfseries, popdark] at (300:1.55) {Politique};
\node[font=\scriptsize, popdark, text width=2.8cm, align=center] at (300:3.35) {Qu'est-ce que le juste ?};
\node[font=\small\bfseries, popdark] at (0:1.55) {Esthétique};
\node[font=\scriptsize, popdark, text width=2.6cm, align=center] at (0:3.3) {Qu'est-ce que le beau ?};
\node[font=\small\bfseries, popdark] at (60:1.55) {Logique};
\node[font=\scriptsize, popdark, text width=2.8cm, align=center] at (60:3.35) {Qu'est-ce que bien raisonner ?};
\end{tikzpicture}
\legende{Les six grandes branches de la philosophie et leur question centrale. Chaque branche découpe un domaine de l'interrogation philosophique, mais toutes convergent vers une même exigence : penser par soi-même, de manière rationnelle et argumentée.}
\end{popfigure}

\courssub{La philosophie africaine et camerounaise}

\textbf{Le problème.} Si le mot « philosophie » est grec, faut-il en conclure que la pensée rationnelle et critique est l'apanage de l'Occident ? Certains auteurs coloniaux l'ont prétendu, niant aux Africains toute capacité d'abstraction. Cette thèse est aujourd'hui réfutée.

\begin{attention}[Erreur fréquente à l'examen]
Erreur fréquente : affirmer que la philosophie n'existe qu'en Occident. Il existe des traditions de pensée critique sur \emph{tous} les continents, même si le mot « philosophie » est grec. L'Afrique possède de riches traditions de sagesse (proverbes, contes, palabres) et une philosophie universitaire contemporaine reconnue (penseurs comme Paulin Hountondji, Marcien Towa, Fabien Eboussi Boulaga).
\end{attention}

En Afrique et au Cameroun, la réflexion critique s'est transmise notamment par :

\begin{itemize}
\item les \textbf{proverbes}, qui condensent une sagesse pratique et invitent à l'interprétation (« Quand un vieillard meurt, c'est une bibliothèque qui brûle ») ;
\item la \textbf{palabre}, institution du dialogue public où la parole circule pour résoudre les conflits et délibérer sur le bien commun.
\end{itemize}

\begin{exemplebox}[La palabre villageoise au Cameroun]
Dans de nombreux villages camerounais, lorsqu'un conflit éclate (litige foncier, dispute familiale), les anciens et les parties se réunissent sous l'\cle{arbre à palabres}. Chacun expose ses raisons, les témoins sont entendus, les sages interrogent, confrontent les arguments, puis la communauté recherche une décision juste et acceptée de tous. Cette pratique illustre une forme africaine du \emph{dialogue rationnel} : on y recherche la vérité et la justice par la parole argumentée, et non par la force --- démarche profondément proche de l'esprit philosophique né sur l'agora grecque.
\end{exemplebox}

\begin{savaistu}[Le débat sur l'ethnophilosophie --- complément Bac]
Au XX\textsuperscript{e} siècle, un débat a agité la philosophie africaine. Des ouvrages comme celui de Placide Tempels (\emph{Philosophie bantoue}, 1945) présentaient la philosophie africaine comme une sagesse \emph{collective} et anonyme, inscrite dans les coutumes : c'est ce qu'on appelle l'\cle{ethnophilosophie}. Des philosophes comme Paulin Hountondji (Bénin) et Marcien Towa (Cameroun) ont critiqué cette conception : la philosophie, disent-ils, suppose un discours \emph{écrit, individuel et critique}, assumé par des auteurs qui débattent entre eux. La sagesse collective en serait la matière, mais non encore la forme accomplie. Ce débat --- la philosophie est-elle nécessairement écrite et individuelle ? --- est un excellent sujet de réflexion pour le Baccalauréat technique.
\end{savaistu}

\begin{exoresolu}[Mythe, logos ou palabre ?]
\textbf{Énoncé.} Pour chacun des énoncés suivants, indique s'il relève du mythe ou du logos, en justifiant ta réponse :
\begin{enumerate}
\item « La crue du Wouri est due à la colère des génies du fleuve. »
\item « La crue du Wouri s'explique par l'accumulation des pluies de saison dans le bassin versant et par l'ensablement du lit du fleuve ; on peut la prévoir en mesurant les précipitations. »
\item « Sous l'arbre à palabres, le notable demande à chaque partie de prouver ses droits sur le terrain contesté. »
\end{enumerate}
\tcblower
\textbf{Solution détaillée.}
\begin{enumerate}
\item Cet énoncé relève du \textbf{mythe} : il explique un phénomène naturel par l'intervention d'êtres surnaturels (les génies), et demande à être cru sans preuve vérifiable.
\item Cet énoncé relève du \textbf{logos} : il invoque des causes naturelles (pluies, ensablement) et propose une méthode de vérification (la mesure des précipitations), donc un discours rationnel, contrôlable et discutable.
\item Cette situation relève du \textbf{logos}, dans sa forme africaine de la palabre : le notable exige des \emph{preuves} et des \emph{raisons}, et la décision naît d'un débat argumenté. On retrouve l'exigence fondatrice de la philosophie : soumettre les opinions à la discussion rationnelle.
\end{enumerate}
\textbf{Conclusion} : le critère de distinction n'est pas géographique (Grèce contre Afrique) mais rationnel : un discours relève du logos dès lors qu'il argumente et accepte d'être discuté, où qu'il soit prononcé.
\end{exoresolu}

\begin{aretenir}[L'essentiel de la section]
\begin{itemize}
\item La philosophie naît au VI\textsuperscript{e} siècle av. J.-C. à \textbf{Milet} (Ionie), avec \textbf{Thalès}, grâce à la rencontre de quatre conditions : la cité démocratique et la \emph{parrhèsia}, les échanges commerciaux, l'écriture alphabétique et le loisir (\emph{skholè}).
\item Elle consiste en un \textbf{passage du mythe au logos} : expliquer le monde par la raison et la preuve, et non plus par les récits fabuleux.
\item Son point de départ est l'\textbf{étonnement} (\emph{thaumazein}), qui transforme l'évidence en problème.
\item Son objet est \textbf{l'interrogation sur le tout} : le monde, l'homme, la connaissance, l'action et les valeurs (bien, juste, beau), répartis en branches : métaphysique, épistémologie, morale, politique, esthétique, logique.
\item La pensée critique existe sur tous les continents : au Cameroun, les proverbes et la \textbf{palabre} témoignent d'une tradition du dialogue rationnel, et le débat sur l'ethnophilosophie interroge la définition même de la philosophie.
\end{itemize}
\end{aretenir}

\coursec{La méthode de la philosophie}

Nous avons vu, dans la section précédente, que la philosophie est née de l'\cle{étonnement} et qu'elle a pour objet la recherche d'une sagesse universelle, fondée sur la raison et non sur le mythe. Mais une question demeure : \emph{comment} le philosophe procède-t-il concrètement ? Quels outils utilise-t-il pour passer de l'étonnement à la vérité ? C'est la question de la \cle{méthode}. La philosophie ne se contente pas d'affirmer : elle examine, elle prouve, elle discute. Cette section présente les grands moments de la démarche philosophique : le doute méthodique, l'analyse, l'argumentation, le dialogue socratique, l'esprit critique, et la visée d'une vérité universelle.

\courssub{Le doute méthodique : suspendre ses certitudes pour examiner}

\textbf{Intuition.}\par Imagine un élève de Première technique à Yaoundé qui entend dire : « Les filles ne réussissent pas en spécialité F (mécanique) ». Il peut accepter cette idée sans réfléchir, comme beaucoup. Mais il peut aussi se dire : « Attends. Est-ce vraiment vrai ? Quelles preuves a-t-on ? » Ce geste de suspension du jugement est le premier acte du philosophe.

\begin{definition}[Le doute méthodique]
Le \cle{doute méthodique} est la démarche qui consiste à suspendre provisoirement toutes ses certitudes afin de les examiner une à une, pour ne retenir que ce qui résiste à l'épreuve de la raison. Il est \emph{provisoire} (on doute pour un temps) et \emph{méthodique} (on doute selon un plan, dans un but précis : trouver un fondement solide).
\end{definition}

C'est le philosophe français René \textbf{Descartes} (1596-1650) qui a donné au doute sa forme la plus célèbre. Dans ses \emph{Méditations}, il décide de faire table rase : tout ce qui peut être mis en doute — les sens qui parfois nous trompent (un bâton plongé dans l'eau paraît cassé), les rêves qui imitent la réalité, les opinions reçues de l'éducation — est provisoirement rejeté. Que reste-t-il alors ? Une vérité indubitable : « Je pense, donc je suis » (\emph{cogito ergo sum}). Même en doutant de tout, je ne peux pas douter que je suis en train de douter, donc que je pense, donc que j'existe. Sur cette première certitude, Descartes reconstruit l'édifice du savoir.

\begin{attention}[Ne pas confondre doute méthodique et scepticisme]
Erreur fréquente à l'examen : présenter le doute cartésien comme un refus définitif de toute vérité. C'est faux. Le \cle{scepticisme} doute de tout \emph{définitivement} et conclut qu'on ne peut rien savoir : c'est une position stérile. Le \cle{doute méthodique}, au contraire, est un \emph{outil provisoire} : on doute comme on démolit une maison bancale pour la reconstruire sur des fondations solides. Douter n'est pas une fin, c'est un moyen pour établir des vérités certaines.
\end{attention}

\courssub{L'analyse : décomposer le problème}

Une fois les certitudes suspendues, il faut examiner. Or un problème philosophique est souvent complexe, comme une machine qu'on ne comprend qu'en la démontant. C'est le travail de l'\cle{analyse}.

\begin{definition}[L'analyse]
L'\cle{analyse} est l'opération intellectuelle qui consiste à décomposer un problème en ses éléments simples : définir précisément les termes employés, distinguer les différentes questions posées, et repérer les \cle{présupposés} (ce qu'on admet sans l'avoir démontré).
\end{definition}

L'analyse comporte trois gestes :

\begin{enumerate}
\item \textbf{Définir les termes.} Avant de discuter « La liberté existe-t-elle ? », il faut s'entendre sur le mot \emph{liberté} : s'agit-il de faire ce qu'on veut, de choisir sans contrainte, ou d'agir selon la raison ? La plupart des disputes viennent de ce que les interlocuteurs n'emploient pas les mots dans le même sens.
\item \textbf{Décomposer la question.} « L'État est-il utile ? » contient plusieurs sous-questions : utile à qui ? à quoi (sécurité, éducation, justice) ? par rapport à quelle alternative ?
\item \textbf{Repérer les présupposés.} Quand quelqu'un demande « Pourquoi les jeunes d'aujourd'hui ne respectent-ils plus rien ? », il présuppose que les jeunes ne respectent rien — ce qui est précisément à prouver.
\end{enumerate}

\begin{exemplebox}[Analyser une question de la vie courante]
Question : « Faut-il interdire les téléphones portables au lycée ? » L'analyse révèle : (1) le terme \emph{interdire} peut signifier interdire en classe seulement ou dans toute l'enceinte ; (2) le problème se décompose : nuisance pour l'attention ? outil de triche ? mais aussi outil de recherche et lien avec les parents ; (3) le présupposé caché est que le portable est d'abord un instrument de distraction, ce qui mérite d'être discuté.
\end{exemplebox}

\courssub{L'argumentation : construire un raisonnement logique}

Analyser ne suffit pas : il faut ensuite \emph{prouver}. C'est le rôle de l'\cle{argumentation}.

\begin{definition}[L'argumentation]
L'\cle{argumentation} est l'ensemble des raisons avancées pour établir ou défendre une thèse. Un raisonnement logique se compose de \cle{prémisses} (les affirmations de départ) et d'une \cle{conclusion} (ce qu'on déduit des prémisses).
\end{definition}

Le modèle le plus ancien de raisonnement rigoureux est le \cle{syllogisme}, étudié par \textbf{Aristote}. Exemple célèbre :

\begin{itemize}
\item Prémisse 1 : Tous les hommes sont mortels.
\item Prémisse 2 : Or Socrate est un homme.
\item Conclusion : Donc Socrate est mortel.
\end{itemize}

Si les deux prémisses sont vraies et si la forme du raisonnement est correcte, la conclusion s'impose nécessairement à tout esprit raisonnable. C'est la force de la logique : elle ne dépend ni de l'autorité, ni de l'émotion, ni de la force physique.

\begin{popfigure}
\begin{center}
\begin{tikzpicture}[
  prem/.style={draw=popblue, fill=popblueL, rounded corners, text width=5.2cm, align=center, font=\small, thick},
  concl/.style={draw=poporange, fill=poporangeL, rounded corners, text width=6.5cm, align=center, font=\small\bfseries, thick},
  fleche/.style={-{Stealth[length=3mm]}, very thick, popdark}]
\node[prem, align=center] (p1) at (-2.9,0){\textbf{Prémisse 1}\\ Tous les hommes sont mortels};
\node[prem, align=center] (p2) at (2.9,0){\textbf{Prémisse 2}\\ Or Socrate est un homme};
\node[concl, align=center] (c) at (0,-3.2){\textbf{Conclusion}\\ Donc Socrate est mortel};
\draw[fleche] (p1.south) -- (c.north west);
\draw[fleche] (p2.south) -- (c.north east);
\end{tikzpicture}
\end{center}
\legende{Structure d'un raisonnement logique (syllogisme) : deux prémisses vraies conduisent nécessairement à la conclusion.}
\end{popfigure}

\begin{methode}[Pour bien argumenter]
\begin{enumerate}
\item Énoncer clairement sa \textbf{thèse} (l'idée qu'on défend).
\item Donner au moins \textbf{deux arguments} (raisons logiques, pas des impressions).
\item Illustrer chaque argument par un \textbf{exemple concret}.
\item Répondre à une \textbf{objection} (ce qu'un adversaire pourrait répliquer).
\item \textbf{Conclure} en rappelant ce qui a été établi.
\end{enumerate}
\end{methode}

\courssub{Le dialogue et la maïeutique socratique}

La philosophie ne se fait pas seul dans sa chambre : elle est d'abord un \cle{dialogue}. \textbf{Socrate} (470-399 av. J.-C.) passait ses journées sur la place publique d'Athènes à interroger ses concitoyens : « Qu'est-ce que la justice ? Qu'est-ce que le courage ? » Par des questions successives, il révélait les contradictions de leurs réponses et les conduisait à chercher mieux.

\begin{definition}[La maïeutique]
La \cle{maïeutique} (du grec \emph{maieutikê}, l'art d'accoucher) est la méthode de Socrate qui, par des questions successives, amène l'interlocuteur à découvrir lui-même la vérité, comme la sage-femme aide la mère à mettre l'enfant au monde. Socrate, fils d'une sage-femme, disait n'accoucher que des esprits.
\end{definition}

La leçon est profonde : le professeur de philosophie ne verse pas un savoir dans une tête vide ; il aide chaque élève à \emph{faire naître} en lui des idées qui s'y trouvent déjà en germe. Apprendre à philosopher, c'est apprendre à penser par soi-même.

\begin{savaistu}[Socrate, le taon et la sage-femme]
Socrate se comparait à un \textbf{taon}, cet insecte qui pique le cheval pour le réveiller : il piquait la cité d'Athènes, endormie dans ses habitudes, pour l'obliger à réfléchir. Il se comparait aussi à une \textbf{sage-femme} qui aide les esprits à accoucher des idées. Les Athéniens, agacés par ses questions, l'ont fait condamner à mort en 399 av. J.-C. : il est mort pour avoir refusé de renoncer à examiner la vie. « Une vie sans examen ne vaut pas la peine d'être vécue », aurait-il déclaré à ses juges.
\end{savaistu}

\courssub{L'esprit critique : distinguer opinion et vérité}

Le dialogue et l'argumentation visent un but : séparer le vrai du faux. C'est l'\cle{esprit critique}.

\begin{definition}[Opinion et vérité]
L'\cle{opinion} (en grec \emph{doxa}) est une croyance non vérifiée, fondée sur l'habitude, l'émotion, la rumeur ou l'autorité d'autrui. La \cle{vérité démontrée} est une affirmation justifiée par des raisons et des preuves que tout être raisonnable peut contrôler. L'\cle{esprit critique} est l'attitude qui refuse les préjugés et les rumeurs, et qui n'accepte une idée qu'après l'avoir examinée.
\end{definition}

\begin{exemplebox}[L'esprit critique en acte à Douala]
Une rumeur se répand un matin : « Le lycée technique de Douala va fermer, les élèves seront transférés ! » Panique dans les classes. Mais l'élève philosophe réagit autrement : (1) il vérifie la \textbf{source} — qui l'a dit ? un message anonyme ou une décision officielle ? ; (2) il demande des \textbf{preuves} — existe-t-il un communiqué du MINESEC ? ; (3) il \textbf{confronte les versions} — que dit l'administration ? les enseignants ? ; (4) seulement alors il \textbf{conclut}, prudemment. Ce geste simple est l'esprit critique en acte : ne jamais croire avant d'avoir examiné.
\end{exemplebox}

\courssub{L'universalité et la rigueur}

Dernière exigence de la méthode : la philosophie ne cherche pas des vérités valables seulement pour moi, pour mon village ou pour mon époque. Elle vise des conclusions \textbf{universelles}, c'est-à-dire valables pour \emph{tout être raisonnable}, partout et toujours. Quand Socrate demande « Qu'est-ce que la justice ? », il ne veut pas la justice telle qu'on la pratique à Athènes ou dans tel quartier : il veut ce qui rend \emph{toute} action juste, à Athènes comme à Bafoussam. Cette exigence impose la \cle{rigueur} : des définitions précises, des raisonnements sans contradiction, des exemples contrôlables. La vérité philosophique ressemble en cela à la vérité mathématique : $2 + 2 = 4$ est vrai à Douala, à Paris et sur la Lune.

\courssub{Les étapes pratiques de la démarche philosophique}

Réunissons tout ce qui précède. Face à un sujet (« La technique nous rend-elle libres ? »), le philosophe suit un parcours ordonné :

\begin{popfigure}
\begin{center}
\begin{tikzpicture}[
  etape/.style={draw=popdark, fill=popgreenL, rounded corners, text width=3.6cm, align=center, font=\small, thick, minimum height=1cm},
  fleche/.style={-{Stealth[length=2.5mm]}, very thick, poporange}]
\node[etape, align=center] (e1) at (0,0){\textbf{1. Étonnement}\\ s'étonner, s'interroger};
\node[etape, align=center] (e2) at (4.6,0){\textbf{2. Doute}\\ suspendre ses certitudes};
\node[etape, align=center] (e3) at (9.2,0){\textbf{3. Analyse}\\ définir, décomposer};
\node[etape, align=center] (e4) at (9.2,-2.4){\textbf{4. Argumentation}\\ avancer des raisons};
\node[etape, align=center] (e5) at (4.6,-2.4){\textbf{5. Discussion}\\ examiner les objections};
\node[etape, align=center] (e6) at (0,-2.4){\textbf{6. Conclusion}\\ thèse justifiée};
\draw[fleche] (e1) -- (e2);
\draw[fleche] (e2) -- (e3);
\draw[fleche] (e3) -- (e4);
\draw[fleche] (e4) -- (e5);
\draw[fleche] (e5) -- (e6);
\end{tikzpicture}
\end{center}
\legende{Organigramme des étapes de la démarche philosophique : de l'étonnement initial à la conclusion justifiée.}
\end{popfigure}

\begin{aretenir}[La démarche philosophique en six étapes]
\begin{enumerate}
\item \textbf{Problématiser} : partir de l'étonnement et formuler la question.
\item \textbf{Douter} : suspendre les opinions reçues.
\item \textbf{Analyser} : définir les termes, décomposer, repérer les présupposés.
\item \textbf{Argumenter} : construire un raisonnement (prémisses $\rightarrow$ conclusion).
\item \textbf{Discuter} : confronter sa thèse aux objections (dialogue).
\item \textbf{Conclure} : formuler une réponse rigoureuse et universelle.
\end{enumerate}
\end{aretenir}

\courssub{Application : rédiger un paragraphe argumenté}

Au Probatoire et au Baccalauréat technique, on te demandera de rédiger un \cle{paragraphe argumenté} : un texte court qui suit le schéma \textbf{thèse — argument — exemple — mini-conclusion}.

\begin{exoresolu}[Rédiger un paragraphe argumenté]
\textbf{Sujet :} Faut-il vérifier une information avant de la partager sur les réseaux sociaux ? Rédige un paragraphe argumenté.
\tcblower
\textbf{Solution.} \emph{Thèse :} Il faut toujours vérifier une information avant de la partager, car la diffusion d'une fausse nouvelle peut causer des torts réels. \emph{Argument 1 :} Une rumeur non vérifiée peut détruire une réputation ou provoquer une panique : celui qui partage sans vérifier devient responsable du mal causé. \emph{Exemple :} Lorsqu'une fausse annonce de fermeture circula dans un lycée technique de Douala, des parents retirèrent leurs enfants des classes avant que l'administration ne démente officiellement. \emph{Argument 2 :} Vérifier une information est facile : il suffit de contrôler la source, de chercher une confirmation officielle et de confronter plusieurs versions, comme le fait l'esprit critique. \emph{Objection et réponse :} On objectera que vérifier prend du temps ; mais quelques minutes de vérification valent mieux que des jours de conséquences d'une erreur. \emph{Mini-conclusion :} Partager sans vérifier, c'est renoncer à penser ; vérifier avant de partager, c'est déjà philosopher.
\end{exoresolu}

\begin{exercice}[À toi de philosopher]
Rédige un paragraphe argumenté (thèse, deux arguments, un exemple camerounais, une réponse à une objection, une mini-conclusion) sur le sujet suivant : « Les notes scolaires reflètent-elles toujours l'intelligence d'un élève ? »
\end{exercice}

\begin{corrige}[Exercice 1]
Éléments de correction. \emph{Thèse possible :} Les notes ne reflètent qu'imparfaitement l'intelligence. \emph{Arguments :} (1) une note mesure une performance à un moment donné, influencée par la santé, le stress ou les conditions matérielles (un élève qui travaille la nuit pour aider sa famille dort en classe) ; (2) l'intelligence a plusieurs formes (technique, pratique, sociale) qu'un devoir écrit ne mesure pas toujours — un excellent ajusteur peut être moyen en rédaction. \emph{Objection :} les notes restent un instrument nécessaire pour classer et orienter. \emph{Réponse :} c'est vrai, mais nécessaire ne veut pas dire parfait : il faut les interpréter avec prudence. \emph{Mini-conclusion :} la note est un outil utile, non un verdict sur l'intelligence. Toute copie cohérente, structurée et illustrée d'un exemple concret est valorisée.
\end{corrige}

\begin{aretenir}[L'essentiel de la section]
La méthode philosophique repose sur cinq piliers : le \cle{doute méthodique} (provisoire, non sceptique), l'\cle{analyse} (définir, décomposer, repérer les présupposés), l'\cle{argumentation} (prémisses et conclusion), le \cle{dialogue} et la \cle{maïeutique} socratique, et l'\cle{esprit critique} qui distingue l'opinion de la vérité démontrée. Elle vise des conclusions universelles, valables pour tout être raisonnable. Sa mise en œuvre pratique suit six étapes : problématiser, douter, analyser, argumenter, discuter, conclure.
\end{aretenir}

\coursec{L'importance de la philosophie}

Dans la section précédente, nous avons découvert la \cle{méthode} de la philosophie : problématiser, analyser les concepts, argumenter, discuter les thèses. Mais une question demeure, que tout élève de Première technique peut légitimement se poser : à quoi sert tout cela ? Un tourneur-fraiseur fabrique des pièces, un électricien installe des circuits, un comptable tient des livres : que \emph{produit} le philosophe ? Cette section répond à cette interrogation en montrant que la philosophie n'est pas un luxe inutile, mais une formation indispensable de l'homme, du citoyen, du technicien et du moraliste. Nous examinerons successivement son importance individuelle, sociale et politique, morale, professionnelle et culturelle, avant d'analyser ses rapports avec la science, de voir en quoi elle est une véritable école de la liberté, puis de discuter les critiques dont elle fait l'objet.

\courssub{L'importance de la philosophie pour l'individu}

\textbf{Intuition.} Imaginez un élève qui répète tout ce qu'on lui dit : « les professeurs disent que... », « les réseaux sociaux affirment que... », « mon quartier pense que... ». Un tel élève ne pense pas : il \emph{rediffuse}. La philosophie commence précisément quand on cesse de rediffuser pour se demander : « Est-ce vrai ? Est-ce juste ? Qu'en pense \emph{moi}, après examen ? »

\begin{definition}[L'autonomie de pensée]
L'\cle{autonomie de pensée} est la capacité à juger par soi-même, sans subir passivement les opinions des autres. Elle ne consiste pas à rejeter systématiquement ce que disent les autres, mais à examiner leurs affirmations à la lumière de la raison avant de les adopter ou de les refuser.
\end{definition}

La première utilité de la philosophie est donc de \textbf{former le jugement}. En nous habituant à définir les mots, à distinguer le vrai du faux, le probable du certain, le préjugé de la connaissance, elle nous donne des instruments pour ne plus être dupes. Platon, dans l'allégorie de la caverne (\emph{République}, Livre VII), décrivait déjà des hommes enchaînés qui prennent des ombres pour la réalité : philosopher, c'est se détourner des ombres pour voir les choses elles-mêmes.

La deuxième utilité est de \textbf{développer l'autonomie}. L'élève qui a appris à problématiser ne demande plus seulement « que faut-il réciter ? » mais « pourquoi est-ce vrai ? ». Il devient sujet de sa pensée au lieu d'en être l'objet.

La troisième utilité, plus profonde, est d'\textbf{apprendre à vivre et à mourir}. Montaigne, philosophe français du XVI\textsuperscript{e} siècle, écrit dans les \emph{Essais} (I, 20) : « philosopher, c'est apprendre à mourir ». Cette formule ne signifie pas que la philosophie rend triste ; elle signifie que réfléchir sur la mort, sur le temps qui passe, sur ce qui compte vraiment, nous délivre de la peur et nous apprend à bien vivre. Celui qui a réfléchi au sens de son existence ne la gaspille pas : il choisit ses priorités, supporte les épreuves et donne une direction à sa vie.

\begin{aretenir}[À retenir]
Pour l'individu, la philosophie a trois fonctions majeures : former le \cle{jugement} (ne plus être dupe), développer l'\cle{autonomie de pensée} (penser par soi-même) et apprendre à \cle{vivre} en réfléchissant sur le sens de l'existence (Montaigne : « philosopher, c'est apprendre à mourir »).
\end{aretenir}

\courssub{L'importance sociale et politique de la philosophie}

\textbf{Intuition.} Une société où personne ne discute, où chacun impose son opinion par la force ou le mensonge, est une société violente ou manipulée. La démocratie suppose au contraire des citoyens capables de débattre, d'écouter, de nuancer.

La philosophie est d'abord une \textbf{éducation à la citoyenneté}. En nous apprenant à justifier nos opinions par des arguments et non par des cris, elle forme le citoyen capable de participer au débat public : voter en connaissance de cause, discuter des lois, juger l'action des gouvernants. Elle est aussi une école de la \cle{tolérance} : celui qui a l'habitude d'examiner plusieurs thèses contradictoires comprend que l'adversaire n'est pas un ennemi, mais un partenaire de recherche.

Enfin, la philosophie est une \textbf{arme contre la manipulation}. À l'ère des réseaux sociaux, les rumeurs et les \emph{fake news} (fausses informations) circulent plus vite que les vérifications. L'esprit critique, formé par la philosophie, pose toujours les mêmes questions : qui parle ? quelle est la source ? quelles preuves ? à qui profite cette information ?

\begin{exemplebox}[La philosophie face aux fausses informations]
Un jeune Camerounais de Yaoundé reçoit sur WhatsApp un message lui demandant de diffuser massivement une information alarmante sur une prétendue pénurie. Avant de partager, il se demande : d'où vient ce message ? est-il confirmé par une source fiable ? quelles seraient les conséquences de sa diffusion (panique, violences, achats précipités) ? Il décide de ne pas le relayer. En réfléchissant aux conséquences de son geste avant d'agir, il applique concrètement la philosophie : elle l'a rendu \cle{responsable}.
\end{exemplebox}

\courssub{L'importance morale de la philosophie}

\textbf{Intuition.} Tout le monde parle du « bien » et du « mal », du « juste » et de l'« injuste ». Mais qui a défini ces mots ? Sur quoi repose le devoir ? Sans réflexion, la morale se réduit à l'habitude ou à la crainte de la punition.

La philosophie nous oblige à \textbf{fonder nos valeurs} au lieu de les subir. Réfléchir sur le \cle{bien}, le \cle{juste} et le \cle{devoir}, c'est se demander : pourquoi dois-je être honnête même quand personne ne me voit ? qu'est-ce qui rend une action juste ? Socrate affirmait déjà, dans l'\emph{Apologie de Socrate} (Platon), qu'« une vie sans examen ne vaut pas la peine d'être vécue » : la morale non réfléchie n'est qu'une soumission aux coutumes.

Cette réflexion fonde la \cle{responsabilité} : je ne suis pleinement responsable de mes actes que si je les ai choisis en connaissance de cause, et non sous l'effet de la peur, de l'imitation ou de l'habitude. La philosophie transforme ainsi l'obéissance aveugle en engagement conscient.

\courssub{L'importance de la philosophie pour le technicien}

\textbf{Intuition.} On pourrait croire que le technicien n'a besoin que de gestes précis et de formules. Erreur : le technicien raisonne, diagnostique, décide — et il est responsable de la sécurité des personnes.

La philosophie apporte au technicien trois biens précieux :

\begin{enumerate}
\item \textbf{La rigueur logique.} Chercher une panne, c'est raisonner : formuler des hypothèses, les tester, éliminer les causes impossibles, conclure. Cette démarche hypothético-déductive est une gymnastique de l'esprit que la philosophie entraîne.
\item \textbf{La résolution de problèmes.} Problématiser — décomposer une difficulté en questions claires — est la première étape de toute solution technique, comme de toute dissertation.
\item \textbf{L'éthique professionnelle.} Le technicien doit répondre de questions morales : puis-je livrer un travail non conforme ? puis-je cacher un défaut ? La déontologie (honnêteté dans le travail, respect des normes de sécurité, sens de la responsabilité) n'est pas un supplément d'âme : elle protège des vies.
\end{enumerate}

\begin{exemplebox}[L'éthique du technicien en génie civil]
Un technicien en génie civil à Douala constate que le ferraillage d'un immeuble en construction ne respecte pas les normes. Le chef de chantier, pressé par les délais, lui demande de signer le procès-verbal de conformité. Malgré la pression et la menace de licenciement, le technicien refuse de signer : un immeuble mal construit peut s'effondrer et tuer. En refusant, il applique une \cle{éthique professionnelle} nourrie de réflexion philosophique : il a compris que sa signature engage sa responsabilité morale et non seulement son emploi.
\end{exemplebox}

\courssub{L'importance culturelle de la philosophie}

La philosophie est aussi une \textbf{ouverture sur les civilisations}. En étudiant les pensées grecque, africaine, asiatique ou moderne, on découvre que chaque culture a posé les grandes questions de l'existence à sa manière. Comprendre les conceptions d'autrui — par exemple la sagesse africaine de la parole, incarnée chez nous par les palabres villageoises où l'on recherche la vérité et la réconciliation par la discussion — permet le \cle{dialogue entre les cultures}. La philosophie combat ainsi l'ethnocentrisme (le fait de croire que sa propre culture est la seule valable) et prépare à la coexistence pacifique dans un pays comme le Cameroun, riche de plus de deux cents groupes ethniques.

\begin{popfigure}
\begin{center}
\begin{tikzpicture}[scale=1.0]
% Diagramme en étoile : 5 dimensions de l'importance de la philosophie
\draw[popline] (0,0) circle (2.2);
\foreach \a in {90,162,234,306,18} {
  \draw[popline, dashed] (0,0) -- (\a:2.2);
}
\foreach \a in {90,162,234,306,18} {
  \foreach \r in {0.55,1.1,1.65} {
    \draw[popline!50] (\a:\r) -- ({(\a+72)}:\r);
  }
}
% Valeurs (sur 4 niveaux)
\draw[thick, popblue, fill=popblueL, fill opacity=0.6]
  (90:2.0) -- (162:1.8) -- (234:1.9) -- (306:2.1) -- (18:1.7) -- cycle;
\foreach \a/\r in {90/2.0, 162/1.8, 234/1.9, 306/2.1, 18/1.7} {
  \fill[popblue] (\a:\r) circle (2pt);
}
% Étiquettes
\node[above, popdark, font=\small\bfseries] at (90:2.45) {Individuelle};
\node[above left, popdark, font=\small\bfseries, align=center] at (162:2.6) {Sociale et\\politique};
\node[below left, popdark, font=\small\bfseries, align=center] at (234:2.6) {Morale};
\node[below right, popdark, font=\small\bfseries, align=center] at (306:2.6) {Professionnelle\\(technicien)};
\node[above right, popdark, font=\small\bfseries, align=center] at (18:2.6) {Culturelle};
\node[popdark, font=\small] at (0,0) {\textbf{Philosophie}};
\end{tikzpicture}
\end{center}
\legende{Diagramme en étoile des cinq dimensions de l'importance de la philosophie : individuelle (jugement, autonomie), sociale et politique (citoyenneté, tolérance), morale (bien, juste, devoir), professionnelle (rigueur, éthique du technicien) et culturelle (dialogue entre civilisations). Chaque branche représente l'intensité de l'apport de la philosophie dans ce domaine.}
\end{popfigure}

\courssub{Philosophie et science : distinction et complémentarité}

\textbf{Intuition.} La science produit des savoirs certains, mesurables, techniques (médicaments, ponts, smartphones). La philosophie ne produit rien de « concret ». Sont-elles rivales ? La science a-t-elle rendu la philosophie obsolète ?

\begin{definition}[Distinction des objets]
Les \cle{sciences} ont pour objet les \emph{faits} : ce qui est observable, mesurable, reproductible, régi par des lois causales (ex. : la loi de la gravité, la structure de l'ADN). La \cle{philosophie} a pour objet les \emph{valeurs}, le \emph{sens} et les \emph{principes ultimes} : ce qui ne se mesure pas mais s'interprète (ex. : la justice, la beauté, le sens de la vie, la nature de la conscience).
\end{definition}

\begin{definition}[Distinction des méthodes]
La méthode scientifique est \emph{empirique et hypothético-déductive} : observation, hypothèse, expérience, validation/falsification. Elle vise l'explication causale (\emph{comment} ?). La méthode philosophique est \emph{rationnelle et critique} : analyse conceptuelle, argumentation logique, problématisation. Elle vise la compréhension de sens (\emph{pourquoi ?}, \emph{pour quoi ?}).
\end{definition}

\textbf{Complémentarité, non opposition.}
\begin{itemize}
\item \textbf{La science a besoin de la philosophie :} pour définir ses concepts (qu'est-ce que le temps ? l'espace ? la vie ?), pour fonder son éthique (bioéthique, éthique de l'IA, responsabilité du savant), pour réfléchir à ses limites (le scientisme est l'erreur qui croit que la science peut tout expliquer).
\item \textbf{La philosophie a besoin de la science :} pour ne pas spéculer dans le vide. La physique moderne (relativité, quantique) renouvelle la métaphysique du temps et de la matière ; les neurosciences éclairent la philosophie de l'esprit ; l'évolutionnisme transforme l'anthropologie philosophique.
\end{itemize}

\begin{aretenir}[À retenir]
Science et philosophie sont distinctes par l'objet (faits vs sens/valeurs) et la méthode (expérience vs raison critique), mais \emph{complémentaires}. La science dit \emph{comment} faire ; la philosophie demande \emph{pourquoi} et \emph{pour qui} faire. Un technicien qui ignore la philosophie risque l'aveuglement technique ; un philosophe qui ignore la science risque le vide spéculatif.
\end{aretenir}

\courssub{La philosophie, école de la liberté}

Toutes les utilités précédentes se résument en une seule : la philosophie \textbf{libère}. Le philosophe allemand Emmanuel Kant, dans son texte \emph{Qu'est-ce que les Lumières ?} (1784), définit les Lumières comme « la sortie de l'homme hors de l'état de minorité dont il est lui-même responsable ». Être mineur, c'est penser sous la direction d'autrui : le prêtre pense à ma place, le chef décide à ma place, la rumeur juge à ma place. La devise des Lumières, selon Kant, est : \emph{Sapere aude}, « aie le courage de te servir de ton propre entendement ».

Philosopher, c'est donc conquérir sa majorité intellectuelle : oser examiner soi-même, au lieu de laisser penser les autres à sa place. C'est pourquoi la philosophie est inséparable de la liberté : un peuple qui ne pense pas par lui-même se laisse gouverner par les préjugés, les rumeurs et les démagogues.

\begin{aretenir}[À retenir]
La philosophie est une \cle{école de la liberté} : selon Kant, elle fait sortir l'homme de la « minorité » en lui apprenant à se servir de son propre entendement. Penser par soi-même est la condition de toutes les autres libertés.
\end{aretenir}

\courssub{Limites et critiques de la philosophie (complément Bac)}

Pour être honnêtes — et la philosophie exige l'honnêteté —, il faut examiner les objections.

\textbf{Première critique : l'abstraction.} On reproche souvent à la philosophie de discuter de questions vagues (« qu'est-ce que l'être ? ») sans jamais conclure, pendant que la science avance et produit des résultats certains. Cette critique est partiellement fondée : la philosophie ne donne pas de réponses définitives et mesurables comme une expérience de physique.

\textbf{Réponse.} Mais toutes les questions importantes de la vie ne sont pas mesurables : « dois-je dénoncer une injustice ? », « quelle vie vaut la peine d'être vécue ? », « jusqu'où va ma responsabilité ? ». Aucune machine ne répond à ces questions ; seule la réflexion le peut. L'utilité de la philosophie est \emph{concrète} : elle produit des jugements éclairés, des décisions responsables, des citoyens libres.

\textbf{Deuxième critique : l'inutilité matérielle.} La philosophie ne fabrique ni pont, ni médicament, ni application. Réponse : elle fabrique ce sans quoi ponts, médicaments et applications peuvent devenir dangereux — des hommes capables de se demander \emph{pour quoi} et \emph{pour qui} ils construisent.

\begin{attention}[Erreur fréquente à éviter]
Ne jamais écrire dans une copie que « la philosophie ne sert à rien parce qu'elle ne produit pas d'objets matériels ». C'est confondre l'utilité \emph{matérielle} avec l'utilité tout court. La philosophie produit des \cle{jugements}, des \cle{valeurs} et des \cle{citoyens libres} : sa production est intellectuelle et morale, non industrielle. À l'examen (Probatoire), une telle affirmation non nuancée est sanctionnée comme un préjugé non examiné — précisément ce que la philosophie apprend à dépasser.
\end{attention}

\begin{exoresolu}[Exercice d'application : l'utilité de la philosophie]
\textbf{Énoncé.} Un élève de Première technique déclare : « Je veux devenir mécanicien ; la philosophie ne me servira jamais, je demande qu'on la supprime de mon emploi du temps. » En vous appuyant sur le cours, rédigez un paragraphe argumenté (8 à 12 lignes) montrant en quoi cette opinion est discutable.
\tcblower
\textbf{Solution détaillée.} L'opinion de cet élève repose sur une confusion entre utilité matérielle et utilité intellectuelle. Certes, la philosophie n'apprend pas à réparer un moteur ; mais elle forme des capacités sans lesquelles le mécanicien resterait un simple exécutant. D'abord, elle développe la rigueur logique : diagnostiquer une panne, c'est raisonner par hypothèses, exactement comme on construit une argumentation. Ensuite, elle fonde l'éthique professionnelle : le mécanicien qui répare mal un frein engage la vie de ses clients ; seule une réflexion sur le devoir et la responsabilité lui apprend à ne jamais sacrifier la sécurité au gain. Enfin, elle forme l'autonomie de pensée : le futur chef de garage devra juger par lui-même, résister aux pressions et prendre des décisions justes. Ainsi, loin d'être inutile au technicien, la philosophie est ce qui le distingue de la machine qu'il répare : elle fait de lui un professionnel responsable et un citoyen libre.
\end{exoresolu}

\begin{savaistu}[Le savais-tu ?]
Au Cameroun, la philosophie est enseignée en classe de Première et de Terminale dans \emph{toutes} les séries, y compris les séries techniques (F, G, E, etc.). Ce choix de l'État n'est pas un hasard : la philosophie y est considérée comme une \cle{formation du citoyen}. Un pays qui forme d'excellents techniciens incapables de juger par eux-mêmes formerait des exécutants ; en enseignant la philosophie à tous, le système éducatif camerounais veut former des techniciens qui sont aussi des hommes libres et responsables.
\end{savaistu}

\begin{exercice}[TP guidé : philosopher sur des situations camerounaises]
\begin{enumerate}
\item \textbf{Situation 1 (Éthique professionnelle).} Un électricien vient de terminer le câblage d'un bâtiment. Il remarque un défaut d'isolement sur une gaine, mais le chef de chantier exige la réception immédiate pour toucher la prime de fin de chantier. Que doit faire l'électricien ? Justifiez en mobilisant les notions de \cle{responsabilité}, \cle{devoir} et \cle{éthique professionnelle}.
\item \textbf{Situation 2 (Citoyenneté et esprit critique).} Une rumeur circule sur Facebook affirmant qu'un produit alimentaire local est empoisonné. La rumeur est anonyme, sans source médicale, mais elle est partagée des milliers de fois. En tant que citoyen ayant suivi des cours de philosophie, quelle démarche adoptez-vous avant de partager (ou non) l'information ? Mobilisez \cle{esprit critique}, \cle{source}, \cle{conséquences}.
\item \textbf{Sujet de dissertation (type Probatoire).} « La philosophie est-elle un luxe réservé aux littéraires ? » Rédigez une introduction complète (accroche, définitions, problématique, annonce de plan) et élaborez un plan détaillé en deux ou trois parties avec arguments et exemples pour chaque partie.
\end{enumerate}
\end{exercice}

\begin{corrige}[Exercices du TP guidé]
\begin{enumerate}
\item \textbf{Situation 1.} L'électricien doit refuser la réception. Argument : la responsabilité professionnelle prime sur la pression hiérarchique ou l'intérêt financier immédiat. Le défaut d'isolement crée un risque d'électrocution ou d'incendie pour les futurs occupants. Le devoir de sécurité (déontologie) s'impose comme un impératif catégorique (Kant) : on ne peut pas traiter des vies humaines comme des moyens. Signer serait une faute morale et pénale (mise en danger de la vie d'autrui).
\item \textbf{Situation 2.} Ne pas partager immédiatement. Démarche : 1) Vérifier la source (qui émet l'info ? Ministère de la Santé ? Laboratoire certifié ?). 2) Croiser les informations (chercher un communiqué officiel). 3) Évaluer les conséquences du partage (panique, ruine économique du producteur, perte de confiance). 4) Appliquer le principe de précaution intellectuelle : tant que la vérité n'est pas établie, s'abstenir de relayer. C'est l'exercice concret de l'autonomie de pensée et de la responsabilité citoyenne.
\item \textbf{Sujet de dissertation.}
\textbf{Introduction.} Accroche : Souvent perçue comme une discipline abstraite, « de lettrés », la philosophie est parfois jugée inutile par les élèves des filières techniques qui visent l'efficacité concrète. Définitions : Philosophie = recherche rationnelle de la sagesse par l'analyse des concepts ; Luxe = superflu, non nécessaire. Problématique : L'utilité se réduit-elle à la production immédiate d'objets techniques ? La formation du jugement, de l'éthique et de la citoyenneté n'est-elle pas la condition \emph{sine qua non} d'une technique maîtrisée et humaine ? Annonce : Nous montrerons d'abord que la philosophie forme des compétences transversales indispensables au technicien (I), puis qu'elle est le fondement de l'éthique professionnelle et de la citoyenneté (II), avant de conclure qu'elle n'est pas un luxe mais la condition de la dignité du métier (III).
\textbf{Plan détaillé.}
I. La philosophie, formation de l'esprit technique (rigueur, problématisation, autonomie).
II. La philosophie, fondement de l'éthique et de la citoyenneté (responsabilité, tolérance, lutte contre la manipulation).
III. Dépasser l'opposition : le technicien-philosophe, figure de l'homme libre et responsable.
\end{enumerate}
\end{corrige}

\coursec{Philosophie et science}

Nous venons de voir que la philosophie est importante : elle forme l'esprit critique, elle éclaire l'action, elle rend libre. Mais une objection surgit aussitôt : la science ne fait-elle pas tout cela, et mieux encore ? Le téléphone portable qui fonctionne, le vaccin qui protège, le satellite qui annonce la pluie : tout cela, c'est la science. À quoi sert alors la philosophie ? Cette question, que beaucoup d'élèves se posent légitimement, suppose que science et philosophie seraient rivales. Toute cette section vise à montrer qu'il n'en est rien : science et philosophie ont des racines communes, des exigences communes, mais des objets et des méthodes différents, et elles ont besoin l'une de l'autre.

\courssub{Une même famille : les points communs entre philosophie et science}

Commençons par une observation simple. Quand un élève de Première technique affirme : « ce moteur chauffe parce que le courant est trop fort », il ne se contente pas de le dire : il doit le \emph{prouver}, par un raisonnement et par une mesure. Quand Socrate affirme : « une vie sans examen ne vaut pas la peine d'être vécue », il ne se contente pas de le dire non plus : il doit le \emph{justifier} par des arguments. Dans les deux cas, on refuse de croire sur parole, on exige des raisons. C'est le premier et le plus profond point commun entre la philosophie et la science.

\begin{propriete}[Ce qui unit philosophie et science (admis)]
La philosophie et la science partagent la même \cle{exigence de rationalité} : l'une et l'autre recherchent la \cle{vérité}, refusent l'\cle{opinion} non fondée (la \emph{doxa}) et la \cle{superstition}, et n'acceptent une affirmation que si elle est justifiée par des raisons. Mais elles diffèrent par leur \cle{objet} (ce qu'elles étudient) et par leur \cle{méthode} (la manière dont elles l'étudient).
\end{propriete}

Précisons ces trois points communs.

\textbf{1. La recherche de la vérité.} Ni le savant ni le philosophe ne cherchent à avoir raison pour le plaisir de gagner : ils cherchent ce qui est \emph{vrai}, c'est-à-dire ce qui est conforme à la réalité ou à la rigueur du raisonnement, même si cette vérité dérange. Galilée soutient que la Terre tourne contre l'opinion de son époque ; Socrate accepte de mourir plutôt que de renoncer à examiner sa vie.

\textbf{2. L'exigence de rationalité.} Dans les deux démarches, toute affirmation doit être accompagnée de ses \emph{raisons}. Le mot grec \emph{logos} signifie à la fois « parole », « raison » et « discours ordonné » : science et philosophie sont toutes deux filles du \emph{logos}, et ennemies du discours magique ou arbitraire.

\textbf{3. Le refus de l'opinion et de la superstition.} Dire « on dit que... », « tout le monde pense que... », « les ancêtres ont toujours cru que... » ne suffit ni au laboratoire ni dans la dissertation. L'opinion (\emph{doxa}) est un savoir non vérifié ; la superstition est une croyance qui attribue à des forces mystérieuses ce qui a des causes réelles. Accuser un voisin de sorcellerie pour expliquer une épidémie, c'est renoncer à chercher la cause véritable (un microbe, l'eau contaminée) : science et philosophie combattent ensemble ce renoncement.

\begin{popfigure}
\begin{tikzpicture}[scale=0.95]
  \draw[thick,popblue,fill=popblue!18] (-1.3,0) circle (2.4);
  \draw[thick,poporange,fill=poporange!18] (1.3,0) circle (2.4);
  \node[popblue,font=\bfseries] at (-2.5,2.1) {PHILOSOPHIE};
  \node[poporange,font=\bfseries] at (2.5,2.1) {SCIENCE};
  \node[align=center,font=\small] at (0,0.55) {\textbf{Points communs}};
  \node[align=center,font=\small] at (0,-0.15) {rationalité};
  \node[align=center,font=\small] at (0,-0.75) {recherche de la vérité};
  \node[align=center,font=\small] at (0,-1.35) {refus de l'opinion};
  \node[align=center,font=\small,popblue!80!black] at (-2.6,0.35) {réflexion sur le tout};
  \node[align=center,font=\small,popblue!80!black] at (-2.6,-0.35) {question du sens};
  \node[align=center,font=\small,popblue!80!black] at (-2.6,-1.05) {argumentation};
  \node[align=center,font=\small,poporange!80!black] at (2.6,0.35) {domaines particuliers};
  \node[align=center,font=\small,poporange!80!black] at (2.6,-0.35) {observation, expérience};
  \node[align=center,font=\small,poporange!80!black] at (2.6,-1.05) {mesure, calcul};
\end{tikzpicture}
\legende{Diagramme de Venn : la zone commune (rationalité, vérité, refus de l'opinion) et les zones propres à chaque discipline.}
\end{popfigure}

\courssub{Des objets différents : le particulier et le tout}

Si science et philosophie partagent la même exigence, pourquoi ne sont-elles pas une seule discipline ? Parce qu'elles ne regardent pas la même chose.

La \cle{science} est divisée en disciplines \emph{particulières} : la biologie étudie le vivant, la physique la matière et l'énergie, les mathématiques les nombres et les figures, la médecine le corps malade. Chaque science découpe dans la réalité un \emph{domaine délimité} et le traite avec des outils précis. Le mot « science » vient du latin \emph{scientia}, « connaissance » : mais c'est toujours la connaissance \emph{d'un objet déterminé}.

La \cle{philosophie}, elle, interroge \emph{le tout} et \emph{le sens}. Elle ne demande pas seulement : « comment fonctionne le vivant ? » mais : « qu'est-ce que vivre ? » Elle ne demande pas seulement : « comment calculer ? » mais : « qu'est-ce que la vérité ? » Elle ne demande pas seulement : « comment soigner ? » mais : « qu'est-ce qu'une vie digne d'être vécue ? » Aristote appelait la philosophie la science des « premières causes et des premiers principes » : elle remonte aux questions que les sciences particulières présupposent sans les examiner.

\begin{definition}[Deux types de vérité]
Une \cle{vérité scientifique} est une proposition \emph{vérifiable par l'expérience et la mesure} : on peut la tester, la reproduire, la chiffrer (exemple : l'eau bout à 100 degrés au niveau de la mer). Une \cle{vérité philosophique} s'établit par la \emph{cohérence du raisonnement} : elle ne se mesure pas, mais elle doit être démontrée sans contradiction (exemple : « tout homme est libre de penser »).
\end{definition}

\begin{exemplebox}[Le satellite et la question du sens]
Un satellite météorologique prévoit les pluies sur le plateau de l'Adamaoua et permet aux éleveurs de préparer la transhumance : c'est un résultat de \textbf{science} (observation, mesure, calcul). Mais se demander si l'homme \emph{doit} maîtriser toute la nature, s'il a le droit de tout transformer, et ce qu'il perd ou gagne à le faire : c'est une question \textbf{philosophique}. La science dit \emph{ce qui est} et \emph{ce qui est possible} ; la philosophie demande \emph{ce qui est bon} et \emph{ce que cela signifie}.
\end{exemplebox}

\courssub{Des méthodes différentes : l'expérience et la réflexion}

La différence d'objet entraîne une différence de \cle{méthode}.

La science procède par \emph{observation}, \emph{expérimentation} et \emph{mesure}. Le savant formule une hypothèse, la soumet à l'épreuve de l'expérience, mesure les résultats, et ne retient que ce qui résiste à la vérification. Claude Bernard, grand physiologiste du XIX\textsuperscript{e} siècle, résumait cette démarche : observer, émettre une hypothèse, expérimenter, conclure. L'instrument (thermomètre, microscope, oscilloscope) est l'outil du savant.

La philosophie procède par \emph{réflexion}, \emph{analyse conceptuelle} et \emph{argumentation}. Le philosophe ne mesure pas : il définit ses concepts (qu'appelle-t-on justice ? liberté ? vérité ?), examine les raisons pour et contre, confronte les thèses, et construit un discours cohérent. Son « laboratoire », c'est le dialogue et l'écriture ; son instrument, c'est le raisonnement.

On peut résumer cette différence par une formule célèbre : \textbf{la science répond au « comment », la philosophie interroge le « pourquoi » et le « à quoi bon »}. La science explique \emph{comment} le virus se transmet ; la philosophie demande \emph{pourquoi} la santé est un bien et \emph{à quoi bon} vivre. La science construit la fusée ; la philosophie demande si conquérir l'espace a un sens pour l'humanité.

\begin{attention}[Piège fréquent : opposer radicalement science et philosophie]
L'erreur la plus répandue consiste à croire que la science a « remplacé » la philosophie et l'a rendue inutile, ou inversement à mépriser la science au nom de la réflexion. C'est une double erreur. D'une part, la science elle-même repose sur des \emph{présupposés philosophiques} qu'elle ne démontre pas : que le monde est intelligible, que la vérité existe, que la raison peut l'atteindre. D'autre part, une philosophie qui ignorerait les découvertes scientifiques parlerait dans le vide. Au Bac, évitez donc la conclusion facile « la science a tué la philosophie » : montrez leur complémentarité.
\end{attention}

\courssub{La filiation historique : les sciences sont nées de la philosophie}

Un fait historique éclaire tout le débat : \emph{les sciences sont des enfants de la philosophie}. Dans la Grèce antique, il n'existait qu'un seul savoir : la \emph{philosophia}, l'amour de la sagesse, qui englobait tout. Aristote écrivait aussi bien sur la logique, sur les animaux, sur le ciel, sur l'âme et sur la politique. Puis, peu à peu, chaque domaine, en trouvant sa méthode propre (l'expérience, la mesure), a quitté la maison maternelle : la physique avec Galilée et Newton, la biologie au XIX\textsuperscript{e} siècle, la psychologie à la fin du XIX\textsuperscript{e} siècle, la logique mathématique au XX\textsuperscript{e} siècle.

\begin{popfigure}
\begin{tikzpicture}[
  grow=right,
  level 1/.style={sibling distance=3.4cm, level distance=3.2cm},
  level 2/.style={sibling distance=1.5cm, level distance=3.2cm},
  every node/.style={font=\small},
  edge from parent/.style={draw=popmuted,thick}
]
\node[fill=popgold!40,rounded corners,draw=popgold,thick,align=center,font=\small\bfseries] {Philosophie\\ antique}
  child { node[fill=popgreen!25,rounded corners,draw=popgreen,thick] {Éthique} }
  child { node[fill=poppurple!20,rounded corners,draw=poppurple,thick] {Logique} }
  child { node[fill=poppink!25,rounded corners,draw=poppink,thick] {Psychologie} }
  child { node[fill=poporange!25,rounded corners,draw=poporange,thick] {Biologie} }
  child { node[fill=popblue!20,rounded corners,draw=popblue,thick] {Physique} };
\node[align=left,font=\footnotesize\itshape,popmuted] at (7.6,4.0) {restées dans la philosophie :};
\node[align=left,font=\footnotesize\itshape,popmuted] at (7.6,3.55) {éthique, logique philosophique};
\node[align=left,font=\footnotesize\itshape,popmuted] at (7.6,-3.1) {devenues sciences autonomes :};
\node[align=left,font=\footnotesize\itshape,popmuted] at (7.6,-3.55) {physique, biologie, psychologie};
\end{tikzpicture}
\legende{Arbre de la filiation : de la philosophie antique sont issues, au fil des siècles, les sciences particulières ; l'éthique et la réflexion sur le sens demeurent philosophiques.}
\end{popfigure}

\begin{savaistu}[La physique s'appelait « philosophie naturelle »]
Isaac Newton publia en 1687 son œuvre majeure sous le titre \emph{Principes mathématiques de la philosophie naturelle} (\emph{Principia}). Au XVII\textsuperscript{e} siècle, la physique s'appelait donc encore « philosophie naturelle », et les savants se nommaient « philosophes de la nature ». Le mot « scientifique » (\emph{scientist}) n'apparaît en anglais qu'au XIX\textsuperscript{e} siècle. La science moderne est littéralement une branche de la philosophie qui a grandi.
\end{savaistu}

Cette filiation a une conséquence importante : quand une science se sépare de la philosophie, celle-ci ne disparaît pas ; elle conserve les questions que la science, par méthode, ne peut pas traiter. La psychologie scientifique mesure les réactions et les comportements ; mais la question « qu'est-ce que la conscience ? » reste philosophique. La physique décrit l'univers ; mais la question « pourquoi y a-t-il quelque chose plutôt que rien ? » reste philosophique.

\courssub{La complémentarité : la philosophie au chevet des sciences}

La science avance vite, parfois plus vite que notre réflexion morale. Or \emph{pouvoir} faire quelque chose ne signifie pas qu'il faille le faire. C'est ici que la philosophie retrouve son rôle : elle réfléchit sur les \emph{fondements} des sciences (qu'est-ce qu'une preuve ? qu'est-ce qu'un fait ?) et sur leurs \emph{conséquences éthiques}.

\textbf{Premier cas : la médecine et la bioéthique.} Au Cameroun comme partout, la médecine progresse : dialyse, greffes d'organes, réanimation. La science répond à la question « comment ? » : comment greffer un rein, comment prolonger la vie. Mais d'autres questions surgissent, auxquelles aucune mesure ne peut répondre : faut-il prolonger la vie à tout prix ? Qui doit recevoir l'organe disponible quand les malades sont nombreux ? Peut-on interrompre les soins d'un patient sans espoir ? Ces questions relèvent de la \cle{bioéthique}, discipline où la philosophie dialogue avec la médecine. Le médecin apporte le savoir technique ; le philosophe (et avec lui le juriste, le religieux, le malade lui-même) apporte l'examen du sens et des valeurs.

\textbf{Deuxième cas : les réseaux sociaux et la vérité.} L'ingénierie informatique a produit WhatsApp, Facebook, TikTok : prouesses techniques indiscutables. Mais ces outils diffusent aussi les rumeurs et les fausses informations (\emph{fake news}), qui ont parfois semé la panique dans nos villes et nos villages. La technique répond au « comment » (comment transmettre un message à des millions de personnes en une seconde) ; la philosophie pose les questions : qu'est-ce qu'une information vraie ? Quelle responsabilité avons-nous quand nous partageons un message ? La liberté d'expression a-t-elle des limites ? Former l'esprit critique de l'utilisateur, c'est un travail philosophique.

\begin{exoresolu}[Science ou philosophie : à vous de classer]
\textbf{Énoncé.} Pour chacune des questions suivantes, indiquez si elle relève de la science, de la philosophie, ou des deux, en justifiant : (1) « Quelle est la cause de la malaria ? » ; (2) « Faut-il obliger la vaccination ? » ; (3) « Comment fabriquer une batterie plus durable ? » ; (4) « L'intelligence artificielle peut-elle penser ? »
\tcblower
\textbf{Solution détaillée.}
(1) \emph{Science} : la cause de la malaria (le paludisme) est un parasite transmis par le moustique ; cela se découvre par l'observation et l'expérience, et se vérifie au microscope.
(2) \emph{Philosophie (et droit)} : la science dit que le vaccin protège (le « comment »), mais décider si l'État peut l'imposer engage des valeurs : liberté individuelle, bien commun. C'est une question de sens et de morale.
(3) \emph{Science et technique} : c'est un problème de « comment », résolu par la recherche expérimentale et la mesure.
(4) \emph{Les deux} : l'informatique décrit ce que fait la machine (calcul, apprentissage) ; mais savoir si calculer, c'est \emph{penser}, exige de définir le concept de pensée : c'est l'analyse philosophique. La bonne réponse montre donc le dialogue entre les deux disciplines.
\end{exoresolu}

\courssub{Synthèse : ni rivalité ni confusion, mais dialogue}

Nous pouvons maintenant répondre à l'objection du début. La science n'a pas rendu la philosophie inutile, et la philosophie ne rivalise pas avec la science. Leur relation obéit à trois principes.

\textbf{1. Ni confusion.} La philosophie ne fait pas d'expériences et la science ne tranche pas les questions de sens. Confondre les deux, c'est exposer la philosophie au ridicule (prétendre démontrer par la seule réflexion ce que seule l'expérience peut établir) et la science à l'aveuglement (croire que mesurer suffit à comprendre l'homme).

\textbf{2. Ni rivalité.} Elles ne sont pas concurrentes, car elles ne répondent pas aux mêmes questions. La science donne des \emph{pouvoirs} ; la philosophie interroge les \emph{fins}. Plus la science est puissante, plus la question philosophique « à quoi bon ? » devient urgente.

\textbf{3. Mais dialogue.} La science a besoin de la philosophie pour examiner ses présupposés (la vérité, la raison, l'objet) et ses conséquences morales (bioéthique, intelligence artificielle, environnement). La philosophie a besoin de la science pour ne pas spéculer dans le vide : une réflexion sur la nature ou sur l'homme qui ignore la physique et la biologie serait une rêverie. Comme l'écrivait le philosophe des sciences Bachelard, la philosophie doit se former « au contact de la science » ; et le savant Einstein reconnaissait que la science sans épistémologie (réflexion sur la science) serait « primitive et confuse ».

\begin{aretenir}[L'essentiel de la section]
\begin{itemize}
\item Philosophie et science partagent la même exigence : la \cle{rationalité}, la recherche de la \cle{vérité}, le refus de l'opinion et de la superstition.
\item Elles diffèrent par l'\cle{objet} : la science étudie des domaines particuliers, la philosophie interroge le tout et le sens.
\item Elles diffèrent par la \cle{méthode} : observation, expérience et mesure d'un côté ; réflexion, analyse des concepts et argumentation de l'autre.
\item La science répond au « \cle{comment} », la philosophie au « \cle{pourquoi} » et au « \cle{à quoi bon} ».
\item Historiquement, les sciences sont \cle{nées de la philosophie} (physique, biologie, psychologie, logique).
\item Aujourd'hui, elles sont \cle{complémentaires} : la philosophie éclaire les fondements et les enjeux éthiques des sciences (bioéthique, réseaux sociaux, intelligence artificielle).
\end{itemize}
\end{aretenir}

\begin{exercice}[Pour s'entraîner à la dissertation]
Sujet : « La science a-t-elle rendu la philosophie inutile ? »
Travail demandé : (1) définissez les termes du sujet ; (2) construisez deux parties : la première montrant en quoi la science semble suffire (ses succès, sa méthode), la seconde montrant ce que la philosophie apporte d'irremplaçable (le sens, l'éthique, les fondements) ; (3) rédigez une introduction de cinq à huit lignes.
\end{exercice}

\begin{corrige}[Exercice : éléments de correction]
(1) \emph{Définitions} : la science est l'ensemble des connaissances vérifiables par l'expérience et la mesure ; la philosophie est la réflexion rationnelle sur le tout et sur le sens. « Rendre inutile » signifie : supprimer la raison d'être.
(2) \emph{Plan} : I. La science semble suffire : elle explique le monde (médecine, technique), elle vérifie ses résultats, elle a absorbé les anciennes questions philosophiques (filiation). II. Mais la philosophie demeure nécessaire : la science présuppose la valeur de la vérité sans la fonder ; elle donne des pouvoirs sans fixer les fins ; les progrès mêmes (greffe, intelligence artificielle) posent des questions éthiques que nulle expérience ne tranche. III. (conclusion) : ni rivalité ni confusion, mais dialogue.
(3) \emph{Introduction possible} : « Les succès de la science sont éclatants : elle soigne, elle prévoit, elle transforme le monde. Dès lors, la philosophie, qui ne mesure rien et ne fabrique rien, semble condamnée. Mais suffit-il de savoir comment agir pour savoir pourquoi agir ? Nous montrerons que si la science répond au comment, la philosophie demeure seule à interroger le sens : loin d'être rivales, elles sont complémentaires. »
\end{corrige}

\coursec{TP guidé : philosopher sur des situations camerounaises}

Dans la section précédente, nous avons établi que la philosophie et la science se distinguent par leur objet et leur méthode, tout en se complétant : la science décrit et explique les faits, la philosophie interroge le sens, les valeurs et les fondements. Mais une leçon de philosophie ne s'achève jamais sur des définitions : elle doit devenir une \cle{pratique}. C'est toute l'ambition de ces travaux pratiques guidés. Comme le technicien apprend en manipulant, le philosophe apprend en \emph{philosophant} : en lisant des textes, en comparant, en débattant, en rédigeant et en exerçant son esprit critique sur des situations concrètes de la vie camerounaise. Chaque TP mobilise un savoir de l'unité et un savoir-faire : analyser, comparer, argumenter, rédiger et justifier.

\courssub{TP 1 — Analyse de documents : trois textes pour définir la philosophie}

\textbf{Objectif.} À partir de trois textes courts, dégager par induction une définition personnelle et justifiée de la philosophie.

\textbf{Documents soumis à l'analyse.}

\begin{exemplebox}[Les trois textes du TP 1]
\textbf{Texte 1 — Platon, \emph{Théétète} :} « L'étonnement est le sentiment d'un philosophe, et la philosophie commence dans l'étonnement. »

\textbf{Texte 2 — Socrate (rapporté par Platon, \emph{Apologie de Socrate}) :} « Une vie sans examen ne vaut pas la peine d'être vécue. »

\textbf{Texte 3 — Kant, \emph{Réponse à la question : qu'est-ce que les Lumières ?} :} « Aie le courage de te servir de ton propre entendement ! »
\end{exemplebox}

Pour traiter ces textes, on applique rigoureusement la méthode d'analyse du texte philosophique.

\begin{methode}[Analyser un texte philosophique en cinq étapes]
\begin{enumerate}
\item \textbf{Identifier} l'auteur, l'œuvre et le thème du texte (de quoi parle-t-on ?).
\item \textbf{Dégager la thèse} : quelle est l'idée principale défendue par l'auteur ?
\item \textbf{Relever les arguments} : quelles raisons, quels mots-clés soutiennent cette thèse ?
\item \textbf{Reformuler} avec ses propres mots, en une ou deux phrases simples.
\item \textbf{Discuter} : la thèse est-elle convaincante ? Peut-on lui opposer un exemple ou une objection ?
\end{enumerate}
\end{methode}

\textbf{Application guidée.} Le texte 1 fait de l'\cle{étonnement} l'origine de la philosophie : devant le ciel étoilé, la naissance d'un enfant ou la foudre sur le mont Cameroun, l'homme cesse de prendre le monde pour évident et se demande « pourquoi ? ». Le texte 2 insiste sur l'\cle{examen de soi} : vivre sans jamais interroger ses choix, c'est subir sa vie au lieu de la conduire. Le texte 3 proclame l'\cle{autonomie de la pensée} : philosopher, c'est refuser de penser par procuration, à travers le chef, le groupe ou la rumeur.

\begin{exoresolu}[Synthèse attendue du TP 1]
Énoncé : à partir des trois textes, propose une définition de la philosophie en deux ou trois phrases.
\tcblower
Solution détaillée : les trois textes convergent. Platon désigne le \emph{point de départ} (l'étonnement devant le monde), Socrate désigne l'\emph{exercice} (l'examen critique de soi et des choses), Kant désigne la \emph{condition} (le courage de penser par soi-même). Définition possible : « La philosophie est une démarche rationnelle née de l'étonnement, qui consiste à examiner de manière critique notre existence, nos savoirs et nos valeurs, en faisant usage de notre propre raison. » Cette définition sera évaluée selon trois critères : fidélité aux textes, clarté de la formulation, caractère personnel de la reformulation.
\end{exoresolu}

\courssub{TP 2 — Construction d'un tableau comparatif philosophie/science}

\textbf{Objectif.} Construire un tableau comparatif à partir de trois exemples fournis : la médecine, l'agriculture, la météorologie. Le travail consiste à montrer, sur chaque cas, ce que la science apporte et ce que la philosophie interroge.

\begin{exemplebox}[Travail guidé sur les trois exemples]
\begin{itemize}
\item \textbf{Médecine :} la science établit que le paludisme se transmet par la piqûre de l'anophèle et se soigne par des antipaludéens. La philosophie demande : qu'est-ce que la santé ? A-t-on le droit de soigner un malade contre sa volonté ? La fin (guérir) justifie-t-elle tous les moyens ?
\item \textbf{Agriculture :} la science améliore les rendements du cacao ou du maïs par les semences sélectionnées et les engrais. La philosophie demande : le progrès technique est-il toujours un bien ? Que devient le lien de l'homme à la terre des ancêtres ?
\item \textbf{Météorologie :} la science prévoit les pluies grâce aux satellites et aux modèles. La philosophie demande : une prévision scientifique est-elle une certitude ? Quelle confiance accorder au savoir, et pourquoi ?
\end{itemize}
\end{exemplebox}

\begin{center}
\begin{tabularx}{\linewidth}{|l|X|X|}
\hline
\rowcolor{popblueL} \textbf{Critère} & \textbf{Science} & \textbf{Philosophie} \\
\hline
Objet & Les faits observables et mesurables & Le sens, les valeurs, les fondements \\
\hline
Question type & Comment ? Pourquoi (causes) ? & Pourquoi (finalité) ? Est-ce bien, vrai, juste ? \\
\hline
Méthode & Observation, expérience, calcul & Réflexion, argumentation, examen critique \\
\hline
Résultat & Lois et théories vérifiables & Thèses discutables et justifiées \\
\hline
Exemple & Le vaccin prévient la maladie & Faut-il rendre la vaccination obligatoire ? \\
\hline
\end{tabularx}
\end{center}

\begin{attention}[Ne pas opposer artificiellement philosophie et science]
Le tableau comparatif ne signifie pas que la philosophie et la science sont ennemies. La science a besoin de la philosophie pour réfléchir à ses fondements et à ses limites éthiques ; la philosophie a besoin de la science pour ne pas spéculer dans le vide. Erreur fréquente : écrire « la science dit le vrai, la philosophie dit n'importe quoi ». Une thèse philosophique n'est pas une opinion quelconque : elle est \emph{argumentée}.
\end{attention}

\courssub{TP 3 — Débat philosophique : « Le technicien a-t-il besoin de la philosophie ? »}

\textbf{Objectif.} Organiser une joute verbale réglée, où deux camps s'affrontent par arguments et non par invectives.

\textbf{Organisation matérielle.} La classe est divisée en deux camps (le « pour » et le « contre ») et trois rôles sont attribués : le \cle{modérateur} (dirige les échanges, donne la parole), le \cle{rapporteur} (note les arguments et rédige la synthèse), le \cle{chronométreur} (veille au temps de parole, deux minutes par intervention).

\begin{methode}[Règles du débat philosophique]
\begin{enumerate}
\item \textbf{Écouter} l'adversaire jusqu'au bout avant de répondre ; ne jamais couper la parole.
\item \textbf{Argumenter sans attaquer la personne} : on réfute une thèse, on ne ridiculise pas celui qui la défend.
\item \textbf{Citer des exemples} concrets pour appuyer chaque argument.
\item \textbf{Respecter la parole de l'autre} et le temps imparti ; le modérateur sanctionne les débordements.
\item \textbf{Conclure par une synthèse} : le rapporteur dégage les points d'accord et de désaccord.
\end{enumerate}
\end{methode}

\textbf{Arguments attendus.} Pour le camp du « pour » : le technicien prend des décisions qui engagent des vies (un pont mal calculé, une installation électrique défectueuse) ; la philosophie forme le sens de la responsabilité, l'honnêteté professionnelle et l'esprit critique face aux ordres injustes. Pour le camp du « contre » : le technicien a besoin de savoir-faire précis, pas de spéculations ; un court-circuit ne se répare pas avec des concepts. La synthèse montrera que la technique répond au « comment » mais que la philosophie éclaire le « pour quoi faire » et le « pour qui » : un technicien compétent mais sans éthique peut devenir dangereux.

\begin{savaistu}[La joute verbale, une tradition africaine]
La joute verbale est une tradition de débat argumenté pratiquée dans plusieurs universités africaines, notamment au Cameroun, où des équipes s'affrontent sur des sujets de société devant un jury. Elle est l'héritière à la fois des \emph{palabres} africaines — où la parole réglée tranchait les conflits du village sous l'arbre à palabres — et de l'\emph{agora} grecque, place publique où Socrate interrogeait ses concitoyens. Philosopher en débattant est donc une pratique à la fois très ancienne et très actuelle.
\end{savaistu}

\courssub{TP 4 — Rédaction guidée : un paragraphe argumenté}

\textbf{Objectif.} Rédiger un paragraphe de 10 à 15 lignes répondant à la question : « Faut-il toujours dire la vérité ? », selon la structure \cle{thèse — arguments — exemples — conclusion}.

\begin{popfigure}
\begin{center}
\begin{tikzpicture}[node distance=0.35cm, every node/.style={font=\small}]
\node[draw=popblue, fill=popblueL, rounded corners, text width=0.86\linewidth, align=left] (t) {\textbf{1. Thèse} : la réponse défendue. \emph{Ex. : Dire la vérité est un devoir, mais il connaît des exceptions légitimes.}};
\node[draw=poporange, fill=poporangeL, rounded corners, text width=0.86\linewidth, align=left, below=of t] (a1) {\textbf{2. Argument 1} : raison principale. \emph{Ex. : Sans vérité, la confiance disparaît et aucune vie sociale n'est possible.}};
\node[draw=poporange, fill=poporangeL, rounded corners, text width=0.86\linewidth, align=left, below=of a1] (a2) {\textbf{3. Argument 2} : raison complémentaire ou objection traitée. \emph{Ex. : Mais une vérité brutale peut détruire inutilement : le mensonge peut parfois protéger.}};
\node[draw=popgreen, fill=popgreenL, rounded corners, text width=0.86\linewidth, align=left, below=of a2] (ex) {\textbf{4. Exemples} : cas concrets. \emph{Ex. : le malade à qui l'on ménage la vérité ; le témoin qui doit dire les faits au tribunal.}};
\node[draw=poppurple, fill=poppurpleL, rounded corners, text width=0.86\linewidth, align=left, below=of ex] (c) {\textbf{5. Conclusion} : bilan nuancé. \emph{Ex. : La vérité est la règle ; seule une raison morale supérieure peut justifier l'exception.}};
\draw[-{Stealth}, thick, popmuted] (t) -- (a1);
\draw[-{Stealth}, thick, popmuted] (a1) -- (a2);
\draw[-{Stealth}, thick, popmuted] (a2) -- (ex);
\draw[-{Stealth}, thick, popmuted] (ex) -- (c);
\end{tikzpicture}
\end{center}
\legende{Grille de la méthode du paragraphe argumenté, remplie sur le sujet « Faut-il toujours dire la vérité ? » : chaque étape enchaîne logiquement avec la suivante.}
\end{popfigure}

\begin{attention}[L'erreur capitale en dissertation]
L'erreur la plus fréquente consiste à donner son opinion sans argument ni exemple : « Je pense qu'il faut toujours dire la vérité parce que c'est bien. » Ce n'est pas philosopher, c'est bavarder. \textbf{Toute affirmation philosophique doit être justifiée} par une raison et illustrée par un exemple. Sans justification, une thèse, même vraie, ne vaut rien dans une copie de philosophie.
\end{attention}

\courssub{TP 5 — Esprit critique : analyser des messages WhatsApp}

\textbf{Objectif.} Appliquer la démarche critique — \cle{source}, \cle{preuve}, \cle{cohérence} — à trois messages reçus sur WhatsApp, situation quotidienne des élèves camerounais.

\textbf{Les trois messages.} (a) Une rumeur : « Transfère ce message à dix personnes, sinon tu auras des malheurs. » (b) Une publicité mensongère : « Boire de l'eau salée guérit toutes les maladies. » (c) Une information vérifiée : « Le ministère de la Santé publique annonce une campagne gratuite de vaccination, dates et centres officiels joints. »

\begin{exemplebox}[Analyse résolue du message (b)]
\textbf{Source :} anonyme, aucun médecin, aucun hôpital ni laboratoire identifié. \textbf{Preuve :} aucune étude, aucun chiffre, aucun témoignage vérifiable. \textbf{Cohérence :} l'affirmation contredit l'ensemble des connaissances médicales établies (aucune substance unique ne guérit « toutes » les maladies) ; un excès de sel est même dangereux pour l'hypertension. \textbf{Verdict :} fausse information, à ne pas transmettre.
\end{exemplebox}

\begin{popfigure}
\begin{center}
\begin{tikzpicture}[node distance=0.5cm, every node/.style={font=\small}]
\node[draw=popink, fill=popcream, rounded corners, text width=0.9\linewidth, align=center] (m) at (0,3) {\textbf{Message reçu} : « Boire de l'eau salée guérit toutes les maladies »};
\node[draw=popblue, fill=popblueL, rounded corners, text width=0.28\linewidth, align=center] (s) at (-4.3,1.2) {\textbf{Source ?}\\Anonyme : suspecte};
\node[draw=poporange, fill=poporangeL, rounded corners, text width=0.28\linewidth, align=center] (p) at (0,1.2) {\textbf{Preuve ?}\\Aucune étude : absente};
\node[draw=popgreen, fill=popgreenL, rounded corners, text width=0.28\linewidth, align=center] (c) at (4.3,1.2) {\textbf{Cohérence ?}\\Contredit la médecine};
\node[draw=popink, fill=poppinkL, rounded corners, text width=0.5\linewidth, align=center] (v) at (0,-0.6) {\textbf{Verdict : fausse information} — ne pas croire, ne pas transférer};
\draw[-{Stealth}, thick, popmuted] (m) -- (s);
\draw[-{Stealth}, thick, popmuted] (m) -- (p);
\draw[-{Stealth}, thick, popmuted] (m) -- (c);
\draw[-{Stealth}, thick, popmuted] (s) -- (v);
\draw[-{Stealth}, thick, popmuted] (p) -- (v);
\draw[-{Stealth}, thick, popmuted] (c) -- (v);
\end{tikzpicture}
\end{center}
\legende{Modèle de fiche d'analyse critique d'un message : trois questions (source, preuve, cohérence) conduisent au verdict. À reproduire pour chaque message du TP.}
\end{popfigure}

\begin{exoresolu}[Analyse des messages (a) et (c)]
Énoncé : applique la fiche d'analyse aux messages (a) et (c), puis conclus sur le comportement à adopter.
\tcblower
Solution détaillée : \textbf{Message (a)} — source anonyme, aucune preuve (aucun lien rationnel entre transférer un message et subir un malheur), incohérent avec toute connaissance causale : verdict, superstition manipulatoire ; on supprime et on ne transfère pas. \textbf{Message (c)} — source identifiable et officielle (ministère de la Santé publique), preuves vérifiables (dates, lieux, possibilité de confirmer auprès du centre de santé), cohérent avec les politiques publiques de vaccination : verdict, information fiable, à condition de vérifier sur les canaux officiels. Conclusion générale : l'esprit critique n'est pas le doute systématique, mais l'exigence de raisons avant de croire et de transmettre.
\end{exoresolu}

\courssub{Restitution et évaluation}

Chaque groupe présente sa production devant la classe (5 minutes par TP) : définition du TP 1, tableau du TP 2, synthèse du rapporteur pour le TP 3, paragraphe du TP 4, fiches d'analyse du TP 5. La correction collective confronte les productions à la grille d'évaluation suivante.

\begin{center}
\begin{tabularx}{\linewidth}{|l|X|c|}
\hline
\rowcolor{popblueL} \textbf{Critère} & \textbf{Attendu} & \textbf{Points} \\
\hline
Exactitude & Les notions (étonnement, examen, argument, thèse) sont correctement employées & /4 \\
\hline
Justification & Chaque affirmation est appuyée par un argument et un exemple & /4 \\
\hline
Méthode & Les étapes (analyse, comparaison, débat, rédaction, critique) sont respectées & /4 \\
\hline
Clarté & Expression correcte, organisation lisible de la production & /4 \\
\hline
Coopération & Répartition des rôles et participation de chaque membre & /4 \\
\hline
\end{tabularx}
\end{center}

\begin{exercice}[Pour aller plus loin]
Choisis un message reçu cette semaine sur ton téléphone (information, rumeur ou publicité) et rédige sa fiche d'analyse critique complète (source, preuve, cohérence, verdict), puis un court paragraphe argumenté répondant à la question : « Faut-il croire tout ce qui circule sur les réseaux sociaux ? »
\end{exercice}

\begin{corrige}[Exercice — Pour aller plus loin]
La fiche doit comporter les quatre rubriques remplies de manière précise : la source est nommée ou qualifiée (anonyme, officielle, commerciale) ; la preuve est recherchée (étude, chiffre, témoin vérifiable) ou son absence est constatée ; la cohérence confronte le message aux connaissances établies ; le verdict est explicite et motivé. Le paragraphe attendu suit la structure thèse — arguments — exemples — conclusion : par exemple, thèse « non, il ne faut pas tout croire, car croire sans vérifier, c'est renoncer à son entendement (Kant) » ; arguments : la facilité de fabriquer une fausse information, les intérêts cachés de ceux qui la diffusent ; exemple : la rumeur de l'eau salée ; conclusion : l'esprit critique, héritier de l'examen socratique, est la condition d'une vie libre et responsable, y compris numérique.
\end{corrige}

\begin{aretenir}[L'essentiel des TP]
Philosopher est une activité qui s'apprend par la pratique : analyser un texte en cinq étapes, comparer méthodiquement philosophie et science, débattre en respectant des règles de parole, rédiger selon la structure thèse — arguments — exemples — conclusion, et soumettre toute information au triple examen de la source, de la preuve et de la cohérence. Ces cinq gestes, appliqués aux situations camerounaises quotidiennes, transforment les définitions de la leçon en compétences durables.
\end{aretenir}

\newpage

\coursec{Activité d'intégration}

\coursec{LEÇON N20 : Qu'est-ce que la Philosophie ?}

\courssub{Définitions de la philosophie}

\begin{definition}[Étymologie et sens premier]
Le mot \emph{philosophie} vient du grec \emph{philein} (aimer, rechercher) et \emph{sophia} (sagesse). Étymologiquement, la philosophie est l'\cle{amour de la sagesse}. Le philosophe n'est pas celui qui possède le savoir (le \emph{sophos}), mais celui qui en est épris et le cherche sans cesse. La sagesse visée n'est pas seulement théorique : elle vise une vie bonne et lucide.
\end{definition}

\begin{definition}[L'examen critique (Socrate)]
D'après Socrate (IV\textsuperscript{e} s. av. J.-C.), rapporté par Platon dans l'\emph{Apologie de Socrate}, la philosophie est un \cle{examen critique} permanent. « Une vie sans examen ne vaut pas la peine d'être vécue. » Philosopher, c'est refuser les opinions toutes faites (\emph{doxa}), interroger les certitudes apparentes et tester la solidité des arguments. C'est la \cle{maïeutique} : accoucher les esprits de leurs contradictions pour faire naître une connaissance plus sûre.
\end{definition}

\begin{definition}[L'autonomie de la raison (Kant)]
Pour Emmanuel Kant (\emph{Réponse à la question : qu'est-ce que les Lumières ?}, 1784), la philosophie est le courage de \cle{penser par soi-même} (\emph{Sapere aude}). « Aie le courage de te servir de ton propre entendement ! » La philosophie est l'activité par laquelle le sujet sort de sa \emph{minorité} (incapacité de se servir de son entendement sans la direction d'autrui) pour accéder à la majorité rationnelle. L'autonomie est la source de la dignité humaine.
\end{definition}

\begin{aretenir}[Synthèse : trois visages de la philosophie]
\begin{itemize}
\item \textbf{Amour de la sagesse} : quête inachevée, désir de comprendre pour mieux vivre.
\item \textbf{Examen critique (Socrate)} : méthode du questionnement, refus du dogmatisme, lucidité sur son ignorance.
\item \textbf{Autonomie rationnelle (Kant)} : usage public de sa raison, responsabilité intellectuelle, fondement de la liberté.
\end{itemize}
\end{aretenir}

\courssub{Origine et objet de la philosophie}

\begin{propriete}[Naissance en Grèce (VI\textsuperscript{e} s. av. J.-C.)]
La philosophie naît dans les cités ioniennes (Milet) avec Thalès, Anaximandre, Anaximène. Elle marque le passage du \emph{mythos} (récit légendaire, arbitraire, autorité de la tradition) au \emph{logos} (discours rationnel, argumenté, universel). On cherche des explications naturelles aux phénomènes (l'eau, l'air, le feu comme principes) au lieu d'en appeler aux volontés divines capricieuses.
\end{propriete}

\begin{propriete}[L'étonnement, principe moteur (Aristote)]
Dans la \emph{Métaphysique}, Aristote affirme : « C'est par l'étonnement que les hommes, maintenant comme d'abord, se sont mis à philosopher. » L'\cle{étonnement} n'est pas simple surprise ; c'est la prise de conscience que l'évident ne va pas de soi. Il déclenche l'interrogation : « Pourquoi y a-t-il quelque chose plutôt que rien ? », « Qu'est-ce que le juste ? », « Comment connaître ? ».
\end{propriete}

\begin{aretenir}[Objet de la philosophie]
La philosophie n'a pas d'objet régional déterminé (contrairement aux sciences). Son objet est \textbf{universel} : elle interroge les \cle{principes premiers} et le \cle{sens} de la réalité totale. Quatre grands domaines :
\begin{enumerate}
\item \textbf{L'homme} : nature, liberté, conscience, société, histoire.
\item \textbf{Le monde} : matière, espace, temps, causalité, existence de Dieu.
\item \textbf{La connaissance} : vérité, erreur, science, opinion, limites de la raison.
\item \textbf{L'action} : morale, droit, politique, technique, art, valeur.
\end{enumerate}
\end{aretenir}

\courssub{La méthode de la philosophie}

\begin{methode}[Démarche philosophique type (exigible au Probatoire)]
La philosophie ne suit pas une méthode unique figée, mais une \cle{démarche rigoureuse} en cinq temps articulés :
\begin{enumerate}
\item \textbf{L'étonnement / La problématisation} : face à une opinion, un fait, un paradoxe, formuler une \cle{question philosophique} précise (ni trop large, ni technique).
\item \textbf{Le doute méthodique (Descartes)} : suspendre son assentiment aux opinions reçues, aux préjugés, aux évidences non fondées. Ce n'est pas un scepticisme destructeur, mais un \cle{outil de tri} (hyperbolique chez Descartes, permanent chez les critiques).
\item \textbf{L'analyse conceptuelle} : définir rigoureusement les termes clés, distinguer les sens, repérer les présupposés, construire des distinctions opérantes (ex. : fait / droit, subjectif / objectif, nécessaire / contingent).
\item \textbf{La discussion argumentée (dialectique)} : confronter les thèses contraires (thèse / antithèse), examiner les objections, tester la cohérence logique, chercher des contre-exemples. L'argument d'autorité ne vaut pas preuve.
\item \textbf{La conclusion provisoire} : formuler une réponse nuancée, justifiée, consciente de ses limites. La vérité philosophique est toujours \cle{revisable} : elle progresse par rectification, non par accumulation dogmatique.
\end{enumerate}
\end{methode}

\begin{attention}[Pièges fréquents]
\begin{itemize}
\item Confondre \emph{doute méthodique} (outil de vérité) et \emph{scepticisme radical} (refus de toute vérité).
\item Remplacer l'argumentation par l'affirmation sentimentale (« Je pense que... », « C'est mon avis »).
\item Croire que la philosophie n'a pas de méthode parce qu'elle n'a pas de laboratoire : son laboratoire, c'est le \cle{langage} et la \cle{raison}.
\end{itemize}
\end{attention}

\courssub{L'importance de la philosophie}

\begin{aretenir}[Pourquoi philosopher ?]
\begin{itemize}
\item \textbf{Former le jugement critique} : distinguer le vrai du plausible, l'argument du sophisme, la preuve de l'assertion. Indispensable à l'ère des réseaux sociaux et de l'infobésité.
\item \textbf{Conquérir l'autonomie intellectuelle} : ne plus subir les opinions dominantes, la publicité, la propagande, la rumeur. Devenir auteur de sa propre pensée (Kant).
\item \textbf{Apprendre à argumenter et dialoguer} : structurer sa pensée, écouter l'adversaire, accepter la contradiction comme progrès. Base de la vie démocratique et professionnelle.
\item \textbf{Donner du sens à l'existence} : clarifier ses valeurs, choisir sa vie en connaissance de cause, affronter l'absurde, la mort, la souffrance (Sagesse antique, Existentialisme).
\item \textbf{Protéger contre la manipulation} : celui qui ne sait pas \emph{pourquoi} il croit est la proie de celui qui sait \emph{comment} le faire croire.
\end{itemize}
\end{aretenir}

\courssub{Philosophie et science}

\begin{propriete}[Points communs]
\begin{itemize}
\item \textbf{Recherche de la vérité} : toutes deux visent l'objectivité, refusent l'arbitraire et l'irrationnel.
\item \textbf{Exigence de justification} : une affirmation scientifique comme philosophique doit être \cle{argumentée}, \cle{démontrée} ou au moins \cle{justifiée} par des raisons communicables.
\item \textbf{Rationalité} : usage de la logique, des concepts, de la déduction et de l'induction.
\end{itemize}
\end{propriete}

\begin{propriete}[Différences fondamentales]
\begin{center}
\begin{tabularx}{\linewidth}{|l|X|X|}
\hline
\rowcolor{popblueL} \textbf{Critère} & \textbf{Science} & \textbf{Philosophie} \\
\hline
Objet & \textbf{Particulier / Régional} : faits observables, lois mathématiques, mécanismes causels (ex. : la cellule, l'atome, l'économie). & \textbf{Universel / Fondamental} : sens de l'être, conditions de possibilité de la science, valeurs, subjectivité. \\
\hline
Méthode de validation & \textbf{Expérience contrôlée} : observation, hypothèse, expérience, reproductibilité, falsifiabilité (Popper). Preuve \emph{empirique}. & \textbf{Analyse rationnelle} : conceptualisation, argumentation logique, critique des présupposés, cohérence. Preuve \emph{discursive}. \\
\hline
Résultat & \textbf{Connaissances cumulatives, techniques, prédictives} : lois, théories, applications technologiques. & \textbf{Sagesse, lucidité, clarification} : pas de « résultat » technique, mais une meilleure compréhension de l'homme et du monde. \\
\hline
Rapport au sujet & Le sujet doit s'effacer (objectivité, protocole). & Le sujet \emph{penseur} est au centre (réflexivité, engagement). \\
\hline
\end{tabularx}
\end{center}
\end{propriete}

\begin{exemplebox}[Exemple concret : l'intelligence artificielle]
\begin{itemize}
\item \textbf{Science (Informatique / IA)} : conçoit des algorithmes, mesure la performance (précision, vitesse), modélise l'apprentissage automatique. Question : « Comment faire pour que la machine reconnaisse un visage ? »
\item \textbf{Philosophie (Éthique / Épistémologie)} : interroge la responsabilité (qui est responsable de l'erreur ?), le biais (la machine reproduit-elle nos préjugés ?), la nature de l'intelligence (simulation vs compréhension), l'impact sur le travail et la liberté. Question : « Que signifie "penser" pour une machine ? Est-ce désirable ? »
\end{itemize}
\end{exemplebox}

\courssub{TP guidé : philosopher sur des situations camerounaises}

\textbf{Objectifs de l'activité}\par
À l'issue de cette activité d'intégration, l'élève sera capable de :
\begin{itemize}
\item mobiliser les \cle{définitions} de la philosophie (amour de la sagesse, examen critique, étonnement) pour analyser une situation concrète de la vie courante ;
\item appliquer la \cle{méthode} philosophique (doute méthodique, analyse, discussion, argumentation) à des documents authentiques ;
\item distinguer une information vérifiée d'une simple rumeur, en montrant l'\cle{importance} de la philosophie dans la formation du jugement ;
\item préciser le rapport entre la \cle{philosophie} et la \cle{science} (recherche de la vérité, exigence de preuve) ;
\item rédiger et justifier une réponse argumentée sous forme d'un paragraphe structuré (thèse, argument, exemple, conclusion), compétence exigée au Probatoire et au Baccalauréat technique.
\end{itemize}

\textbf{Situation}\par
\textbf{La rumeur du marché : un concours de jeunes philosophes.}

Nous sommes un lundi matin au Lycée technique de Deïdo, à Douala. À 7 h 15, un message anonyme circule massivement dans les groupes WhatsApp des classes de Première technique. Le voici :

\begin{quote}
\textbf{Document 1 (message WhatsApp anonyme)} : « Alerte !!! Info sûre d'un ami qui travaille à Yaoundé : TOUTES les notes du Probatoire technique de cette année seront ANNULÉES à cause des fuites. Inutile de réviser, l'examen sera refait avec d'autres sujets. Partagez vite avant que le message soit supprimé !!! »
\end{quote}

En moins de deux heures, le message a été transféré dans 12 groupes de classe. Sur les 48 élèves de la classe de Première F4, 31 déclarent avoir « un peu » ou « beaucoup » peur, et 9 affirment qu'ils vont arrêter de réviser « puisque c'est inutile ». Une rumeur semblable avait déjà circulé l'année précédente, sans qu'aucune note ne soit jamais annulée.

À 10 h 30, le proviseur affiche à la porte de la salle des professeurs le document suivant :

\begin{quote}
\textbf{Document 2 (communiqué officiel)} : « L'Office du Baccalauréat du Cameroun (OBC) informe l'opinion publique qu'aucune annulation d'épreuve du Probatoire technique n'est envisagée. Les épreuves se dérouleront conformément au calendrier publié. Toute information contraire relève de l'intoxication. Fait à Yaoundé, le directeur général de l'OBC. »
\end{quote}

Le professeur de philosophie saisit l'occasion pour organiser le « concours de jeunes philosophes ». Il distribue deux textes supplémentaires :

\begin{quote}
\textbf{Document 3 (Socrate, d'après l'\emph{Apologie de Socrate} de Platon)} : « Une vie sans examen ne vaut pas la peine d'être vécue. » Socrate passait ses journées sur la place publique d'Athènes à interroger ceux qui prétendaient savoir, afin de montrer que leurs certitudes n'étaient souvent que des opinions non fondées.

\textbf{Document 4 (Kant, \emph{Réponse à la question : qu'est-ce que les Lumières ?}, 1784)} : « Aie le courage de te servir de ton propre entendement ! » Penser par soi-même, c'est refuser de recevoir passivement les opinions des autres et soumettre toute affirmation à l'examen de sa propre raison.
\end{quote}

\textbf{Problème} : comment la philosophie peut-elle aider les élèves de Première F4 à démêler le vrai du faux et à prendre une décision raisonnable ?

\textbf{Consignes}\par
Par groupes de 4 élèves, puis individuellement pour la rédaction finale, réponds aux tâches suivantes :

\begin{enumerate}
\item \textbf{Analyse critique des messages.} À l'aide de la grille \emph{source -- preuve -- cohérence}, compare le document 1 et le document 2 : qui parle ? quelle preuve est fournie ? le message est-il cohérent avec les faits connus (notamment la rumeur de l'année précédente) ?
\item \textbf{Définition de la philosophie.} À partir des documents 3 et 4, dégage deux définitions de la philosophie : l'une liée à l'\emph{examen} (Socrate), l'autre liée à l'\emph{autonomie de la pensée} (Kant). Formule chaque définition en une phrase avec tes propres mots.
\item \textbf{Méthode philosophique.} Indique quelles étapes de la méthode philosophique (étonnement, doute, analyse, discussion, conclusion) les élèves de la classe auraient dû suivre avant de paniquer.
\item \textbf{Mini-débat.} En 10 minutes, deux groupes s'opposent sur la question : « Faut-il croire ce qui circule sur les réseaux sociaux ? » Chaque groupe doit avancer au moins deux arguments et un exemple camerounais.
\item \textbf{Rédaction individuelle.} Rédige un paragraphe argumenté d'environ 15 lignes répondant à la question : « Pourquoi faut-il examiner les informations avant d'y croire ? » Ton paragraphe doit contenir une thèse, deux arguments, un exemple tiré de la situation et une conclusion.
\end{enumerate}

\textbf{Ressources à mobiliser}\par

\begin{aretenir}[Ce qu'il faut se rappeler]
\begin{itemize}
\item \textbf{Définitions} : le mot \emph{philosophie} vient du grec \emph{philein} (aimer) et \emph{sophia} (sagesse) : la philosophie est l'\cle{amour de la sagesse}, non la possession du savoir. Elle est aussi \cle{examen critique} des opinions (Socrate) et usage autonome de la raison (Kant).
\item \textbf{Origine et objet} : née en Grèce au VI\textsuperscript{e} siècle av. J.-C., la philosophie naît de l'\cle{étonnement} (Aristote) et a pour objet la recherche de la \cle{vérité} sur l'homme, le monde, la connaissance et l'action.
\item \textbf{Méthode} : étonnement $\rightarrow$ doute méthodique $\rightarrow$ analyse des notions $\rightarrow$ discussion argumentée $\rightarrow$ conclusion provisoire. Le doute n'est pas un refus de tout croire, mais un \emph{instrument} pour trier le vrai du faux.
\item \textbf{Importance} : la philosophie forme le jugement, rend libre, apprend à argumenter et protège contre la manipulation.
\item \textbf{Philosophie et science} : toutes deux recherchent la vérité et exigent des preuves ; la science vérifie par l'expérience, la philosophie interroge le sens et les fondements. Une information sérieuse, comme un résultat scientifique, doit pouvoir être \emph{justifiée}.
\end{itemize}
\end{aretenir}

\textbf{Démarche et corrigé commenté}\par

\begin{exoresolu}[Corrigé commenté de l'activité]
\textbf{Tâche 1 : analyse source -- preuve -- cohérence.}

\begin{center}
\begin{tabularx}{\linewidth}{|l|X|X|}
\hline
\rowcolor{popblueL} \textbf{Critère} & \textbf{Document 1 (WhatsApp)} & \textbf{Document 2 (OBC)} \\
\hline
Source & Anonyme, « un ami qui travaille à Yaoundé » : non identifiable, donc non vérifiable. & Institution officielle et compétente : l'OBC organise l'examen. \\
\hline
Preuve & Aucune : simple affirmation assortie d'un appel à partager « vite ». & Référence au calendrier officiel publié ; document signé et affiché. \\
\hline
Cohérence & Contradite par les faits : la rumeur de l'année précédente était fausse ; l'argument « partagez avant suppression » vise à court-circuiter la réflexion. & Cohérent avec le fonctionnement normal de l'examen et vérifiable auprès de l'administration. \\
\hline
\end{tabularx}
\end{center}

\emph{Commentaire} : la grille montre que la crédibilité d'un message dépend de trois conditions : une source identifiable et compétente, une preuve accessible, une cohérence avec les faits établis. Le document 1 ne remplit aucune de ces conditions.

\tcblower

\textbf{Tâche 2 : deux définitions de la philosophie.}

\begin{itemize}
\item D'après Socrate (document 3) : \emph{la philosophie est l'examen critique de soi-même et des opinions reçues} ; croire sans examiner, c'est vivre sans réfléchir.
\item D'après Kant (document 4) : \emph{la philosophie est le courage de penser par soi-même}, c'est-à-dire de soumettre toute affirmation au tribunal de sa propre raison au lieu de répéter ce que disent les autres.
\end{itemize}

\emph{Commentaire} : ces deux définitions se complètent : l'examen (Socrate) est l'outil, l'autonomie (Kant) est le but. Appliquées à la situation, elles interdisent de croire la rumeur sans l'avoir soumise à la critique.

\textbf{Tâche 3 : la méthode appliquée.}

\begin{enumerate}
\item \textbf{Étonnement} : « D'où vient ce message ? Est-ce possible ? » — la surprise doit devenir question, non panique.
\item \textbf{Doute méthodique} : suspendre son jugement tant qu'aucune preuve n'est fournie.
\item \textbf{Analyse} : appliquer la grille source -- preuve -- cohérence (tâche 1).
\item \textbf{Discussion} : confronter les points de vue en classe, écouter le communiqué officiel.
\item \textbf{Conclusion provisoire} : continuer à réviser, car l'information officielle, seule vérifiable, confirme le calendrier.
\end{enumerate}

\textbf{Tâche 4 : éléments du débat.}

\emph{Pour la méfiance} : les réseaux diffusent des messages anonymes, invérifiables, conçus pour émouvoir (ex. : fausses annonces de concours, rumeurs de pénurie au marché Mokolo). \emph{Nuance} : les réseaux peuvent aussi relayer des informations vraies et utiles (communiqués officiels, alertes météo) ; il ne faut donc pas tout rejeter, mais tout \emph{examiner} — c'est exactement la position philosophique.

\textbf{Tâche 5 : modèle de paragraphe argumenté (corrigé type).}

\begin{quote}
Il faut examiner les informations avant d'y croire, parce que croire sans examen, c'est renoncer à penser par soi-même. D'abord, toute affirmation n'est vraie que si elle peut être justifiée : comme la science n'accepte un résultat qu'avec une preuve expérimentale, la philosophie n'accepte une opinion qu'avec des raisons ; or le message anonyme reçu au Lycée technique de Deïdo ne cite ni source identifiable ni preuve, contrairement au communiqué de l'OBC. Ensuite, l'examen nous protège de la manipulation : Socrate montrait déjà que ceux qui prétendent savoir ne font souvent que répéter des opinions, et Kant nous invite à « avoir le courage de nous servir de notre propre entendement » ; les 9 élèves prêts à abandonner leurs révisions montrent le danger réel d'une rumeur non examinée : elle peut gâcher une année de travail. Enfin, la rumeur de l'année précédente, qui s'était révélée fausse, prouve que la répétition d'un message ne le rend pas vrai. Ainsi, examiner avant de croire n'est pas du scepticisme stérile : c'est l'exercice même de la liberté de penser, et la première leçon de la philosophie.
\end{quote}

\emph{Commentaire méthodologique} : le paragraphe suit le plan exigé à l'examen : thèse annoncée, argument 1 (exigence de preuve, lien philosophie-science), argument 2 (protection contre la manipulation, appui sur Socrate et Kant), exemple chiffré tiré de la situation, conclusion qui répond exactement à la question posée.
\end{exoresolu}

\textbf{Bilan de l'activité}\par
Cette activité a permis de mettre en évidence que la philosophie n'est pas un savoir abstrait réservé aux livres : elle est un \cle{outil pratique} de survie intellectuelle dans la vie quotidienne camerounaise, où rumeurs et messages viraux circulent chaque jour. Les élèves ont d'abord vérifié que la méthode philosophique — étonnement, doute, analyse, discussion, conclusion — suffit à démonter une fausse information, sans aucun recours à la force ou à la peur. Ils ont ensuite compris que les définitions de la philosophie (amour de la sagesse, examen socratique, autonomie kantienne) se complètent : philosopher, c'est refuser de croire sur parole et exiger des raisons. Enfin, le rapprochement avec la science a montré que vérité et preuve vont ensemble : une information, comme un résultat scientifique, n'a de valeur que si elle peut être justifiée. La rédaction du paragraphe argumenté a en outre entraîné les élèves à la dissertation philosophique telle qu'elle sera exigée au Probatoire et au Baccalauréat technique : une thèse claire, des arguments fondés sur des auteurs, un exemple précis et une conclusion personnelle. L'affichage des meilleurs paragraphes au tableau valorise le travail collectif et installe durablement le réflexe : \emph{avant de croire, examine}.

\newpage

\coursec{Méthodes et exercices résolus}

\coursec{Leçon 20 : Qu'est-ce que la Philosophie ?}

\begin{definition}[La Philosophie : définition étymologique et conceptuelle]
Étymologiquement, le mot \emph{philosophie} vient du grec \emph{philein} (aimer, rechercher, désirer) et \emph{sophia} (sagesse, savoir suprême). La philosophie est donc littéralement « \cle{amour de la sagesse} ». Selon la tradition (Diogène Laërce, \emph{Vies des philosophes}), c'est \textbf{Pythagore} (VI\textsuperscript{e} s. av. J.-C.) qui aurait forgé ce terme : interrogé sur s'il était un sage, il aurait répondu qu'il n'était qu'un « philosophe », c'est-à-dire un ami de la sagesse, car la sagesse parfaite n'appartient qu'aux dieux. L'homme ne peut que la \emph{rechercher}.

Conceptuellement, la philosophie se définit comme une \cle{réflexion rationnelle, critique et universelle} sur l'homme, le monde et les valeurs, qui vise la \cle{sagesse} (conduite de vie lucide) et la \cle{vérité}. Elle se distingue de l'opinion (\emph{doxa}) par l'exigence de justification argumentée.
\end{definition}

\begin{aretenir}[Formule synthétique à retenir pour l'examen]
« La philosophie est née de l'\cle{étonnement}, procède par le \cle{doute} et le questionnement rationnel, et vise la \cle{sagesse}. »
\end{aretenir}

\courssub{Fiche méthode 1 : Appliquer les notions de l'unité à une situation concrète (Savoir-faire)}

\begin{methode}[Analyser une situation de la vie courante avec les notions de la leçon]
Pour mobiliser correctement les notions (\cle{étonnement}, \cle{doute}, \cle{esprit critique}, \cle{réflexion}, \cle{sagesse}) face à une situation concrète :
\begin{enumerate}
\item \textbf{Lire attentivement la situation} et souligner les faits, les paroles ou les comportements décrits.
\item \textbf{Identifier le problème philosophique} : s'agit-il d'une croyance acceptée sans examen (\emph{doxa}) ? D'un préjugé ? D'une question sur le sens de la vie, la justice, la vérité, le bonheur ?
\item \textbf{Choisir la ou les notions pertinentes} du cours : étonnement (Platon, Aristote), doute méthodique (Descartes), examen critique (Kant), amour de la sagesse (Pythagore).
\item \textbf{Appliquer la notion au cas concret} : montrer \emph{comment} le personnage philosophe (ou devrait philosopher), en citant des éléments précis de l'énoncé.
\item \textbf{Conclure} en dégageant la portée philosophique : en quoi cette situation montre que la philosophie est une attitude utile dans la vie quotidienne.
\end{enumerate}
\textbf{Réflexe rédactionnel (obligatoire au Probatoire) :} toujours articuler la réponse en trois temps : \emph{Notion définie} $\rightarrow$ \emph{Application au cas} $\rightarrow$ \emph{Conclusion philosophique}.
\end{methode}

\begin{exoresolu}[Application : le marché et la rumeur (Contexte camerounais)]
\textbf{Énoncé.} Au marché de Mokolo à Yaoundé, une rumeur affirme que les prix du gari vont doubler la semaine suivante. Sans vérifier, Mama Béa achète en panique dix sacs. Son fils Étienne, élève en classe de Première technique, lui demande d'abord : « Qui l'a dit ? Où est la preuve ? ». Identifiez la notion philosophique illustrée par l'attitude d'Étienne et montrez en quoi cette attitude est philosophique.
\tcblower
\textbf{Solution détaillée.}
\begin{enumerate}
\item \textbf{Notion définie.} L'attitude d'Étienne illustre le \cle{doute méthodique} et l'\cle{esprit critique}, qui sont au cœur de la démarche philosophique. Philosopher, c'est refuser d'accepter une affirmation sans l'avoir examinée rationnellement : \textbf{Descartes} (\emph{Discours de la méthode}) fait du doute le point de départ de la recherche de la vérité. Le philosophe \emph{suspend son jugement} (\emph{epochè}) tant qu'il n'a pas de raisons suffisantes de croire.
\item \textbf{Application au cas.} Mama Béa agit sous l'empire de la \emph{doxa} (l'opinion commune, la rumeur) : elle croit sans preuve et agit par peur (passion). Étienne, au contraire, pose deux questions rationnelles : la question de la \emph{source} (« Qui l'a dit ? » $\rightarrow$ critique de l'autorité) et la question de la \emph{justification} (« Où est la preuve ? » $\rightarrow$ exigence de raison). Il refuse de prendre l'opinion pour la vérité : c'est exactement le passage de la croyance au savoir rationnel qui définit l'esprit philosophique depuis \textbf{Socrate}.
\item \textbf{Conclusion.} Cette situation montre que philosopher n'est pas réservé aux livres : c'est une attitude concrète d'examen critique qui permet, dans la vie quotidienne camerounaise (marché, réseaux sociaux, politique), de ne pas être victime des rumeurs (\emph{fake news}) et de prendre des décisions fondées sur des raisons plutôt que sur la peur ou l'imitation.
\end{enumerate}
\end{exoresolu}

\courssub{Fiche méthode 2 : Définir la philosophie et mobiliser les définitions des auteurs}

\begin{methode}[Construire une définition argumentée de la philosophie (Type dissertation / Question de cours)]
\begin{enumerate}
\item \textbf{Rappeler l'étymologie} : \emph{philein} (aimer/rechercher) + \emph{sophia} (sagesse) $\rightarrow$ « Amour de la sagesse » (Pythagore). Préciser : le philosophe ne \emph{possède} pas la sagesse, il la \emph{cherche}.
\item \textbf{Distinguer les définitions selon les points de vue} :
    \begin{itemize}
    \item Par l'\textbf{origine} : l'\cle{étonnement} (Platon, \emph{Théétète} ; Aristote, \emph{Métaphysique}).
    \item Par la \textbf{méthode} : le \cle{doute} (Descartes), l'\cle{examen critique} (Kant, « Sapere aude »).
    \item Par la \textbf{finalité} : la sagesse (conduite de vie), le bonheur, la vérité, l'autonomie.
    \end{itemize}
\item \textbf{Citer brièvement un auteur} pour appuyer chaque aspect (nom + œuvre + idée clé), sans réciter de longues citations.
\item \textbf{Synthétiser} en une définition complète : « La philosophie est une réflexion rationnelle, critique et universelle sur l'homme, le monde et les valeurs, qui vise la sagesse. »
\end{enumerate}
\end{methode}

\begin{exoresolu}[Exercice : définir la philosophie (Sujet type Probatoire)]
\textbf{Énoncé.} En vous appuyant sur l'étymologie du mot « philosophie » et sur deux auteurs étudiés, proposez une définition complète de la philosophie (10 à 15 lignes).
\tcblower
\textbf{Solution détaillée.}

Étymologiquement, le mot \emph{philosophie} vient du grec \emph{philein} (aimer, désirer, rechercher) et \emph{sophia} (sagesse, savoir suprême) : la philosophie est littéralement « \cle{amour de la sagesse} ». Selon la tradition rapportée par Diogène Laërce, c'est \textbf{Pythagore} qui aurait forgé ce terme : le philosophe ne prétend pas \emph{posséder} la sagesse (seuls les dieux la possèdent), il la \emph{recherche} humblement. La philosophie est donc d'abord une \cle{quête}, non un savoir achevé ou un dogme.

Cette quête naît de l'\textbf{étonnement} : selon \textbf{Platon} (\emph{Théétète}, 155d) et \textbf{Aristote} (\emph{Métaphysique}, A, 2), « l'étonnement est le commencement de la philosophie ». Devant le monde, l'existence, la mort, la justice, l'homme ne prend pas les choses pour acquises ; il s'interroge : « Pourquoi y a-t-il quelque chose plutôt que rien ? », « Qu'est-ce que le juste ? ».

Mais philosopher exige une méthode rigoureuse : le \textbf{doute rationnel}. \textbf{Descartes} (\emph{Méditations métaphysiques}) montre qu'il faut douter de tout ce qui peut l'être (sens, tradition, autorité) afin de parvenir à une vérité certaine (\emph{Cogito}). La philosophie est ainsi une \cle{réflexion critique} qui refuse les préjugés et les opinions reçues (\emph{doxa}).

\textbf{Conclusion :} on peut définir la philosophie comme une \emph{réflexion rationnelle et critique}, née de l'étonnement, qui interroge radicalement l'homme, le monde et les valeurs, et qui vise la \cle{sagesse} (lucidité sur l'existence) et la \cle{vérité}.
\end{exoresolu}

\courssub{Fiche méthode 3 : Expliquer l'origine et l'objet de la philosophie}

\begin{methode}[Traiter une question sur l'origine ou l'objet de la philosophie]
\begin{enumerate}
\item \textbf{Pour l'origine}, distinguer trois plans complémentaires :
    \begin{itemize}
    \item \textbf{Historique :} Grèce, VI\textsuperscript{e} s. av. J.-C., cités ioniennes (Milet), passage du \cle{mythe} au \cle{logos} (Thalès, Anaximandre, Héraclite). Naissance de la raison naturelle.
    \item \textbf{Psychologique :} L'\cle{étonnement} (Platon, Aristote) face au monde, à l'être, aux valeurs.
    \item \textbf{Social / Politique :} La \cle{liberté de parole} (\emph{parrhésia}) dans la cité (Athènes), le \cle{loisir} (\emph{scholè}) permettant de se soustraire aux nécessités vitales pour penser.
    \end{itemize}
\item \textbf{Pour l'objet}, préciser que la philosophie interroge ce que les sciences particulières ne traitent pas : le \cle{sens} de la vie, la \cle{vérité}, le \cle{bien}, le \cle{beau}, la \cle{justice}, \cle{Dieu}, la \cle{liberté}, la \cle{conscience}. Ce sont des questions \emph{universelles} (concernent tous les hommes) et \emph{radicales} (vont aux racines/racines).
\item \textbf{Illustrer} par un exemple concret opposant explication mythique et questionnement philosophique (ex: mythe africain de la pluie vs explication météorologique + question du sens de la sécheresse).
\item \textbf{Conclure} : la philosophie est un questionnement universel, rationnel et radical.
\end{enumerate}
\end{methode}

\begin{exoresolu}[Exercice : le passage du mythe au logos (Contexte Extrême-Nord)]
\textbf{Énoncé.} Dans un village de l'Extrême-Nord, on explique la sécheresse par la colère des ancêtres. Un élève affirme : « C'est une explication philosophique ». A-t-il raison ? Justifiez votre réponse en vous appuyant sur l'origine de la philosophie.
\tcblower
\textbf{Solution détaillée.}

\begin{enumerate}
\item \textbf{Rappel du cours (Origine historique).} La philosophie est née en Grèce au VI\textsuperscript{e} siècle av. J.-C. par le passage du \cle{mythe} au \cle{logos}.
    \begin{itemize}
    \item Le \emph{mythe} (\emph{muthos}) : explication par un récit mettant en scène des volontés divines, des esprits, des forces surnaturelles arbitraires. Il se \emph{raconte}, s'impose par la tradition/autorité, ne se discute pas rationnellement.
    \item Le \emph{logos} : explication par la \emph{raison}, la \emph{cause naturelle}, l'argumentation. Il se \emph{démontre}, se vérifie, se discute. \textbf{Thalès de Milet} cherche le principe (\emph{archè}) de toutes choses dans l'\emph{eau} (élément naturel), non dans la volonté de Poséidon.
    \end{itemize}
\item \textbf{Application au cas.} L'explication de la sécheresse par la « colère des ancêtres » relève du \emph{mythe} (ou de la croyance religieuse traditionnelle) : elle invoque une cause surnaturelle, intentionnelle, qui ne se soumet pas à la vérification rationnelle ni à la discussion critique. Ce n'est donc pas, à proprement parler, une explication \emph{philosophique}.
\item \textbf{Nuance essentielle (Origine psychologique).} Cependant, si les villageois (ou l'élève) se mettent à \emph{discuter} cette croyance, à en demander les preuves, à la confronter à l'explication climatique scientifique, alors une \emph{démarche philosophique} commence : philosopher, c'est précisément soumettre les croyances reçues au questionnement rationnel (étonnement $\rightarrow$ doute $\rightarrow$ recherche de causes).
\item \textbf{Conclusion.} L'élève a \textbf{tort} : l'explication mythique n'est pas philosophique \emph{en elle-même} (elle relève du mythe). Mais elle peut devenir le \emph{point de départ} (la matière) d'une réflexion philosophique si on l'interroge rationnellement, car la philosophie naît de l'étonnement et du doute face aux explications toutes faites.
\end{enumerate}
\end{exoresolu}

\courssub{Fiche méthode 4 : Décrire la méthode philosophique}

\begin{methode}[Présenter la méthode de la philosophie en étapes logiques (Savoir-faire : Rédiger et justifier une démarche)]
La méthode philosophique n'est pas une recette figée, mais une \cle{démarche rigoureuse} qu'il faut savoir énumérer, ordonner et illustrer :
\begin{enumerate}
\item \textbf{L'étonnement (Origine) :} Prise de conscience d'un problème qui ne va pas de soi. Refus de l'évidence. (Platon, Aristote).
\item \textbf{Le doute (Outil critique) :} Mise entre parenthèses (\emph{epochè}) des opinions reçues, préjugés, autorités, sens commun. \textbf{Descartes} : doute méthodique, hyperbolique, pour trouver un point d'appui certain.
\item \textbf{L'analyse conceptuelle (Rigueur) :} Définir les termes du problème, distinguer les sens (polysémie), décomposer le problème en sous-questions. \textbf{Socrate} : « Qu'est-ce que... ? » (Recherche de l'essence / définition).
\item \textbf{L'argumentation (Preuve) :} Avancer des \cle{raisons} (arguments logiques), des \cle{exemples} concrets, des \cle{contre-exemples}, examiner les \cle{objections} (antithèse). Distinguer argument d'autorité (faible) et argument de raison (fort).
\item \textbf{Le dialogue / La dialectique (Vérification) :} Confronter sa pensée à celle d'autrui. \textbf{Socrate} : \cle{maïeutique} (art d'accoucher les esprits par le questionnement). La vérité se construit dans la contradiction maîtrisée.
\item \textbf{La conclusion provisoire (Résultat) :} Formuler une \cle{thèse} argumentée, en conscience de sa fallibilité (on peut se tromper, la philosophie est \emph{interminable}).
\end{enumerate}
\textbf{À retenir :} La méthode philosophique est \emph{rationnelle} (procède par raisons, non par sentiments/autorité) et \emph{critique} (examine tout, y compris ses propres présupposés).
\end{methode}

\begin{exoresolu}[Exercice : identifier les étapes de la méthode (Cas concret)]
\textbf{Énoncé.} Lisez le texte suivant : « Ngono se demande pourquoi il faut obéir aux lois. D'abord, il refuse la réponse facile : "parce que les parents l'ont dit". Ensuite, il définit ce qu'est une loi et la distingue d'une coutume. Puis il avance des raisons : sans lois, la vie en société serait impossible. Enfin, il discute avec ses camarades qui ne sont pas d'accord. » Relevez dans ce texte les étapes de la méthode philosophique en les nommant.
\tcblower
\textbf{Solution détaillée.}

On identifie dans le comportement de Ngono les étapes suivantes de la démarche philosophique :
\begin{enumerate}
\item \textbf{L'étonnement / La problématisation :} « Ngono \emph{se demande} pourquoi il faut obéir aux lois. » Le questionnement initial, la mise en problème d'une évidence sociale, est le point de départ de toute philosophie (Platon, Aristote : l'étonnement).
\item \textbf{Le doute méthodique / Critique de l'autorité :} « Il \emph{refuse la réponse facile} : "parce que les parents l'ont dit". » Ngono met en doute l'opinion reçue (\emph{doxa}), l'argument d'autorité parentale ou traditionnelle. C'est l'attitude cartésienne : ne rien accepter comme vrai qu'on ne l'ait reconnu évidente.
\item \textbf{L'analyse conceptuelle / Définition :} « Il \emph{définit ce qu'est une loi et la distingue d'une coutume}. » La rigueur philosophique exige de clarifier les concepts : distinguer le droit (loi écrite, sanctionnée) de la coutume (usage social), définir l'obligation. Socrate exigeait toujours la définition (\emph{ti estin ?}) avant de discuter.
\item \textbf{L'argumentation / La justification :} « Il \emph{avance des raisons} : sans lois, la vie en société serait impossible. » Philosopher, c'est justifier sa thèse par des arguments logiques (ici : argument pragmatique / condition de possibilité de la vie sociale), non par des sentiments ou des ordres.
\item \textbf{Le dialogue / La confrontation dialectique :} « Il \emph{discute avec ses camarades qui ne sont pas d'accord}. » La confrontation des points de vue (thèse/antithèse) fait progresser la recherche de vérité (maïeutique socratique, dialectique hégélienne). Accepter la contradiction est une marque d'honnêteté intellectuelle.
\end{enumerate}
\textbf{Conclusion :} ce petit texte illustre la totalité de la démarche philosophique : étonnement, doute, analyse, argumentation, dialogue. Ngono, sans le savoir, philosophe : cela montre que la méthode philosophique est accessible à tout esprit rigoureux, hors de tout cadre scolaire.
\end{exoresolu}

\courssub{Fiche méthode 5 : Traiter le rapport entre Philosophie et Science (Point de programme spécifique)}

\begin{methode}[Construire un développement structuré sur Philosophie et Sciences]
\begin{enumerate}
\item \textbf{Définir les deux termes} dans l'introduction :
    \begin{itemize}
    \item \cle{Science} : Savoir \emph{objectif}, \emph{vérifiable} (expérimentalement), \emph{particulier} (objets définis : matière, vie, société), visant l'explication causale (\emph{comment ?}) et la prévision. Mathématisation fréquente.
    \item \cle{Philosophie} : Réflexion \emph{critique}, \emph{universelle} (porte sur le tout, les fondements, les valeurs), \emph{rationnelle} (argumentative), visant la \emph{compréhension du sens} (\emph{pourquoi ?}, \emph{pour quoi ?}) et la sagesse.
    \end{itemize}
\item \textbf{Dégager les ressemblances (Communauté d'origine et d'esprit)} : Toutes deux naissent de l'\cle{étonnement} (Aristote), rejettent le mythe/la croyance non fondée, procèdent par \cle{raison} (logos), visent la \cle{vérité} (certitude rationnelle), exigent la \cle{démonstration} / l'argumentation.
\item \textbf{Dégager les différences (Spécificité irréductible)} :
    \begin{center}
    \begin{tabularx}{\linewidth}{|l|X|X|}\hline
    \rowcolor{popblueL} \textbf{Critère} & \textbf{Science} & \textbf{Philosophie} \\ \hline
    \textbf{Objet} & Particulier, délimité (faits, phénomènes) & Universel, radical (sens, valeurs, fondements, totalité) \\ \hline
    \textbf{Méthode} & Expérimentale, observation, mesure, modélisation math. & Conceptuelle, analyse, dialectique, argumentation logique \\ \hline
    \textbf{Vérification} & Inter-subjective, reproductible en laboratoire & Logique (cohérence), pertinence argumentative, universalisable \\ \hline
    \textbf{Questions types} & \emph{Comment ?} (Mécanismes, causes efficientes) & \emph{Pourquoi ? / Pour quoi ?} (Finalité, sens, valeur, fondement) \\ \hline
    \textbf{Statut du résultat} & Provisoire, falsifiable (Popper), cumulatif & Ouvert, discuté, non cumulatif de la même façon \\ \hline
    \end{tabularx}
    \end{center}
\item \textbf{Montrer la complémentarité historique et actuelle} :
    \begin{itemize}
    \item \textbf{Histoire :} La philosophie est la « \cle{mère des sciences} » (Aristote, Descartes). Les sciences se sont détachées successivement (Physique $\rightarrow$ Galilée/Newton ; Biologie $\rightarrow$ Darwin ; Psychologie $\rightarrow$ Wundt/Freud ; Sociologie $\rightarrow$ Comte/Durkheim).
    \item \textbf{Épistémologie :} La philosophie des sciences interroge \emph{aujourd'hui} les fondements de la science : Qu'est-ce qu'une preuve ? Une loi scientifique ? Le hasard ? Le déterminisme ? (Ex: Mécanique quantique).
    \item \textbf{Éthique / Axiologie :} La science dit \emph{ce qui est possible} (pouvoir technique : clonage, IA, arme nucléaire, OGM). La philosophie demande \emph{ce qui est souhaitable / juste} (devoir éthique : faut-il le faire ?). Ex: Bioéthique, éthique de l'IA au Cameroun.
    \end{itemize}
\item \textbf{Conclure} par une position nuancée : Ni \emph{confusion} (scientisme : la science répond à tout), ni \emph{opposition totale} (irrationalisme : la philosophie n'a rien à voir avec le réel). \textbf{Complémentarité} : La science éclaire le philosophe sur le réel ; la philosophie éclaire le scientifique sur le sens et la responsabilité de son savoir.
\end{enumerate}
\end{methode}

\begin{exoresolu}[Exercice : Philosophie et Science (Sujet type Bac/Probatoire)]
\textbf{Énoncé.} Un élève déclare : « La science ayant tout expliqué, la philosophie ne sert plus à rien. » Discutez cette affirmation en vous appuyant sur le rapport entre philosophie et science.
\tcblower
\textbf{Solution détaillée.}

\begin{enumerate}
\item \textbf{Ce qui est vrai dans l'affirmation (Thèse partielle).} Il est historiquement exact que la science a progressivement expliqué de nombreux phénomènes autrefois abandonnés aux mythes ou à la spéculation philosophique : l'origine des espèces (Darwin), la structure de la matière (Physique quantique), les maladies (Microbiologie), le fonctionnement du cerveau (Neurosciences). La science apporte des réponses \emph{précises}, \emph{mesurables}, \emph{vérifiables} et \emph{efficaces} techniquement. En ce sens, elle a réduit le domaine de l'ignorance brute et relégué certaines « philosophies de la nature » au musée.
\item \textbf{Ce qui est faux dans l'affirmation (Antithèse : les limites internes de la science).} Mais la science ne peut pas « tout expliquer », par définition de son objet et de sa méthode :
    \begin{itemize}
    \item \textbf{Les questions de sens et de valeurs :} \emph{Que dois-je faire ? Qu'est-ce qu'une vie réussie ? La justice vaut-elle mieux que l'injustice ? La vie a-t-elle un sens ?} Ces questions \emph{axiologiques} (valeurs) et \emph{existentielles} ne relèvent pas de l'expérience mesurable. La science décrit \emph{ce qui est} (fait), elle ne prescrit pas \emph{ce qui doit être} (valeur) $\rightarrow$ \textbf{Hume} (guillotine de Hume : on ne déduit pas un « il faut » d'un « est »).
    \item \textbf{Les fondements de la science elle-même :} Qu'est-ce qu'une cause ? Une loi de la nature ? Le temps ? L'espace ? La scientificité d'une théorie ? (Problème de la démarcation, \textbf{Popper}). Ce sont des questions \emph{épistémologiques} (philosophie des sciences), non scientifiques.
    \item \textbf{L'indétermination éthique du progrès :} La technique (issue de la science) permet de manipuler le génome humain (CRISPR-Cas9), de créer des IA autonomes, de surveiller les populations. Mais \emph{faut-il} le faire ? Seule la réflexion philosophique (éthique, politique) peut répondre.
    \end{itemize}
\item \textbf{Synthèse (Complémentarité).} Philosophie et Sciences sont donc \emph{complémentaires} et non concurrentes :
    \begin{itemize}
    \item La science nous dit \emph{comment} le monde fonctionne (faits, mécanismes, pouvoir technique).
    \item La philosophie nous demande \emph{pourquoi} (sens, finalité) et \emph{pour quoi} (valeurs, responsabilité) nous en servons.
    \item La philosophie est historiquement la « mère des sciences » et reste leur « conscience critique ».
    \end{itemize}
\end{enumerate}
\textbf{Conclusion :} L'affirmation de l'élève relève du \cle{scientisme} (croyance que la science est le seul savoir valide). Elle est inacceptable : loin de rendre la philosophie inutile, le progrès fulgurant de la science et de la technique rend la réflexion philosophique sur le \emph{sens}, les \emph{valeurs} et la \emph{responsabilité} plus nécessaire que jamais pour éviter que l'homme ne devienne l'esclave de ses propres créations.
\end{exoresolu}

\courssub{Fiche méthode 6 : Rédiger et justifier une réponse argumentée (Type Probatoire / Bac)}

\begin{methode}[Rédiger une dissertation philosophique ou une question de réflexion structurée]
\begin{enumerate}
\item \textbf{Analyser le sujet (Brouillon) :} Définir \emph{chaque terme important} du sujet. Identifier l'opposition ou la tension. Reformuler le sujet en \textbf{question problématique} (ex: « La philosophie ne fait pas le pain » $\rightarrow$ « L'utilité de la philosophie se mesure-t-elle seulement à sa productivité matérielle ? »).
\item \textbf{Construire le plan (2 ou 3 parties) :} Schéma classique : \emph{Thèse} (l'évidence du sujet) / \emph{Antithèse} (les limites, le contre-argument) / \emph{Synthèse} (dépassement, nuance, articulation). Ou plan progressif : Analyser $\rightarrow$ Comprendre $\rightarrow$ Dépasser. \textbf{Annoncer le plan clairement en introduction.}
\item \textbf{Rédiger l'introduction (Modèle type) :}
    \begin{enumerate}
    \item Accroche / Contexte (citation, fait de société, opposition courante).
    \item Définition des termes clés du sujet.
    \item Problématisation (Question précise).
    \item Annonce du plan (« Nous montrerons d'abord que..., ensuite que..., enfin que... »).
    \end{enumerate}
\item \textbf{Développer les parties (Corps du devoir) :} Un paragraphe = Une idée force.
    \textbf{Schéma paragraphe obligatoire :} \emph{Idée directrice (phrase thèse)} $\rightarrow$ \emph{Argumentation (explication logique, référence auteur)} $\rightarrow$ \emph{Exemple concret (de préférence camerounais/africain)} $\rightarrow$ \emph{Mini-transition vers paragraphe suivant}.
\item \textbf{Justifier chaque affirmation :} \textbf{Interdiction d'affirmer sans raison.} « La philosophie est utile car... » $\rightarrow$ Argument $\rightarrow$ Exemple. Une suite d'affirmations gratuites vaut 0 point en argumentation.
\item \textbf{Rédiger la conclusion :} \textbf{Répondre clairement à la question problématique.} Reprendre le fil du plan. Ouvrir \emph{légèrement} (perspective) sans nouveau développement. « On peut donc conclure que... »
\end{enumerate}
\textbf{Formules de liaison à maîtriser :} « En premier lieu... », « Certes... mais... », « Par conséquent... », « Cela montre que... », « En définitive... », « Force est de constater que... ».
\end{methode}

\begin{exoresolu}[Question de réflexion : L'importance de la philosophie (Sujet classique Probatoire)]
\textbf{Énoncé.} « La philosophie ne fait pas le pain. » Montrez, en un développement structuré d'une vingtaine de lignes, que la philosophie est pourtant utile à l'homme et à la société.
\tcblower
\textbf{Solution détaillée.}

\emph{Introduction.} L'adage populaire « la philosophie ne fait pas le pain » oppose la réflexion philosophique, jugée abstraite et stérile, aux activités concrètes et productives (agriculture, artisanat, industrie) qui assurent la survie matérielle. Faut-il en conclure que la philosophie est un luxe inutile ? Nous montrerons que la philosophie est utile d'abord à l'individu pour sa lucidité et sa liberté, ensuite à la société pour sa justice et sa cohésion.

\emph{1. Utilité pour l'individu : former l'esprit critique et la liberté.}
La philosophie est l'école de l'\cle{esprit critique}. Elle apprend à ne pas croire sur parole, à distinguer le \cle{vrai} du \cle{faux}, le \cle{savoir} de l'\cle{opinion} (\emph{doxa}). L'homme qui philosophe échappe aux rumeurs (ex: fausses informations sur les réseaux sociaux au Cameroun), aux manipulations politiques et aux préjugés ethniques ou sexistes. Ainsi, le citoyen qui examine les programmes et les bilans avant de voter, au lieu de voter par fidélité tribale ou par cadeau (corruption électorale), agit en philosophe : il use de sa \cle{raison} pour choisir. De plus, la philosophie aide à \cle{vivre} : en nous interrogeant sur le bonheur, la mort, le sens de l'existence (Stoïciens, Épicuriens, Existentialistes), elle nous rend \cle{libres} et \cle{responsables} de nos choix. Selon \textbf{Socrate} (\emph{Apologie de Socrate}), « une vie sans examen ne vaut pas la peine d'être vécue ».

\emph{2. Utilité pour la société : fonder le droit, la justice et la démocratie.}
La philosophie est indispensable à la vie collective. Elle est la mère du \cle{Droit} et de la \cle{Justice} : réfléchir sur ce qui est \emph{juste} (Platon, \emph{La République} ; Rawls, \emph{Théorie de la justice}) permet de construire des lois meilleures, plus équitables, protégeant le faible contre le fort. Elle est la gardienne de la \cle{démocratie} : un peuple capable de discuter rationnellement, d'accepter la contradiction et le compromis (Habermas, \emph{éthique de la discussion}) résiste à la violence, à la démagogie et au totalitarisme. Au Cameroun, le dialogue national ou les concertations politiques gagnent à être menés avec une exigence philosophique d'écoute et d'argumentation. Enfin, elle éclaire le \cle{progrès scientifique et technique} : face aux machines, à l'IA, aux biotechnologies, elle pose la question morale cruciale : « Tout ce qui est \emph{techniquement possible} est-il \emph{moralement souhaitable} ? » (Jonas, \emph{Principe responsabilité}).

\emph{Conclusion.} Il est vrai que la philosophie ne produit ni pain, ni routes, ni médicaments directement. Mais elle produit quelque chose de plus fondamental : des \cle{esprits libres}, \cle{critiques} et \cle{responsables}, sans lesquels aucune société humaine ne peut véritablement prospérer, ni même survivre à ses propres excès techniques. L'utilité de la philosophie est donc d'un autre ordre que celle du pain : elle \emph{nourrit l'esprit}, condition de la dignité humaine.
\end{exoresolu}

\courssub{TP Guidé : Philosopher sur des situations camerounaises (Mise en situation - Savoir-faire)}

\begin{exercice}[Situation 1 : La tradition et le droit moderne]
\textbf{Contexte.} Dans une chefferie du Nord-Ouest, un notable affirme : « La coutume veut que la femme n'hérite pas de la terre. C'est notre philosophie, on ne change pas. » Une jeune fille, étudiante en droit, répond : « Cette coutume est injuste, elle viole l'égalité devant la loi. »
\begin{enumerate}
\item Identifiez dans les deux propos ce qui relève de l'\emph{opinion} (\emph{doxa}) et ce qui relève de la \emph{démarche philosophique}.
\item Quelle notion philosophique (égalité, justice, droit naturel, droit positif, tradition) permet de dépasser l'opposition ?
\item Rédigez un paragraphe argumenté (10 lignes) pour trancher : la coutume a-t-elle valeur de loi morale absolue ?
\end{enumerate}
\end{exercice}

\begin{exercice}[Situation 2 : Science et éthique médicale]
\textbf{Contexte.} Un hôpital de Douala propose une sélection embryonnaire (DPI) pour éviter une maladie génétique grave. Un parent refuse au nom de « la volonté de Dieu ». Un médecin répond : « La science nous permet d'éviter la souffrance, il faut l'utiliser. »
\begin{enumerate}
\item Distinguez les trois registres de discours présents : \emph{religieux/théologique}, \emph{scientifique/technique}, \emph{philosophique/éthique}.
\item Pourquoi la science seule ne peut pas trancher le dilemme : « Faut-il sélectionner les embryons ? »
\item Quelle est la valeur ajoutée de la philosophie (bioéthique) dans ce débat ?
\end{enumerate}
\end{exercice}

\begin{exercice}[Situation 3 : L'étonnement philosophique au quotidien]
\textbf{Contexte.} Assis dans un taxi « Clap » à Yaoundé, vous observez le chauffeur qui refuse un billet de 1000 FCFA déchiré mais accepte un billet neuf de 500 FCFA. Vous vous demandez : « Qu'est-ce qui donne de la valeur à l'argent ? La confiance ? L'État ? Le papier ? »
\begin{enumerate}
\item Expliquez pourquoi cette question (« Qu'est-ce que la valeur de la monnaie ? ») est une question \emph{philosophique} (origine, objet, méthode).
\item Proposez deux thèses opposées (Thèse / Antithèse) pour y répondre.
\item Indiquez quel auteur (étudié en Première ou culture générale) pourrait éclairer ce débat (ex: Locke, Marx, Simmel, Aristote).
\end{enumerate}
\end{exercice}

\begin{attention}[Les 5 erreurs fatales qui coûtent des points au Probatoire / Bac Tech]
\begin{enumerate}
\item \textbf{Confondre philosophie et sagesse.} La philosophie n'est pas la sagesse elle-même (contenu, état abouti), mais l'\emph{amour} / la \emph{recherche} / le \emph{désir} de la sagesse (étymologie : \emph{philein} + \emph{sophia}). Écrire « La philosophie est la sagesse » est une \textbf{erreur de définition} sanctionnée.
\item \textbf{Confondre mythe, religion et philosophie.} Une croyance, un dogme, un récit traditionnel ou une explication par les esprits/ancêtres n'est pas une philosophie. La philosophie exige l'\cle{argumentation rationnelle}, l'\cle{examen critique} et l'\cle{universalité} (passage du \emph{mythos} au \emph{logos}). Ne dites jamais : « La philosophie africaine, c'est les mythes » ou « L'explication par les ancêtres est philosophique ».
\item \textbf{Affirmer sans justifier (Opinion gratuite).} Toute thèse doit être \emph{soutenue par un argument logique} et \emph{illustrée par un exemple concret} (idéalement camerounais/africain). Une suite d'affirmations du type « La philosophie est importante car elle est importante » ou « Elle aide à vivre car elle aide à vivre » vaut \textbf{zéro point} en argumentation.
\item \textbf{Réciter le cours sans répondre à la question (Hors-sujet).} Il faut d'abord \emph{analyser le sujet} (définir, problématiser) et construire une réponse \emph{adaptée}. Un devoir parfait sur « La méthode de Descartes » alors que le sujet est « L'utilité de la philosophie » = \textbf{Note très faible (souvent < 4/20)}.
\item \textbf{Opposer totalement philosophie et science (Scientisme ou Irrationalisme).} Erreur classique : dire « La science est vraie, la philosophie est fausse / inutile » (Scientisme) ou « La science détruit le mystère, la philosophie le préserve » (Irrationalisme). Il faut montrer leurs \emph{différences de méthode et d'objet} \textbf{ET} leur \emph{complémentarité} (Philosophie = mère des sciences + conscience critique de la technoscience).
\end{enumerate}
\end{attention}

\begin{savaistu}[Histoire des idées \& Contexte camerounais]
\textbf{Les prémices de la philosophie en Afrique.} Bien que le terme « philosophie » soit grec, l'\cle{activité philosophique} (questionnement rationnel sur l'homme, le monde, Dieu, la justice) est universelle.
\begin{itemize}
\item \textbf{Égypte antique (Kémèt) :} Ptahhotep (V\textsuperscript{e} dynastie, ~2400 av. J.-C.), \emph{Maximes} : sagesse, justice (\emph{Maât}), écoute, mesure. Antécédent direct de la philosophie grecque (Platon a séjourné en Égypte).
\item \textbf{Philosophie africaine moderne :} Débat sur l'existence d'une « philosophie africaine » (Placide Tempels, \emph{La Philosophie bantoue}, 1945 ; Paulin Hountondji, \emph{Sur la philosophie africaine}, 1973 ; Kwasi Wiredu, \emph{Philosophie et culture africaine}).
\item \textbf{Cameroun :} \textbf{Fabien Eboussi Boulaga} (1934-2018), jésuite, philosophe de la libération, \emph{Christianisme sans fétiche} (1981) : critique de l'aliénation culturelle, appel à une « anthropologie philosophique » africaine. \textbf{Marcien Towa} (1931-2014), \emph{Essai sur la problématique philosophique dans l'Afrique actuelle} (1971) : exige une philosophie critique, rationnelle, décolonisée.
\item \textbf{Programme MINESEC :} L'étude de la philosophie en Première Technique vise à former des \emph{esprits critiques} capables de penser les défis du Cameroun (développement, démocratie, éthique biomédicale, identité culturelle, vivre-ensemble) avec les outils de la raison universelle.
\end{itemize}
\end{savaistu}

\newpage

\coursec{Exercices}

\courssub{Je vérifie mes connaissances}

\begin{exercice}[QCM : les fondamentaux]
Pour chaque question, une seule réponse est correcte. Recopie la lettre correspondante.
\begin{enumerate}
\item Le mot « philosophie » vient du grec \emph{philein} et \emph{sophia} et signifie :
\begin{enumerate}
\item la science de la sagesse ;
\item l'amour de la sagesse ;
\item la possession de la vérité.
\end{enumerate}
\item La philosophie est née, selon la tradition, en Grèce :
\begin{enumerate}
\item au \textsc{ii}\textsuperscript{e} siècle avant J.-C. ;
\item au \textsc{vi}\textsuperscript{e} siècle avant J.-C. ;
\item au \textsc{xv}\textsuperscript{e} siècle après J.-C.
\end{enumerate}
\item La maïeutique, méthode d'accouchement des esprits par le questionnement, est attribuée à :
\begin{enumerate}
\item Aristote ;
\item Platon ;
\item Socrate.
\end{enumerate}
\item Selon Platon, le commencement de la philosophie est :
\begin{enumerate}
\item l'étonnement ;
\item la peur ;
\item le calcul.
\end{enumerate}
\end{enumerate}
\end{exercice}

\begin{exercice}[Vrai ou faux justifié]
Pour chaque affirmation, indique si elle est vraie ou fausse, puis justifie ta réponse en une ou deux phrases.
\begin{enumerate}
\item La philosophie et la science poursuivent exactement le même objet.
\item La méthode de la philosophie repose sur le doute méthodique, l'analyse et l'argumentation.
\item La philosophie n'a aucune utilité pratique dans la vie quotidienne.
\item Toute opinion personnelle est une vérité philosophique.
\end{enumerate}
\end{exercice}

\begin{exercice}[Définitions]
Définis en deux ou trois phrases chacun des termes suivants, tels qu'ils ont été étudiés dans la leçon :
\begin{enumerate}
\item la philosophie ;
\item l'étonnement ;
\item la maïeutique ;
\item le doute méthodique.
\end{enumerate}
\end{exercice}

\begin{exercice}[Complétion]
Recopie et complète le texte suivant avec les mots de la liste : \emph{Grèce} ; \emph{sagesse} ; \emph{questionnement} ; \emph{science} ; \emph{universelle}.

« La philosophie est née en \dots\dots, au \textsc{vi}\textsuperscript{e} siècle avant J.-C. Étymologiquement, elle désigne l'amour de la \dots\dots. Sa méthode repose sur le \dots\dots rationnel et l'argumentation. Contrairement à la \dots\dots, qui étudie des objets particuliers de manière expérimentale, la philosophie pose des questions \dots\dots sur l'existence, la vérité et les valeurs. »
\end{exercice}

\courssub{Je m'entraîne}

\begin{exercice}[Classer et justifier]
Classe chacune des six affirmations suivantes selon qu'elle relève de la \cle{science}, de la \cle{philosophie} ou de l'\cle{opinion}, et justifie ton choix en une phrase.
\begin{enumerate}
\item « $2+2=4$. »
\item « Il faut être juste envers les autres. »
\item « L'eau bout à \SI{100}{\celsius} au niveau de la mer. »
\item « La vie a-t-elle un sens ? »
\item « Le ndolé est le meilleur plat du Cameroun. »
\item « Qu'est-ce qu'une société juste ? »
\end{enumerate}
\end{exercice}

\begin{exercice}[Philosopher au quotidien]
Pour chacune des situations camerounaises suivantes, formule une question philosophique pertinente (sur la justice, la vérité, la liberté ou le bonheur), puis explique en deux phrases pourquoi cette question relève de la philosophie et non de la science.
\begin{enumerate}
\item Un élève de Première technique découvre qu'un camarade a triché à une évaluation.
\item Un jeune technicien hésite entre un emploi bien payé à l'étranger et un emploi modeste au service de son village.
\item Un chef de quartier doit répartir équitablement l'eau d'un forage entre les habitants pendant la saison sèche.
\end{enumerate}
\end{exercice}

\begin{exercice}[Analyse critique d'un message viral]
Sur un groupe WhatsApp de quartier à Douala circule le message suivant : « Un célèbre médecin a révélé que boire de l'eau de mer guérit toutes les maladies. Partagez vite à tous vos contacts ! »

Rédige un paragraphe d'une dizaine de lignes dans lequel tu appliques la méthode du \cle{doute méthodique} à ce message, en examinant successivement : la \cle{source} (qui parle ? est-elle identifiable ?), la \cle{preuve} (quels faits vérifiables sont donnés ?), la \cle{cohérence} (le message est-il logique et compatible avec les connaissances établies ?) et la \cle{conclusion} (que faut-il raisonnablement en penser et faire ?).
\end{exercice}

\begin{exercice}[L'étonnement en action]
\begin{enumerate}
\item Rappelle la thèse de Platon : « l'étonnement est le commencement de la philosophie ».
\item À partir de deux phénomènes de ta vie quotidienne (par exemple : le lever du soleil, le fonctionnement d'un téléphone portable, la mort d'un proche), montre comment l'étonnement peut conduire à une question philosophique.
\item Distingue, pour chacun de tes exemples, la question scientifique (qui cherche une explication expérimentale) de la question philosophique (qui interroge le sens).
\end{enumerate}
\end{exercice}

\courssub{Je me prépare au Bac}

\begin{exercice}[Dissertation courte — type Probatoire]
\textbf{Sujet :} La philosophie est-elle utile au technicien ?

Rédige une dissertation courte en suivant la démarche ci-dessous.
\begin{enumerate}
\item \textbf{Introduction :} définis la philosophie à partir de son étymologie et de son objet ; dégage la problématique (un technicien, formé à la précision et à l'efficacité, a-t-il besoin de réfléchir sur le sens, les valeurs et les finalités de son activité ?) ; annonce ton plan.
\item \textbf{Développement :} présente deux arguments étayés chacun par un exemple camerounais concret. Par exemple : la philosophie forme l'esprit critique face aux rumeurs et aux fausses informations ; elle éclaire les choix éthiques du technicien (sécurité des constructions, honnêteté dans les devis, respect de l'environnement).
\item \textbf{Conclusion :} réponds clairement à la question posée et ouvre sur une question nouvelle.
\end{enumerate}
\end{exercice}

\begin{exercice}[Commentaire de texte guidé — Platon]
Lis attentivement l'extrait suivant, inspiré du \emph{Théétète} de Platon :

« L'étonnement est le sentiment d'un philosophe, et la philosophie commence par l'étonnement. Celui qui s'étonne et qui doute se reconnaît ignorant ; c'est pourquoi il se met à chercher. »

Réponds aux questions suivantes.
\begin{enumerate}
\item Identifie la thèse de l'auteur et reformule-la en une phrase avec tes propres mots.
\item Explique le sens de l'affirmation : « l'étonnement est le commencement de la philosophie ». En quoi l'étonnement suppose-t-il la conscience de l'ignorance ?
\item Illustre cette thèse par un exemple tiré de la vie courante ou de ta formation technique.
\item Discute : l'étonnement suffit-il à faire un philosophe, ou faut-il y ajouter une méthode ? Justifie ta réponse.
\end{enumerate}
\end{exercice}

\begin{exercice}[Débat argumenté — type Baccalauréat]
\textbf{Sujet :} Peut-on se passer de la philosophie à l'ère de la science et de la technologie ?

Construis une réponse dialectique en suivant le plan guidé ci-dessous.
\begin{enumerate}
\item \textbf{Thèse :} présente deux raisons pour lesquelles on pourrait penser que la science et la technologie rendent la philosophie inutile (la science donne des résultats vérifiables ; la technologie résout les problèmes concrets). Appuie chaque raison sur un exemple (médecine, téléphonie mobile, agriculture au Cameroun).
\item \textbf{Antithèse :} montre, par deux arguments, que la science ne répond pas à toutes les questions : elle dit \emph{comment} agir, mais non \emph{pourquoi} ni \emph{s'il faut} agir (exemples : usage des données personnelles, intelligence artificielle, OGM). Rappelle à cet égard la différence entre l'objet de la science et celui de la philosophie.
\item \textbf{Synthèse :} propose une position personnelle nuancée : science et philosophie sont-elles rivales ou complémentaires ? Justifie ta réponse et conclus par une ouverture.
\end{enumerate}
\end{exercice}

\newpage

\coursec{Corrigés des exercices}

\begin{corrige}[Exercice 1 — QCM : les fondamentaux]
\begin{enumerate}
\item \textbf{Réponse b : l'amour de la sagesse.} Le mot « philosophie » vient du grec \emph{philein} (aimer) et \emph{sophia} (sagesse, savoir). Le philosophe n'est donc pas celui qui possède la sagesse, mais celui qui la recherche : c'est ce qui distingue le \emph{philosophos} du \emph{sophos} (le sage).
\item \textbf{Réponse b : au VI\textsuperscript{e} siècle avant J.-C.} Selon la tradition, la philosophie est née en Grèce (plus précisément en Ionie, à Milet, avec Thalès) au VI\textsuperscript{e} siècle avant J.-C.
\item \textbf{Réponse c : Socrate.} La maïeutique (du grec \emph{maieutikè}, art d'accoucher) est la méthode de Socrate : par le questionnement, il aide son interlocuteur à « accoucher » de la vérité qu'il porte en lui, comme sa mère, sage-femme, aidait les femmes à accoucher.
\item \textbf{Réponse a : l'étonnement.} Pour Platon (dans le \emph{Théétète}), l'étonnement (\emph{thaumazein}) est le commencement de la philosophie : c'est parce que l'homme s'étonne devant le monde qu'il se met à interroger et à chercher.
\end{enumerate}
\textbf{Point de vigilance :} ne pas confondre « amour de la sagesse » (philosophie) et « possession de la sagesse » : le philosophe est un chercheur, non un détenteur de la vérité.
\end{corrige}

\begin{corrige}[Exercice 2 — Vrai ou faux justifié]
\begin{enumerate}
\item \textbf{Faux.} La science étudie des objets particuliers (la matière, le vivant, les nombres) par l'expérience et le calcul, tandis que la philosophie pose des questions universelles sur le sens de l'existence, la vérité et les valeurs. Leurs objets et leurs méthodes diffèrent, même si elles recherchent toutes deux la vérité.
\item \textbf{Vrai.} La méthode philosophique repose sur le doute méthodique (ne rien admettre sans l'avoir examiné), l'analyse des notions et l'argumentation rationnelle : philosopher, c'est douter, examiner, puis justifier par des raisons.
\item \textbf{Faux.} La philosophie a une utilité pratique : elle forme l'esprit critique face aux rumeurs et aux fausses informations, éclaire les choix moraux (honnêteté, justice) et aide chacun à donner un sens à sa vie et à son travail.
\item \textbf{Faux.} Une opinion (\emph{doxa}) est un jugement personnel, non justifié et variable selon les individus ; une vérité philosophique doit être argumentée, justifiée par des raisons et soumise à l'examen critique. Toute opinion n'est donc pas une vérité.
\end{enumerate}
\textbf{Point de vigilance :} une justification d'une ou deux phrases suffit, mais elle doit mobiliser une notion précise du cours (objet, méthode, distinction opinion/vérité).
\end{corrige}

\begin{corrige}[Exercice 3 — Définitions]
\begin{enumerate}
\item \textbf{La philosophie} : étymologiquement, « amour de la sagesse » (du grec \emph{philein}, aimer, et \emph{sophia}, sagesse). C'est une réflexion rationnelle et critique qui interroge les questions fondamentales de l'existence humaine : la vérité, la justice, la liberté, le bonheur, le sens de la vie. Elle procède par le doute, l'analyse et l'argumentation.
\item \textbf{L'étonnement} : sentiment de surprise et d'admiration devant ce qui paraît étrange ou incompréhensible. Selon Platon, il est le commencement de la philosophie, car celui qui s'étonne prend conscience de son ignorance et se met à chercher.
\item \textbf{La maïeutique} : méthode de Socrate, littéralement « art d'accoucher les esprits ». Par un questionnement habile, le philosophe aide son interlocuteur à découvrir par lui-même la vérité qu'il porte en lui, et à abandonner ses fausses certitudes.
\item \textbf{Le doute méthodique} : démarche qui consiste à suspendre provisoirement son jugement et à ne rien tenir pour vrai avant de l'avoir examiné rationnellement. Ce n'est pas un doute sceptique (qui doute pour douter), mais un instrument de recherche de la vérité.
\end{enumerate}
\textbf{Point de vigilance :} une définition correcte comporte le genre proche (sentiment, méthode, démarche) et la différence spécifique ; elle ne doit pas être une simple paraphrase du mot.
\end{corrige}

\begin{corrige}[Exercice 4 — Complétion]
« La philosophie est née en \textbf{Grèce}, au VI\textsuperscript{e} siècle avant J.-C. Étymologiquement, elle désigne l'amour de la \textbf{sagesse}. Sa méthode repose sur le \textbf{questionnement} rationnel et l'argumentation. Contrairement à la \textbf{science}, qui étudie des objets particuliers de manière expérimentale, la philosophie pose des questions \textbf{universelles} sur l'existence, la vérité et les valeurs. »

\textbf{Justifications :} « Grèce » est le berceau historique de la philosophie ; « sagesse » traduit le grec \emph{sophia} ; « questionnement » désigne la démarche interrogative propre au philosophe ; « science » est le terme à opposer à la philosophie dans la dernière phrase ; « universelles » qualifie les questions philosophiques, qui concernent tout homme et non un objet particulier.
\end{corrige}

\begin{corrige}[Exercice 5 — Classer et justifier]
\begin{enumerate}
\item « $2+2=4$ » : \textbf{science} (mathématiques). C'est une vérité démontrable, nécessaire et universelle, établie par le raisonnement.
\item « Il faut être juste envers les autres » : \textbf{philosophie}. C'est une réflexion sur une valeur morale (la justice) et sur ce qu'il faut faire, qui appelle une justification rationnelle.
\item « L'eau bout à \SI{100}{\celsius} au niveau de la mer » : \textbf{science} (physique). C'est un fait vérifiable par l'expérience et la mesure.
\item « La vie a-t-elle un sens ? » : \textbf{philosophie}. C'est une question universelle sur le sens de l'existence, à laquelle aucune expérience ne peut répondre seule.
\item « Le ndolé est le meilleur plat du Cameroun » : \textbf{opinion}. C'est un jugement de goût personnel, variable selon les individus, qui ne peut être démontré.
\item « Qu'est-ce qu'une société juste ? » : \textbf{philosophie}. C'est une interrogation rationnelle sur la définition et les conditions de la justice sociale.
\end{enumerate}
\textbf{Point de vigilance :} le critère de classement doit être explicité : vérité démontrable ou expérimentale (science), question de sens ou de valeur argumentée (philosophie), jugement personnel non justifié (opinion).
\end{corrige}

\begin{corrige}[Exercice 6 — Philosopher au quotidien]
\begin{enumerate}
\item \textbf{Question philosophique :} « Est-il juste de dénoncer un camarade qui a triché ? » ou « Pourquoi la tricherie est-elle une faute ? ». Cette question porte sur la justice et l'honnêteté, c'est-à-dire sur des valeurs morales : aucune expérience scientifique ne peut dire ce qui est juste ou injuste. Elle relève donc de la réflexion philosophique, qui examine les principes de l'action.
\item \textbf{Question philosophique :} « Doit-on privilégier son bonheur personnel ou le devoir envers sa communauté ? ». Ce dilemme oppose deux conceptions du bien (bonheur individuel, devoir moral) : il demande de définir ce qu'est une vie réussie, ce qui est une question de sens et de valeurs, non un problème technique mesurable.
\item \textbf{Question philosophique :} « Qu'est-ce qu'un partage équitable : donner la même part à chacun, ou donner selon les besoins de chacun ? ». Cette question interroge le concept même de justice (égalité stricte ou équité) ; la science peut mesurer les volumes d'eau, mais non décider du critère juste de répartition.
\end{enumerate}
\textbf{Point de vigilance :} la question formulée doit être une vraie question (avec point d'interrogation) portant sur une valeur (justice, liberté, vérité, bonheur), et la justification doit opposer clairement la question de \emph{sens} (philosophie) à la question de \emph{fait} (science).
\end{corrige}

\begin{corrige}[Exercice 7 — Analyse critique d'un message viral]
\textbf{Corrigé rédigé (paragraphe attendu) :}

Appliquons le doute méthodique à ce message. D'abord, examinons la \cle{source} : le message évoque « un célèbre médecin » sans le nommer, sans indiquer ni son pays, ni son hôpital, ni la publication où il aurait fait cette « révélation ». Une source non identifiable n'offre aucune garantie : n'importe qui peut inventer une telle affirmation. Ensuite, la \cle{preuve} : le message ne donne aucun fait vérifiable, aucune expérience, aucune statistique, aucun résultat d'étude ; il se contente d'affirmer. Or, en science comme en philosophie, une affirmation sans preuve ne vaut rien. Quant à la \cle{cohérence}, l'affirmation est incompatible avec les connaissances établies : l'eau de mer, très salée, déshydrate l'organisme et peut aggraver l'état d'un malade ; prétendre qu'elle guérit « toutes » les maladies est absurde, car les maladies ont des causes très différentes et ne sauraient avoir un remède unique. Enfin, la \cle{conclusion} s'impose : il faut raisonnablement rejeter ce message comme une fausse information, ne pas le partager, et avertir ceux qui l'ont reçu. Ce travail de doute méthodique montre l'utilité pratique de la philosophie : elle nous protège des rumeurs et des manipulations qui circulent sur les réseaux sociaux.

\textbf{Barème indicatif (sur 10) :} source examinée (2 pts) ; preuve examinée (2 pts) ; cohérence examinée (2 pts) ; conclusion raisonnable (2 pts) ; clarté et organisation du paragraphe (2 pts).

\textbf{Points de vigilance :} les quatre étapes (source, preuve, cohérence, conclusion) doivent apparaître explicitement ; il ne suffit pas de dire « c'est faux », il faut montrer \emph{pourquoi} par un raisonnement ordonné.
\end{corrige}

\begin{corrige}[Exercice 8 — L'étonnement en action]
\begin{enumerate}
\item \textbf{Thèse de Platon :} dans le \emph{Théétète}, Platon affirme que « l'étonnement est le sentiment d'un philosophe, et la philosophie commence par l'étonnement ». Autrement dit, philosopher naît de la surprise devant le monde : celui qui s'étonne reconnaît son ignorance et se met à chercher la vérité.
\item \textbf{Exemples d'étonnement conduisant à philosopher :}
\begin{itemize}
\item \emph{Le lever du soleil :} chaque matin, à Yaoundé comme ailleurs, le soleil se lève. Au lieu de le considérer comme banal, je m'étonne : pourquoi y a-t-il un ordre régulier dans la nature ? Cet étonnement peut conduire à la question philosophique : « Le monde a-t-il un sens, une finalité ? »
\item \emph{La mort d'un proche :} la disparition d'un être cher provoque un étonnement douloureux : comment une personne pleine de vie peut-elle cesser d'exister ? Cet étonnement conduit à la question philosophique : « Qu'est-ce que la mort, et que vaut une vie humaine ? »
\end{itemize}
\item \textbf{Distinction science / philosophie :}
\begin{itemize}
\item Pour le soleil : la question \emph{scientifique} est « Quels mécanismes physiques expliquent la rotation de la Terre et le lever du soleil ? » (explication expérimentale) ; la question \emph{philosophique} est « Pourquoi y a-t-il quelque chose plutôt que rien, et cet ordre a-t-il un sens ? » (interrogation sur le sens).
\item Pour la mort : la question \emph{scientifique} est « Quels processus biologiques provoquent l'arrêt des fonctions vitales ? » ; la question \emph{philosophique} est « La mort donne-t-elle son prix à la vie ? Comment bien vivre en sachant que l'on va mourir ? »
\end{itemize}
\end{enumerate}
\textbf{Point de vigilance :} la question scientifique cherche une cause ou un mécanisme vérifiable (le \emph{comment}) ; la question philosophique interroge le sens, la valeur ou la finalité (le \emph{pourquoi}). Les deux exemples doivent être traités selon cette grille.
\end{corrige}

\begin{corrige}[Exercice 9 — Dissertation courte : La philosophie est-elle utile au technicien ?]
\textbf{Corrigé rédigé (modèle de copie) :}

\emph{Introduction.} La philosophie, étymologiquement « amour de la sagesse » (du grec \emph{philein} et \emph{sophia}), est une réflexion rationnelle et critique sur les questions fondamentales de l'existence : la vérité, la justice, le devoir. Le technicien, lui, est formé à la précision, au calcul et à l'efficacité : il sait \emph{comment} faire. On peut alors se demander si la philosophie lui est utile : un technicien, dont le métier est de résoudre des problèmes concrets, a-t-il besoin de réfléchir sur le sens, les valeurs et les finalités de son activité ? Nous montrerons que la philosophie est utile au technicien, d'abord parce qu'elle forme son esprit critique, ensuite parce qu'elle éclaire ses choix éthiques.

\emph{Développement.} Premièrement, la philosophie forme l'esprit critique, indispensable au technicien. Par le doute méthodique, elle apprend à ne rien admettre sans preuve et à vérifier les sources. Or le technicien camerounais est entouré de fausses informations : rumeurs sur les matériaux, « astuces » de construction non vérifiées, messages viraux sur WhatsApp. Un électricien qui croirait une rumeur sur une installation « miracle » mettrait des vies en danger. Grâce à la philosophie, il examine, doute, vérifie avant d'agir : sa rigueur technique s'en trouve renforcée.

Deuxièmement, la philosophie éclaire les choix éthiques du technicien. La technique dit comment construire, mais non s'il faut tricher sur les matériaux ou respecter les normes. Prenons l'exemple d'un chef de chantier à Douala tenté d'utiliser moins de ciment que prévu pour augmenter son bénéfice : seule une réflexion morale sur l'honnêteté et la responsabilité peut le guider, car la sécurité des futurs habitants est en jeu. De même, le technicien agricole qui choisit des méthodes respectueuses de l'environnement accomplit un choix de valeurs, pas seulement un calcul de rendement. La philosophie donne ainsi un sens moral au savoir-faire.

\emph{Conclusion.} La philosophie est donc bien utile au technicien : elle aiguise son esprit critique face aux rumeurs et oriente ses décisions vers la justice et la responsabilité. Le bon technicien n'est pas seulement celui qui sait faire, mais celui qui sait pourquoi et pour qui il fait. Reste alors une question : comment former davantage les techniciens camerounais à cette réflexion critique ?

\textbf{Barème indicatif (sur 20) :} introduction (définitions, problématique, annonce du plan) : 5 pts ; premier argument étayé par un exemple camerounais : 5 pts ; deuxième argument étayé par un exemple camerounais : 5 pts ; conclusion (réponse claire et ouverture) : 3 pts ; langue et organisation : 2 pts.

\textbf{Points de vigilance :} la problématique doit être formulée sous forme de question ; chaque argument doit être suivi d'un exemple concret ; la conclusion doit répondre explicitement à la question posée (« oui, elle est utile ») et s'achever par une ouverture.
\end{corrige}

\begin{corrige}[Exercice 10 — Commentaire de texte guidé : Platon]
\begin{enumerate}
\item \textbf{Thèse de l'auteur :} Platon affirme que la philosophie naît de l'étonnement. Reformulée : c'est parce que l'homme s'étonne devant ce qu'il ne comprend pas, et qu'il prend ainsi conscience de son ignorance, qu'il se met à chercher la vérité.
\item \textbf{Explication :} dire que « l'étonnement est le commencement de la philosophie » signifie que la réflexion philosophique ne part pas de certitudes, mais d'une surprise : devant le monde, la mort, la justice, l'homme s'aperçoit qu'il ne sait pas. L'étonnement suppose donc la conscience de l'ignorance : celui qui croit tout savoir ne s'étonne de rien et ne cherche rien ; celui qui s'étonne, au contraire, « se reconnaît ignorant » et « se met à chercher ». L'étonnement est ainsi le moteur de la recherche de la vérité.
\item \textbf{Illustration :} un élève de Première technique découvre qu'un téléphone portable transmet instantanément la voix à des centaines de kilomètres. S'il s'étonne (« comment est-ce possible ? »), il peut être conduit à une question plus profonde : « cette technique qui relie les hommes les rapproche-t-elle vraiment ? ». Son étonnement, en révélant son ignorance, l'a fait passer de la simple utilisation de l'objet à la réflexion.
\item \textbf{Discussion :} l'étonnement ne suffit pas à faire un philosophe : il n'en est que le « commencement », le déclencheur. Beaucoup s'étonnent sans jamais chercher sérieusement. Il faut y ajouter une méthode : le doute méthodique, l'analyse des notions et l'argumentation, comme le montre la maïeutique de Socrate qui transforme la surprise en recherche rationnelle. L'étonnement fournit la question ; la méthode permet d'y répondre de manière justifiée.
\end{enumerate}
\textbf{Barème indicatif (sur 10) :} thèse identifiée et reformulée (2 pts) ; explication du lien étonnement/ignorance (3 pts) ; exemple pertinent (2 pts) ; discussion argumentée (3 pts).

\textbf{Points de vigilance :} la reformulation doit être personnelle (pas de simple recopie) ; la discussion doit trancher avec une justification, non rester indécise.
\end{corrige}

\begin{corrige}[Exercice 11 — Débat argumenté : Peut-on se passer de la philosophie à l'ère de la science et de la technologie ?]
\textbf{Corrigé rédigé (plan dialectique complété) :}

\emph{Thèse : on pourrait croire la philosophie devenue inutile.} Première raison : la science donne des résultats vérifiables et certains, là où la philosophie semble débattre sans jamais conclure. La médecine, par exemple, guérit des maladies grâce à des traitements expérimentalement testés, et non grâce à des discussions abstraites. Deuxième raison : la technologie résout directement les problèmes concrets de la vie quotidienne. Au Cameroun, le téléphone mobile permet de communiquer et de faire des transactions (Mobile Money), et les techniques agricoles améliorées augmentent les récoltes : à quoi bon alors philosopher ?

\emph{Antithèse : la science ne répond pas à toutes les questions.} Premier argument : la science et la technologie disent \emph{comment} agir, mais non \emph{pourquoi} ni \emph{s'il faut} agir. La technique de collecte des données personnelles existe, mais seule une réflexion philosophique peut dire s'il est juste de l'utiliser sans le consentement des personnes ; de même, l'intelligence artificielle peut prendre des décisions, mais c'est à la philosophie de demander si on doit les lui confier. Deuxième argument : l'objet de la science diffère de celui de la philosophie. La science étudie des objets particuliers de manière expérimentale ; la philosophie pose des questions universelles sur le sens, les valeurs et les finalités. Les OGM, par exemple, relèvent de la science pour leur fabrication, mais de la philosophie pour la question : est-il bon de transformer ainsi le vivant ? On ne peut donc pas se passer de la philosophie, car les questions les plus importantes pour l'homme échappent à l'expérience.

\emph{Synthèse : science et philosophie sont complémentaires, non rivales.} La science accroît notre pouvoir d'agir ; la philosophie éclaire l'usage de ce pouvoir. Sans science, la philosophie serait une rêverie sans prise sur le réel ; sans philosophie, la science serait un pouvoir aveugle, capable du meilleur comme du pire. Le progrès véritable, pour le Cameroun comme pour le monde, suppose donc les deux : des techniciens compétents \emph{et} des citoyens capables de réfléchir au sens de leurs actions. Reste une question ouverte : comment faire dialoguer davantage scientifiques et philosophes dans nos sociétés ?

\textbf{Barème indicatif (sur 20) :} thèse avec deux arguments et exemples (5 pts) ; antithèse avec deux arguments et distinction science/philosophie (6 pts) ; synthèse personnelle et nuancée (5 pts) ; conclusion avec ouverture (2 pts) ; langue et cohérence (2 pts).

\textbf{Points de vigilance :} le plan dialectique impose que les trois moments (thèse, antithèse, synthèse) soient clairement distingués ; la distinction entre le \emph{comment} (science) et le \emph{pourquoi / s'il faut} (philosophie) doit être explicitement rappelée ; la synthèse doit être une position personnelle justifiée, non une simple répétition de l'antithèse.
\end{corrige}

\newpage

\coursec{Fiche bilan}

\begin{popfigure}
\begin{tikzpicture}[mindmap, grow cyclic, every node/.style={concept, text=white, font=\small\bfseries},
root concept/.append style={concept color=popink, font=\large\bfseries},
level 1/.append style={level distance=3.4cm, sibling angle=60},
level 1 concept/.append style={font=\small\bfseries}]
\node[root concept]{La Philosophie}
child[concept color=popblue]{node[align=center]{Définitions\\Amour de la sagesse\\Recherche de la vérité}}
child[concept color=poporange]{node[align=center]{Origine\\Grèce, VI\textsuperscript{e} s. av. J.-C.\\Étonnement, doute}}
child[concept color=popgreen]{node[align=center]{Objet\\L'homme, le monde,\\Dieu, la morale}}
child[concept color=poppurple]{node[align=center]{Méthode\\Doute, analyse,\\argumentation, critique}}
child[concept color=poppink]{node[align=center]{Importance\\Esprit critique,\\liberté, tolérance}}
child[concept color=popgold]{node[align=center]{Philosophie\\et Science\\Parenté et différences}};
\end{tikzpicture}
\legende{Carte mentale de la leçon : les six idées maîtresses autour de la question \og Qu'est-ce que la philosophie ? \fg{}}
\end{popfigure}

\begin{bilan}
\textbf{Définitions clés à connaître}
\begin{itemize}
\item \cle{Philosophie} : du grec \emph{philein} (aimer) et \emph{sophia} (sagesse, savoir) : littéralement \og amour de la sagesse \fg{}. Recherche rationnelle, critique et désintéressée de la vérité.
\item \cle{Étonnement} (\emph{thaumazein}) : selon Platon et Aristote, sentiment d'admiration et de surprise devant le monde, origine de la philosophie.
\item \cle{Doute} : attitude qui consiste à suspendre son jugement pour examiner les opinions ; point de départ de la réflexion (Descartes).
\item \cle{Sagesse} : manière de vivre conforme à la raison, visée du philosophe.
\item \cle{Argumentation} : ensemble de raisons logiquement ordonnées pour justifier ou réfuter une thèse.
\item \cle{Esprit critique} : capacité à examiner, discuter et juger par soi-même au lieu d'accepter les opinions reçues.
\end{itemize}

\textbf{Les thèses des auteurs à retenir}
\begin{center}
\begin{tabularx}{\linewidth}{|l|X|}
\hline
\rowcolor{popblueL} \textbf{Auteur} & \textbf{Thèse essentielle} \\
\hline
Pythagore & Seul le sage est philosophe : l'homme \emph{aime} la sagesse sans la posséder pleinement. \\
\hline
Platon & La philosophie naît de l'étonnement ; le philosophe s'arrache aux opinions (allégorie de la caverne). \\
\hline
Aristote & \og Tous les hommes désirent naturellement connaître \fg{} ; la philosophie commence par l'étonnement. \\
\hline
Descartes & Le doute méthodique est le premier pas vers la vérité : \og Je pense, donc je suis \fg{}. \\
\hline
Kant & La philosophie répond à quatre questions : que puis-je savoir ? que dois-je faire ? que m'est-il permis d'espérer ? qu'est-ce que l'homme ? \\
\hline
\end{tabularx}
\end{center}

\textbf{Méthodes express}
\begin{itemize}
\item \cle{Méthode philosophique} : (1) s'étonner et douter des évidences ; (2) analyser et définir les notions ; (3) problématiser (poser une question) ; (4) argumenter et discuter les thèses ; (5) conclure de façon nuancée.
\item \cle{Philosophie / Science} : même exigence de rigueur et de rationalité ; mais la science étudie des objets précis par l'expérience et le calcul, la philosophie interroge le sens, les valeurs et les fondements. La philosophie est la \og mère des sciences \fg{}.
\item \cle{Importance} : former le jugement, lutter contre les préjugés, apprendre à débattre, devenir un citoyen libre et responsable.
\end{itemize}
\end{bilan}

\begin{aretenir}[Checklist avant le Bac]
Avant l'épreuve, vérifie que tu sais :
\begin{itemize}
\item[$\square$] Donner l'étymologie grecque du mot \emph{philosophie} (\emph{philein} + \emph{sophia}).
\item[$\square$] Citer au moins deux définitions de la philosophie avec leurs auteurs.
\item[$\square$] Situer l'origine de la philosophie en Grèce au VI\textsuperscript{e} siècle avant J.-C.
\item[$\square$] Expliquer le rôle de l'étonnement et du doute dans la naissance de la philosophie.
\item[$\square$] Énumérer les grands objets de la philosophie (l'homme, le monde, Dieu, la morale, la vérité).
\item[$\square$] Décrire les étapes de la méthode philosophique (doute, analyse, problématisation, argumentation).
\item[$\square$] Distinguer clairement opinion (\emph{doxa}) et vérité.
\item[$\square$] Comparer philosophie et science : points communs et différences.
\item[$\square$] Justifier l'importance de la philosophie par au moins trois arguments.
\item[$\square$] Illustrer chaque notion par un exemple concret de la vie courante camerounaise.
\item[$\square$] Rédiger une introduction de dissertation : accroche, définition du sujet, problématique, annonce du plan.
\item[$\square$] Conclure en répondant clairement à la problématique posée.
\end{itemize}
\end{aretenir}

\findecours

\end{document}
